\documentclass[12pt]{article}

% ---------------------------------------------------------------------
%  Datos de la práctica
% ---------------------------------------------------------------------
\newcommand{\NumPractica}{1}
\newcommand{\TituloPractica}{Operaciones Básicas sobre Lenguajes}
\newcommand{\PracticaCorta}{Práctica \NumPractica}
\newcommand{\Alumno}{Rico Espluga Juan Diego}
\newcommand{\Boleta}{2025630474}
\newcommand{\Grupo}{4CV4}
\newcommand{\Docente}{Gabriel Hurtado Avilés}
\newcommand{\FechaEntrega}{22 de septiembre de 2026}

% ======================================================================
%  Preámbulo común de los reportes de Teoría de la Computación
%  Compilar con XeLaTeX + Biber:  latexmk -xelatex main.tex
% ======================================================================
\usepackage{fontspec}
\IfFontExistsTF{Times New Roman}
  {\setmainfont{Times New Roman}}
  {\setmainfont{TeX Gyre Termes}}
\setmonofont{DejaVu Sans Mono}[Scale=0.82]
\newfontfamily\simbolos{DejaVu Sans}[Scale=0.92]

\usepackage[spanish,mexico,es-noshorthands]{babel}
\usepackage{csquotes}

\usepackage[letterpaper,margin=2.54cm,headheight=15pt]{geometry}
\usepackage{setspace}
\doublespacing

\usepackage{amsmath,amssymb}
\usepackage{graphicx}
\usepackage[table]{xcolor}
\usepackage{booktabs,longtable,array,calc,multirow}
\usepackage{caption}
\captionsetup{font={small,singlespacing},labelfont=bf,labelsep=period}
\usepackage{enumitem}
\setlist{nosep,leftmargin=*}
\usepackage{tikz}
\usepackage{fancyhdr}
\usepackage{titlesec}
\usepackage[style=apa,backend=biber]{biblatex}
\usepackage[hidelinks]{hyperref}

% --- Símbolos Unicode que la fuente de texto no incluye ------------------
\usepackage{newunicodechar}
\newunicodechar{⁰}{\textsuperscript{0}}
\newunicodechar{⁴}{\textsuperscript{4}}
\newunicodechar{⁵}{\textsuperscript{5}}
\newunicodechar{⁶}{\textsuperscript{6}}
\newunicodechar{⁷}{\textsuperscript{7}}
\newunicodechar{⁸}{\textsuperscript{8}}
\newunicodechar{⁹}{\textsuperscript{9}}
\newunicodechar{⁺}{\textsuperscript{+}}
\newunicodechar{ⁿ}{\textsuperscript{\textit{n}}}
\newunicodechar{ᵏ}{\textsuperscript{\textit{k}}}
\newunicodechar{ᴿ}{\textsuperscript{\textit{R}}}
\newunicodechar{₀}{\textsubscript{0}}
\newunicodechar{₁}{\textsubscript{1}}
\newunicodechar{₂}{\textsubscript{2}}
\newunicodechar{₃}{\textsubscript{3}}
\newunicodechar{∅}{\ensuremath{\varnothing}}
\newunicodechar{∈}{\ensuremath{\in}}
\newunicodechar{∉}{\ensuremath{\notin}}
\newunicodechar{∣}{\ensuremath{\mid}}
\newunicodechar{∩}{\ensuremath{\cap}}
\newunicodechar{∪}{\ensuremath{\cup}}
\newunicodechar{⊆}{\ensuremath{\subseteq}}
\newunicodechar{⋃}{\ensuremath{\bigcup}}
\newunicodechar{⋯}{\ensuremath{\cdots}}
\newunicodechar{≤}{\ensuremath{\le}}
\newunicodechar{≥}{\ensuremath{\ge}}
\newunicodechar{≠}{\ensuremath{\neq}}
\newunicodechar{→}{\ensuremath{\to}}
\newunicodechar{↦}{\ensuremath{\mapsto}}
\newunicodechar{×}{\ensuremath{\times}}
\newunicodechar{·}{\ensuremath{\cdot}}
\newunicodechar{𝒫}{\ensuremath{\mathcal{P}}}
\newunicodechar{✓}{{\simbolos ✓}}
\newunicodechar{✗}{{\simbolos ✗}}
\newunicodechar{↔}{\ensuremath{\leftrightarrow}}
\newunicodechar{⇔}{\ensuremath{\Leftrightarrow}}
\newunicodechar{∋}{\ensuremath{\ni}}
\newunicodechar{⊥}{\ensuremath{\bot}}

\newunicodechar{ᵐ}{\textsuperscript{\textit{m}}}
\newunicodechar{ᵖ}{\textsuperscript{\textit{p}}}
\newunicodechar{ᵢ}{\textsubscript{\textit{i}}}
\newunicodechar{ⱼ}{\textsubscript{\textit{j}}}
\newunicodechar{ₓ}{\textsubscript{\textit{x}}}
\newunicodechar{⁻}{\textsuperscript{−}}
\newunicodechar{≡}{\ensuremath{\equiv}}
\newunicodechar{¹}{\textsuperscript{1}}
\newunicodechar{ᵃ}{\textsuperscript{\textit{a}}}
\newunicodechar{ᵇ}{\textsuperscript{\textit{b}}}
\newunicodechar{ᶜ}{\textsuperscript{\textit{c}}}
\newunicodechar{∃}{\ensuremath{\exists}}
\newunicodechar{∀}{\ensuremath{\forall}}

\newunicodechar{ˣ}{\textsuperscript{\textit{x}}}
\newunicodechar{ⁱ}{\textsuperscript{\textit{i}}}
\newunicodechar{ₙ}{\textsubscript{\textit{n}}}
\newunicodechar{⟶}{\ensuremath{\longrightarrow}}
\newunicodechar{⇒}{\ensuremath{\Rightarrow}}
\newunicodechar{λ}{\ensuremath{\lambda}}
\newunicodechar{δ}{\ensuremath{\delta}}
\newunicodechar{Σ}{\ensuremath{\Sigma}}

% --- Colores institucionales ------------------------------------------
\definecolor{guinda}{HTML}{6C1D45}
\definecolor{azulescom}{HTML}{006699}
\definecolor{gristenue}{HTML}{F4F4F6}

% --- Títulos ----------------------------------------------------------
\titleformat{\section}{\normalfont\large\bfseries\color{guinda}}{\thesection.}{0.6em}{}
\titleformat{\subsection}{\normalfont\normalsize\bfseries}{\thesubsection.}{0.6em}{}
\titleformat{\subsubsection}{\normalfont\normalsize\bfseries\itshape}{\thesubsubsection.}{0.6em}{}
\titlespacing*{\section}{0pt}{1.2em}{0.4em}
\titlespacing*{\subsection}{0pt}{0.9em}{0.3em}
\titlespacing*{\subsubsection}{0pt}{0.7em}{0.2em}

% --- Encabezado y pie -------------------------------------------------
\pagestyle{fancy}
\fancyhf{}
\fancyhead[L]{\small\textcolor{guinda}{Teoría de la Computación}}
\fancyhead[R]{\small\PracticaCorta}
\fancyfoot[C]{\small\thepage}
\renewcommand{\headrulewidth}{0.4pt}

% --- Tablas y listas que genera pandoc --------------------------------
\providecommand{\tightlist}{\setlength{\itemsep}{0pt}\setlength{\parskip}{0pt}}
\providecommand{\passthrough}[1]{#1}
\AtBeginEnvironment{longtable}{\singlespacing\small}
\setlength{\LTpre}{0.6em}\setlength{\LTpost}{0.6em}
\renewcommand{\arraystretch}{1.25}

% --- Código (resaltado de pandoc, estilo tango) -----------------------
\usepackage{fancyvrb}
\usepackage{fvextra}
\newcommand{\VerbBar}{|}
\newcommand{\VERB}{\Verb[commandchars=\\\{\}]}
\DefineVerbatimEnvironment{Highlighting}{Verbatim}{commandchars=\\\{\},fontsize=\small,breaklines,breakanywhere}
\usepackage{framed}
\colorlet{shadecolor}{gristenue}
\newenvironment{Shaded}{\begin{singlespace}\begin{snugshade}}{\end{snugshade}\end{singlespace}}
\newcommand{\AlertTok}[1]{\textcolor[rgb]{0.94,0.16,0.16}{#1}}
\newcommand{\AnnotationTok}[1]{\textcolor[rgb]{0.56,0.35,0.01}{\textbf{\textit{#1}}}}
\newcommand{\AttributeTok}[1]{\textcolor[rgb]{0.13,0.29,0.53}{#1}}
\newcommand{\BaseNTok}[1]{\textcolor[rgb]{0.00,0.00,0.81}{#1}}
\newcommand{\BuiltInTok}[1]{#1}
\newcommand{\CharTok}[1]{\textcolor[rgb]{0.31,0.60,0.02}{#1}}
\newcommand{\CommentTok}[1]{\textcolor[rgb]{0.56,0.35,0.01}{\textit{#1}}}
\newcommand{\CommentVarTok}[1]{\textcolor[rgb]{0.56,0.35,0.01}{\textbf{\textit{#1}}}}
\newcommand{\ConstantTok}[1]{\textcolor[rgb]{0.56,0.35,0.01}{#1}}
\newcommand{\ControlFlowTok}[1]{\textcolor[rgb]{0.13,0.29,0.53}{\textbf{#1}}}
\newcommand{\DataTypeTok}[1]{\textcolor[rgb]{0.13,0.29,0.53}{#1}}
\newcommand{\DecValTok}[1]{\textcolor[rgb]{0.00,0.00,0.81}{#1}}
\newcommand{\DocumentationTok}[1]{\textcolor[rgb]{0.56,0.35,0.01}{\textbf{\textit{#1}}}}
\newcommand{\ErrorTok}[1]{\textcolor[rgb]{0.64,0.00,0.00}{\textbf{#1}}}
\newcommand{\ExtensionTok}[1]{#1}
\newcommand{\FloatTok}[1]{\textcolor[rgb]{0.00,0.00,0.81}{#1}}
\newcommand{\FunctionTok}[1]{\textcolor[rgb]{0.13,0.29,0.53}{\textbf{#1}}}
\newcommand{\ImportTok}[1]{#1}
\newcommand{\InformationTok}[1]{\textcolor[rgb]{0.56,0.35,0.01}{\textbf{\textit{#1}}}}
\newcommand{\KeywordTok}[1]{\textcolor[rgb]{0.13,0.29,0.53}{\textbf{#1}}}
\newcommand{\NormalTok}[1]{#1}
\newcommand{\OperatorTok}[1]{\textcolor[rgb]{0.81,0.36,0.00}{\textbf{#1}}}
\newcommand{\OtherTok}[1]{\textcolor[rgb]{0.56,0.35,0.01}{#1}}
\newcommand{\PreprocessorTok}[1]{\textcolor[rgb]{0.56,0.35,0.01}{\textit{#1}}}
\newcommand{\RegionMarkerTok}[1]{#1}
\newcommand{\SpecialCharTok}[1]{\textcolor[rgb]{0.81,0.36,0.00}{\textbf{#1}}}
\newcommand{\SpecialStringTok}[1]{\textcolor[rgb]{0.31,0.60,0.02}{#1}}
\newcommand{\StringTok}[1]{\textcolor[rgb]{0.31,0.60,0.02}{#1}}
\newcommand{\VariableTok}[1]{\textcolor[rgb]{0.00,0.00,0.00}{#1}}
\newcommand{\VerbatimStringTok}[1]{\textcolor[rgb]{0.31,0.60,0.02}{#1}}
\newcommand{\WarningTok}[1]{\textcolor[rgb]{0.56,0.35,0.01}{\textbf{\textit{#1}}}}
\RecustomVerbatimEnvironment{verbatim}{Verbatim}{fontsize=\small,breaklines,breakanywhere}

% --- Evidencias -------------------------------------------------------
% \evidencia[ancho]{ruta}{leyenda}: inserta la captura si existe; si aún
% no se ha agregado, dibuja un recuadro que indica la ruta esperada.
\newcommand{\evidencia}[3][0.95]{%
  \begin{figure}[htbp]\centering
  \IfFileExists{#2}{%
    \fbox{\includegraphics[width=#1\linewidth,height=0.42\textheight,keepaspectratio]{#2}}%
  }{%
    \fcolorbox{guinda}{guinda!6}{\parbox[c][3.2cm][c]{0.8\linewidth}{\centering\singlespacing
      \textbf{\textcolor{guinda}{Captura pendiente}}\\[2pt]
      \footnotesize\ttfamily\detokenize{#2}}}%
  }
  \caption{#3}
  \end{figure}}
% \evidenciaaqui: igual que \evidencia, pero sin flotar (se queda junto al texto que la explica).
\newcommand{\evidenciaaqui}[3][0.95]{%
  \par\begin{center}
  \IfFileExists{#2}{%
    \fbox{\includegraphics[width=#1\linewidth,height=0.4\textheight,keepaspectratio]{#2}}%
  }{%
    \fcolorbox{guinda}{guinda!6}{\parbox[c][3.2cm][c]{0.8\linewidth}{\centering\singlespacing
      \textbf{\textcolor{guinda}{Captura pendiente}}\\[2pt]
      \footnotesize\ttfamily\detokenize{#2}}}%
  }
  \captionof{figure}{#3}
  \end{center}}

\addbibresource{referencias.bib}
\nocite{gribkoff2013,lunabenoso2022,turing1936,mavridou2018,davis2019}

\begin{document}
% Portada institucional (IPN · ESCOM)
\begin{titlepage}
\singlespacing
\thispagestyle{empty}
\noindent
\begin{minipage}[c]{0.18\textwidth}\centering
  \includegraphics[width=\linewidth]{img/logo-ipn.png}
\end{minipage}\hfill
\begin{minipage}[c]{0.58\textwidth}\centering
  {\Large\bfseries Instituto Politécnico Nacional\par}
  \vspace{4pt}
  {\large\bfseries Escuela Superior de Cómputo\par}
\end{minipage}\hfill
\begin{minipage}[c]{0.20\textwidth}\centering
  \includegraphics[width=\linewidth]{img/logo-escom.png}
\end{minipage}

\vspace{0.8cm}
{\color{guinda}\hrule height 1.6pt}
\vspace{0.6cm}

\begin{center}
  {\large Ingeniería en Sistemas Computacionales · Plan 2020\par}
  \vspace{0.3cm}
  {\large Unidad de aprendizaje: \textbf{Teoría de la Computación}\par}

  \vfill
  {\LARGE\bfseries\color{guinda} Práctica \NumPractica\par}
  \vspace{0.4cm}
  {\Large\bfseries \TituloPractica\par}
  \vfill

  \renewcommand{\arraystretch}{1.5}
  \begin{tabular}{>{\bfseries}r l}
    Alumno:  & \Alumno \\
    Boleta:  & \Boleta \\
    Grupo:   & \Grupo \\
    Docente: & \Docente \\
    Fecha de entrega: & \FechaEntrega \\
  \end{tabular}
\end{center}

\vspace{0.8cm}
{\color{guinda}\hrule height 1.6pt}
\end{titlepage}


\tableofcontents
\clearpage

\section{Objetivo}
Aplicar los conceptos de la Unidad Temática I al diseño de autómatas finitos en JFLAP y al desarrollo de una aplicación con interfaz gráfica que implemente las operaciones básicas sobre cadenas y lenguajes, y construir el entorno de trabajo, con control de versiones y contenedores, que se usará el resto del curso.

La práctica tiene cinco ejercicios: el primero prepara el entorno; el segundo y el tercero dan el sustento conceptual, y el cuarto y el quinto lo aplican. El código completo está en \texttt{src/} y \texttt{tests/}, los autómatas en \texttt{automatas/} y la versión en Markdown de este reporte en \texttt{docs/}.

\hypertarget{ejercicio-1.-entorno-de-trabajo-control-de-versiones-y-contenedores}{%
\section{Ejercicio 1. Entorno de trabajo: control de versiones y contenedores}\label{ejercicio-1.-entorno-de-trabajo-control-de-versiones-y-contenedores}}

\hypertarget{parte-a.-investigaciuxf3n}{%
\subsection{Parte A. Investigación}\label{parte-a.-investigaciuxf3n}}

\hypertarget{control-de-versiones}{%
\subsubsection{Control de versiones}\label{control-de-versiones}}

Un sistema de control de versiones registra cada estado de un conjunto de archivos junto con su autor, su fecha y una explicación del cambio, de modo que cualquier versión anterior puede recuperarse o compararse con la actual \autocite[cap. 1]{chacon2014}. En un equipo resuelve tres problemas que aparecen en cuanto dos personas tocan el mismo proyecto: saber quién cambió qué y por qué, trabajar en paralelo sin sobrescribir el trabajo ajeno y volver atrás cuando un cambio introduce un error. Git es \emph{distribuido}: cada copia del repositorio contiene la historia completa, por lo que se puede confirmar, crear ramas y consultar el historial sin conexión.

\textbf{Git y GitHub.} Git es el programa que administra la historia en la computadora de cada persona. GitHub es un servicio en la nube que aloja repositorios Git y añade lo que Git no trae: Pull Requests con revisión, integración continua (GitHub Actions), control de acceso e incidencias. Se puede usar Git sin GitHub, pero GitHub no tiene sentido sin Git.

\begin{longtable}[]{@{}
  >{\raggedright\arraybackslash}p{(\columnwidth - 4\tabcolsep) * \real{0.1111}}
  >{\raggedright\arraybackslash}p{(\columnwidth - 4\tabcolsep) * \real{0.4444}}
  >{\raggedright\arraybackslash}p{(\columnwidth - 4\tabcolsep) * \real{0.4444}}@{}}
\caption{Términos de control de versiones.}\tabularnewline
\toprule\noalign{}
\begin{minipage}[b]{\linewidth}\raggedright
Concepto
\end{minipage} & \begin{minipage}[b]{\linewidth}\raggedright
Definición
\end{minipage} & \begin{minipage}[b]{\linewidth}\raggedright
Ejemplo en esta práctica
\end{minipage} \\
\midrule\noalign{}
\endfirsthead
\toprule\noalign{}
\begin{minipage}[b]{\linewidth}\raggedright
Concepto
\end{minipage} & \begin{minipage}[b]{\linewidth}\raggedright
Definición
\end{minipage} & \begin{minipage}[b]{\linewidth}\raggedright
Ejemplo en esta práctica
\end{minipage} \\
\midrule\noalign{}
\endhead
\bottomrule\noalign{}
\endlastfoot
Repositorio & Carpeta cuya historia completa guarda Git en el subdirectorio oculto \texttt{.git}. & \texttt{practica1-tc}, con \texttt{src/}, \texttt{tests/}, \texttt{docs/} y \texttt{entorno/}. \\
Confirmación (\emph{commit}) & Fotografía de los archivos preparados, con un mensaje y un identificador único (hash SHA-1). & \texttt{Agrega\ el\ enumerador\ de\ Σⁿ\ con\ límite\ de\ cadenas}. \\
Rama (\emph{branch}) & Puntero móvil a una confirmación; permite desarrollar sin afectar la rama principal. & \texttt{feature/interfaz}, para la interfaz Flet. \\
Fusión (\emph{merge}) & Integra en una rama los cambios de otra; si ambas avanzaron, crea una confirmación con dos padres. & \texttt{feature/interfaz} se fusionó en \texttt{main}. \\
Conflicto de fusión & Ocurre cuando las dos ramas modificaron las mismas líneas; Git no elige y marca el archivo con \texttt{\textless{}\textless{}\textless{}\textless{}\textless{}\textless{}\textless{}}, \texttt{=======} y \texttt{\textgreater{}\textgreater{}\textgreater{}\textgreater{}\textgreater{}\textgreater{}\textgreater{}} para que una persona decida. & Dos versiones distintas del mensaje de error del límite en \texttt{lenguajes.py}. \\
Pull Request & Solicitud en GitHub para fusionar una rama, con descripción, revisión, comentarios y resultado de las pruebas \autocite{githubpr2026}. & «Interfaz con menú lateral y ayuda por sección». \\
\texttt{.gitignore} & Lista de patrones de archivos que Git no debe rastrear, normalmente porque se regeneran. & \texttt{venv/}, \texttt{\_\_pycache\_\_/}, \texttt{*.pyc}, \texttt{salidas/}. \\
\texttt{README} & Documento de entrada del repositorio; GitHub lo muestra en la página principal. & Portada con datos del alumno, índice e instrucciones. \\
\end{longtable}

\hypertarget{contenedores}{%
\subsubsection{Contenedores}\label{contenedores}}

Un contenedor es un proceso aislado del sistema operativo anfitrión que ve su propio sistema de archivos, su propia red y sus propios procesos, pero \textbf{comparte el núcleo} del anfitrión. Una máquina virtual, en cambio, emula hardware completo y arranca su propio sistema operativo sobre un hipervisor. De ahí las tres diferencias que pide la práctica \autocite{felter2015}:

\begin{itemize}
\tightlist
\item
  \textbf{Arranque.} Un contenedor inicia en fracciones de segundo, porque sólo lanza un proceso; una máquina virtual tarda lo que tarda en arrancar un sistema operativo.
\item
  \textbf{Tamaño.} Una imagen de contenedor contiene únicamente la aplicación y sus bibliotecas (decenas o cientos de MB); una máquina virtual incluye un sistema operativo completo (varios GB).
\item
  \textbf{Aislamiento.} La máquina virtual aísla más, porque cada una tiene su propio núcleo; en el contenedor el aislamiento se logra con \emph{namespaces} y \emph{cgroups} del mismo núcleo, que es una frontera más delgada.
\end{itemize}

Los cuatro términos de Docker que se usan en el ejercicio son \autocite{dockerdocs2026}: la \textbf{imagen}, plantilla de sólo lectura formada por capas (por ejemplo, \texttt{python:3.12-slim} más las dependencias); el \textbf{contenedor}, una instancia en ejecución de una imagen, con una capa escribible propia; el \textbf{volumen}, almacenamiento que vive fuera del ciclo de vida del contenedor y que aquí monta la carpeta del proyecto en \texttt{/app} para que los cambios al código se vean sin reconstruir; y el \textbf{puerto publicado}, que conecta un puerto del equipo con uno del contenedor (\texttt{8550:8550}), necesario para abrir la interfaz desde el navegador.

\textbf{Entorno virtual de Python.} Es una carpeta con su propio \texttt{site-packages} y enlaces al intérprete que lo creó; aísla las \emph{bibliotecas} de cada proyecto, no el intérprete \autocite{pythonvenv2026}. Por eso un entorno virtual no cambia la versión de Python: \texttt{python\ -m\ venv} con Python 3.12 siempre produce un entorno 3.12. Para probar tres versiones hacen falta tres intérpretes, y cada contenedor aporta el suyo.

\textbf{Por qué fijar la imagen.} \texttt{python:latest} apunta a la versión más reciente en el momento de construir; un día es 3.13 y al siguiente 3.14, con lo que la misma receta produce entornos distintos y un fallo puede aparecer sin que nadie haya cambiado el código. \texttt{python:3.12-slim} fija la versión menor (y la variante ligera), de modo que la construcción es reproducible; las correcciones de la versión de parche siguen llegando.

\hypertarget{parte-b.-repositorio}{%
\subsection{Parte B. Repositorio}\label{parte-b.-repositorio}}

El repositorio \texttt{practica1-tc} contiene \texttt{README.md} (portada con nombre, boleta, grupo, carrera, unidad de aprendizaje, profesor y fecha, índice enlazado e instrucciones), un \texttt{.gitignore} que excluye \texttt{venv/}, \texttt{.venv/}, \texttt{\_\_pycache\_\_/}, \texttt{*.pyc}, \texttt{.pytest\_cache/}, \texttt{salidas/} y los auxiliares de LaTeX, y que \textbf{no} excluye los \texttt{.jff}, y la historia de confirmaciones. Se hicieron más de cinco confirmaciones con mensajes en imperativo que dicen qué cambia; la interfaz se desarrolló en la rama \texttt{feature/interfaz} y entró a \texttt{main} mediante un Pull Request descrito y fusionado.

\evidencia[0.95]{../evidencias/git/git-log-graph.png}{Salida de \texttt{git\ log\ -\/-oneline\ -\/-graph\ -\/-all}.}

\evidencia[0.95]{../evidencias/git/pull-request.png}{Pull Request de \texttt{feature/interfaz} fusionado en \texttt{main}.}

\hypertarget{parte-c.-los-tres-entornos-de-python}{%
\subsection{Parte C. Los tres entornos de Python}\label{parte-c.-los-tres-entornos-de-python}}

\hypertarget{archivos-versionados}{%
\subsubsection{Archivos versionados}\label{archivos-versionados}}

El \texttt{Dockerfile} (carpeta \texttt{entorno/}) es único y recibe la versión como argumento. La tabla explica cada instrucción.

\begin{longtable}[]{@{}
  >{\raggedright\arraybackslash}p{(\columnwidth - 2\tabcolsep) * \real{0.3462}}
  >{\raggedright\arraybackslash}p{(\columnwidth - 2\tabcolsep) * \real{0.6538}}@{}}
\caption{Instrucciones del \texttt{Dockerfile}.}\tabularnewline
\toprule\noalign{}
\begin{minipage}[b]{\linewidth}\raggedright
Instrucción
\end{minipage} & \begin{minipage}[b]{\linewidth}\raggedright
Función
\end{minipage} \\
\midrule\noalign{}
\endfirsthead
\toprule\noalign{}
\begin{minipage}[b]{\linewidth}\raggedright
Instrucción
\end{minipage} & \begin{minipage}[b]{\linewidth}\raggedright
Función
\end{minipage} \\
\midrule\noalign{}
\endhead
\bottomrule\noalign{}
\endlastfoot
\texttt{ARG\ PYTHON\_VERSION=3.12} & Declara el argumento de construcción; \texttt{compose.yml} le da 3.11, 3.12 o 3.13. El valor por omisión permite construir la imagen sin Compose. \\
\texttt{FROM\ python:\$\{PYTHON\_VERSION\}-slim} & Imagen base oficial en su variante ligera, con la versión elegida. \\
\texttt{RUN\ python\ -m\ venv\ /opt/venv} & Crea el entorno virtual fuera de \texttt{/app}, porque el volumen de Compose reemplaza \texttt{/app} con la carpeta del equipo y borraría un entorno creado ahí. \\
\texttt{ENV\ PATH="/opt/venv/bin:\$PATH"} & Antepone el \texttt{bin/} del entorno al \texttt{PATH}; equivale a activarlo para todas las órdenes siguientes. \\
\texttt{WORKDIR\ /app} & Carpeta de trabajo dentro del contenedor. \\
\texttt{COPY\ requirements.txt\ .} & Copia sólo la lista de dependencias, para que la capa siguiente se reutilice mientras esa lista no cambie. \\
\texttt{RUN\ pip\ install\ …\ -r\ requirements.txt} & Actualiza pip e instala las versiones fijadas; \texttt{-\/-no-cache-dir} evita guardar los paquetes descargados en la imagen. \\
\texttt{COPY\ .\ .} & Copia el proyecto (útil si se ejecuta sin volumen). \\
\texttt{ENV\ FLET\_FORCE\_WEB\_SERVER=true} y \texttt{FLET\_SERVER\_PORT=8550} & Sin escritorio, Flet se sirve como aplicación web en el puerto 8550. \\
\texttt{CMD\ {[}"python",\ "src/app.py"{]}} & Orden por omisión: abrir la interfaz. \\
\end{longtable}

\texttt{compose.yml} declara los servicios \texttt{py311}, \texttt{py312} y \texttt{py313} con el mismo \texttt{Dockerfile} \autocite{composedocs2026} y distinto \texttt{PYTHON\_VERSION}; los tres montan la carpeta del proyecto en \texttt{/app}, \texttt{py311} y \texttt{py313} ejecutan \texttt{pytest\ -q} por omisión y sólo \texttt{py312} publica el puerto 8550 para la interfaz. \texttt{requirements.txt} fija \texttt{flet{[}web{]}==0.86.5}, \texttt{fastapi==0.115.12} y \texttt{pytest==8.3.4}. FastAPI se fija porque \texttt{flet-web} lo requiere sin acotar su versión y las publicaciones de FastAPI posteriores a la 0.115 eliminaron un parámetro que \texttt{flet-web} 0.86.5 todavía usa; sin esa línea la interfaz no arranca en un entorno nuevo.

\hypertarget{justificaciuxf3n-de-las-versiones}{%
\subsubsection{Justificación de las versiones}\label{justificaciuxf3n-de-las-versiones}}

\begin{longtable}[]{@{}
  >{\raggedright\arraybackslash}p{(\columnwidth - 4\tabcolsep) * \real{0.1579}}
  >{\raggedright\arraybackslash}p{(\columnwidth - 4\tabcolsep) * \real{0.1579}}
  >{\raggedright\arraybackslash}p{(\columnwidth - 4\tabcolsep) * \real{0.6842}}@{}}
\caption{Elección de versiones.}\tabularnewline
\toprule\noalign{}
\begin{minipage}[b]{\linewidth}\raggedright
Servicio
\end{minipage} & \begin{minipage}[b]{\linewidth}\raggedright
Versión
\end{minipage} & \begin{minipage}[b]{\linewidth}\raggedright
Motivo
\end{minipage} \\
\midrule\noalign{}
\endfirsthead
\toprule\noalign{}
\begin{minipage}[b]{\linewidth}\raggedright
Servicio
\end{minipage} & \begin{minipage}[b]{\linewidth}\raggedright
Versión
\end{minipage} & \begin{minipage}[b]{\linewidth}\raggedright
Motivo
\end{minipage} \\
\midrule\noalign{}
\endhead
\bottomrule\noalign{}
\endlastfoot
\texttt{py311} & 3.11 & Flet 0.86.5 exige Python ≥ 3.10 \autocite{fletpypi2026}, pero 3.10 llegó al fin de su soporte en octubre de 2026 \autocite{pythonversions2026}; 3.11 es la más antigua que conserva parches de seguridad (hasta octubre de 2027) y que Flet admite. \\
\texttt{py312} & 3.12 & Versión de referencia del curso; recibe parches de seguridad hasta octubre de 2028. Sirve la interfaz. \\
\texttt{py313} & 3.13 & Versión estable reciente (soporte hasta octubre de 2029); detecta dependencias de comportamientos que cambiaron. \\
\end{longtable}

\hypertarget{verificaciuxf3n}{%
\subsubsection{Verificación}\label{verificaciuxf3n}}

Cada contenedor se construye con \texttt{docker\ compose\ build} y se comprueba con \texttt{docker\ compose\ run\ -\/-rm\ py311\ python\ -\/-version} y sus equivalentes; las salidas son \texttt{Python\ 3.11.x}, \texttt{Python\ 3.12.x} y \texttt{Python\ 3.13.x}.

\evidencia[0.95]{../evidencias/docker/python-version.png}{\texttt{python\ -\/-version} en los contenedores \texttt{py311},
\texttt{py312} y \texttt{py313}.}

\hypertarget{ejercicio-2.-investigaciuxf3n-quuxe9-es-la-teoruxeda-de-la-computaciuxf3n}{%
\section{Ejercicio 2. Investigación: qué es la Teoría de la Computación}\label{ejercicio-2.-investigaciuxf3n-quuxe9-es-la-teoruxeda-de-la-computaciuxf3n}}

\hypertarget{objeto-de-estudio-y-preguntas-centrales}{%
\subsection{Objeto de estudio y preguntas centrales}\label{objeto-de-estudio-y-preguntas-centrales}}

La Teoría de la Computación estudia, con modelos matemáticos, qué pueden calcular las máquinas y con qué recursos. \textcite[cap. 0]{sipser2013} la organiza alrededor de una sola pregunta: \emph{cuáles son las capacidades y las limitaciones fundamentales de las computadoras}, y la descompone en tres áreas que responden a versiones cada vez más finas de esa pregunta. \textcite[cap. 1]{hopcroft2007} parten de un punto distinto: definen la teoría de autómatas como el estudio de dispositivos de cómputo abstractos, y llegan a la computabilidad y a la complejidad como consecuencias de preguntar qué lenguajes reconoce cada dispositivo. \textcite[cap. 1]{linz2017} adopta un enfoque todavía más concreto: los lenguajes formales, las gramáticas y los autómatas son el lenguaje común con el que se describen y comparan los modelos.

Las definiciones coinciden en el objeto (modelos abstractos de cómputo) y en el método (definición formal y demostración), pero difieren en el orden y en el énfasis. Sipser parte de la pregunta filosófica y presenta los modelos como herramientas para contestarla; Hopcroft, Motwani y Ullman, y Linz, parten de los modelos y dejan que las preguntas surjan de sus límites. Para un curso que empieza por los autómatas, la segunda forma es la más práctica; la primera explica por qué vale la pena el esfuerzo.

Las preguntas que plantea la teoría pueden resumirse en tres: ¿qué problemas se pueden resolver con un algoritmo?, ¿cuáles se pueden resolver eficientemente? y ¿qué modelo mínimo basta para un problema dado?

\hypertarget{oruxedgenes}{%
\subsection{Orígenes}\label{oruxedgenes}}

A principios del siglo XX la lógica matemática atravesaba una crisis: la paradoja de Russell (el conjunto de los conjuntos que no se contienen a sí mismos) mostró que la teoría intuitiva de conjuntos era contradictoria, y otras paradojas de autorreferencia, como la del mentiroso, apuntaban al mismo problema \autocite[cap. 3]{kleene1952}. David Hilbert propuso entonces un programa para dar a la matemática una base segura: formalizarla por completo en sistemas axiomáticos y demostrar, con medios finitos, que esos sistemas son consistentes y completos.

Una pieza del programa era el \emph{Entscheidungsproblem}, planteado por Hilbert y Ackermann en 1928: encontrar un procedimiento mecánico que, dada una fórmula de la lógica de primer orden, decida en un número finito de pasos si es válida. Para responder negativamente había que decir con precisión qué es un «procedimiento mecánico», y esa necesidad dio origen a la teoría. Gödel demostró en 1931 que todo sistema consistente suficientemente expresivo contiene enunciados verdaderos que no puede demostrar, y en 1936 Church \autocite*{church1936} y Turing \autocite*{turing1936}, de manera independiente, mostraron que el \emph{Entscheidungsproblem} no tiene solución: Church con el cálculo λ y Turing con sus máquinas.

\hypertarget{las-tres-ramas}{%
\subsection{Las tres ramas}\label{las-tres-ramas}}

\begin{longtable}[]{@{}
  >{\raggedright\arraybackslash}p{(\columnwidth - 4\tabcolsep) * \real{0.1250}}
  >{\raggedright\arraybackslash}p{(\columnwidth - 4\tabcolsep) * \real{0.4167}}
  >{\raggedright\arraybackslash}p{(\columnwidth - 4\tabcolsep) * \real{0.4583}}@{}}
\caption{Las tres ramas de la teoría \autocite[cap. 0]{sipser2013}.}\tabularnewline
\toprule\noalign{}
\begin{minipage}[b]{\linewidth}\raggedright
Rama
\end{minipage} & \begin{minipage}[b]{\linewidth}\raggedright
Pregunta
\end{minipage} & \begin{minipage}[b]{\linewidth}\raggedright
Ejemplo de resultado
\end{minipage} \\
\midrule\noalign{}
\endfirsthead
\toprule\noalign{}
\begin{minipage}[b]{\linewidth}\raggedright
Rama
\end{minipage} & \begin{minipage}[b]{\linewidth}\raggedright
Pregunta
\end{minipage} & \begin{minipage}[b]{\linewidth}\raggedright
Ejemplo de resultado
\end{minipage} \\
\midrule\noalign{}
\endhead
\bottomrule\noalign{}
\endlastfoot
Teoría de autómatas y lenguajes formales & ¿Qué puede reconocer o generar cada modelo de cómputo restringido? & Los autómatas finitos reconocen exactamente los lenguajes regulares; \{aⁿbⁿ\} no es regular. \\
Computabilidad & ¿Qué problemas se pueden resolver con un algoritmo? & El problema de la parada es indecidible. \\
Complejidad & ¿Qué problemas se pueden resolver con recursos razonables de tiempo y memoria? & La pregunta P = NP sigue abierta. \\
\end{longtable}

\hypertarget{la-tesis-de-church-turing}{%
\subsection{La tesis de Church-Turing}\label{la-tesis-de-church-turing}}

La tesis afirma que \textbf{toda función que puede calcularse mediante un procedimiento efectivo, en el sentido intuitivo, puede calcularse con una máquina de Turing}. En la década de 1930 se propusieron varios modelos independientes: las funciones recursivas generales de Herbrand y Gödel, el cálculo λ de Church, las máquinas de Turing y los sistemas de Post; todos resultaron equivalentes, en el sentido de que calculan exactamente las mismas funciones \autocites[cap. 13]{kleene1952}[sec.~3.3]{sipser2013}. Turing lo demostró para su modelo y el de Church en el apéndice de su artículo.

Se llama \emph{tesis} y no teorema porque uno de sus dos lados, «procedimiento efectivo», es una noción intuitiva y no un objeto matemático; no puede demostrarse una equivalencia entre algo definido y algo que no lo está. Lo que sí está demostrado es la equivalencia entre los modelos formales, y esa coincidencia es la mejor evidencia a favor de la tesis.

\hypertarget{alfabetos-cadenas-y-lenguajes}{%
\subsection{Alfabetos, cadenas y lenguajes}\label{alfabetos-cadenas-y-lenguajes}}

Un \textbf{alfabeto} Σ es un conjunto finito y no vacío de símbolos. Una \textbf{cadena} (o palabra) sobre Σ es una sucesión finita de símbolos de Σ; su longitud \textbar ω\textbar{} es el número de símbolos, y la \textbf{cadena vacía} λ es la única de longitud 0. Σⁿ es el conjunto de las cadenas de longitud n, Σ* = ⋃ₙ≥₀ Σⁿ el de todas las cadenas y Σ⁺ = ⋃ₙ≥₁ Σⁿ el de las no vacías. Un \textbf{lenguaje} sobre Σ es cualquier subconjunto de Σ* \autocite[sec.~1.5]{hopcroft2007}.

\textbf{Operaciones sobre cadenas.} La \emph{concatenación} ωπ escribe π a continuación de ω, es asociativa y tiene a λ como neutro; \textbar ωπ\textbar{} = \textbar ω\textbar{} + \textbar π\textbar. La \emph{potencia} se define por recursión: ω⁰ = λ y ωⁿ⁺¹ = ωⁿω. La \emph{reflexión} ωᴿ invierte el orden: λᴿ = λ y (aω)ᴿ = ωᴿa, con (ωπ)ᴿ = πᴿωᴿ.

\textbf{Operaciones sobre lenguajes.} Como conjuntos, los lenguajes admiten \emph{unión}, \emph{intersección} y \emph{diferencia}. La \emph{concatenación} es L₁L₂ = \{xy : x ∈ L₁, y ∈ L₂\}; la \emph{potencia}, L⁰ = \{λ\} y Lⁿ⁺¹ = LⁿL; la \emph{reflexión}, Lᴿ = \{ωᴿ : ω ∈ L\}. La \textbf{cerradura de Kleene} es L* = ⋃ₙ≥₀ Lⁿ y la \textbf{cerradura positiva}, L⁺ = ⋃ₙ≥₁ Lⁿ \autocite[sec.~1.2]{linz2017}.

\textbf{Por qué Σ⁰ = \{λ\}.} Σ⁰ es el conjunto de las cadenas de longitud 0, y hay exactamente una: λ. Además, \{λ\} es el neutro de la concatenación de lenguajes (L\{λ\} = L), igual que λ lo es de la de cadenas; si Σ⁰ fuera ∅, la identidad Σᵐ⁺ⁿ = ΣᵐΣⁿ fallaría para m = 0, porque ∅L = ∅. \textbf{Σ* frente a Σ⁺.} La única diferencia es λ: Σ* = Σ⁺ ∪ \{λ\} y Σ⁺ = ΣΣ\emph{. En general L⁺ = L} sólo si λ ∈ L.

\hypertarget{la-jerarquuxeda-de-chomsky}{%
\subsection{La jerarquía de Chomsky}\label{la-jerarquuxeda-de-chomsky}}

\textcite{chomsky1956} clasificó las gramáticas según la forma de sus producciones; cada tipo genera una clase de lenguajes que contiene a la siguiente, y a cada clase corresponde una máquina que la reconoce.

\begin{longtable}[]{@{}
  >{\raggedright\arraybackslash}p{(\columnwidth - 6\tabcolsep) * \real{0.1154}}
  >{\raggedright\arraybackslash}p{(\columnwidth - 6\tabcolsep) * \real{0.1154}}
  >{\raggedright\arraybackslash}p{(\columnwidth - 6\tabcolsep) * \real{0.4231}}
  >{\raggedright\arraybackslash}p{(\columnwidth - 6\tabcolsep) * \real{0.3462}}@{}}
\caption{Jerarquía de Chomsky \autocite{hopcroft2007,lewis1998,sipser2013}.}\tabularnewline
\toprule\noalign{}
\begin{minipage}[b]{\linewidth}\raggedright
Tipo
\end{minipage} & \begin{minipage}[b]{\linewidth}\raggedright
Lenguajes
\end{minipage} & \begin{minipage}[b]{\linewidth}\raggedright
Forma de las producciones
\end{minipage} & \begin{minipage}[b]{\linewidth}\raggedright
Máquina reconocedora
\end{minipage} \\
\midrule\noalign{}
\endfirsthead
\toprule\noalign{}
\begin{minipage}[b]{\linewidth}\raggedright
Tipo
\end{minipage} & \begin{minipage}[b]{\linewidth}\raggedright
Lenguajes
\end{minipage} & \begin{minipage}[b]{\linewidth}\raggedright
Forma de las producciones
\end{minipage} & \begin{minipage}[b]{\linewidth}\raggedright
Máquina reconocedora
\end{minipage} \\
\midrule\noalign{}
\endhead
\bottomrule\noalign{}
\endlastfoot
0 & Recursivamente enumerables & α → β, con α no vacía (sin restricción) & Máquina de Turing \\
1 & Sensibles al contexto & αAβ → αγβ, con γ no vacía & Autómata linealmente acotado \\
2 & Independientes del contexto & A → γ & Autómata de pila (no determinista) \\
3 & Regulares & A → aB o A → a (lineal por la derecha) & Autómata finito \\
\end{longtable}

Las inclusiones son estrictas: \{aⁿbⁿ : n ≥ 0\} es independiente del contexto pero no regular, \{aⁿbⁿcⁿ : n ≥ 0\} es sensible al contexto pero no independiente del contexto, y el lenguaje de la parada es recursivamente enumerable pero no recursivo.

\hypertarget{autuxf3matas-finitos-y-expresiones-regulares}{%
\subsection{Autómatas finitos y expresiones regulares}\label{autuxf3matas-finitos-y-expresiones-regulares}}

Un \textbf{autómata finito determinista} es una quíntupla M = (Q, Σ, δ, q₀, F), con Q un conjunto finito de estados, Σ el alfabeto, δ: Q × Σ → Q la función de transición, q₀ ∈ Q el estado inicial y F ⊆ Q los de aceptación. δ se extiende a cadenas con δ̂(q, λ) = q y δ̂(q, ωa) = δ(δ̂(q, ω), a); M acepta ω si δ̂(q₀, ω) ∈ F, y L(M) es el conjunto de las cadenas aceptadas \autocite[sec.~2.2]{hopcroft2007}. Un \textbf{autómata finito no determinista} sólo cambia la función de transición: δ: Q × Σ → 𝒫(Q) devuelve un conjunto de estados, posiblemente vacío, y la cadena se acepta si \emph{alguna} de las sucesiones posibles de estados termina en F.

Los dos modelos son \textbf{equivalentes}: para todo AFN existe un AFD que reconoce el mismo lenguaje \autocite{kozen1997}. \textcite{rabin1959} lo demostraron con la \emph{construcción de subconjuntos}, en la que cada estado del AFD es un conjunto de estados del AFN; el precio es que el AFD puede tener hasta 2ⁿ estados. \textcite{kleene1956} probó además que los lenguajes que reconocen los autómatas finitos son exactamente los que describen las \textbf{expresiones regulares}, construidas a partir de ∅, λ y los símbolos con unión, concatenación y estrella.

Tres problemas reales en los que se aplican expresiones regulares:

\begin{enumerate}
\def\labelenumi{\arabic{enumi}.}
\tightlist
\item
  \textbf{Análisis léxico.} El primer paso de un compilador divide el código en fichas (identificadores, números, palabras reservadas); cada clase de ficha se describe con una expresión, como \texttt{{[}A-Za-z\_{]}{[}A-Za-z0-9\_{]}*} para los identificadores, y el analizador se construye como un AFD.
\item
  \textbf{Validación de datos.} Un formulario rechaza un correo o un código postal mal escritos comparando el texto con una expresión antes de enviarlo.
\item
  \textbf{Búsqueda en texto y registros.} Herramientas como \texttt{grep} encuentran todas las líneas de un archivo de bitácora que contienen, por ejemplo, una dirección IP o un código de error, sin recorrer el archivo más de una vez.
\end{enumerate}

Las expresiones de los lenguajes de programación añaden operadores que no son regulares, como las referencias hacia atrás, y sus motores pueden tardar un tiempo exponencial; el artículo de \textcite{davis2019} del Ejercicio 3 estudia esas diferencias.

\hypertarget{ejercicio-3.-estado-del-arte-cinco-artuxedculos}{%
\section{Ejercicio 3. Estado del arte: cinco artículos}\label{ejercicio-3.-estado-del-arte-cinco-artuxedculos}}

Cada texto se resume en una ficha de seis puntos; para los tres asignados se responden además las preguntas del enunciado.

\hypertarget{texto-1.-gribkoff-2013}{%
\subsection{Texto 1. Gribkoff (2013)}\label{texto-1.-gribkoff-2013}}

\begin{enumerate}
\def\labelenumi{\arabic{enumi}.}
\tightlist
\item
  \textbf{Cita.} Gribkoff, E. (2013). \emph{Applications of deterministic finite automata} {[}Documento de curso, ECS 120{]}. University of California, Davis. \url{https://www.cs.ucdavis.edu/~rogaway/classes/120/spring13/eric-dfa.pdf}
\item
  \textbf{Problema.} Mostrar a estudiantes de un primer curso de teoría que el AFD no es sólo un objeto abstracto, sino un modelo que aparece en sistemas cotidianos.
\item
  \textbf{Método.} Tras repasar la definición formal del AFD y la de la máquina de Mealy, analiza cuatro ejemplos: una máquina expendedora, el comportamiento de los fantasmas de Pac-Man, el diagrama de estados de TCP y el autocompletado de Apache Lucene.
\item
  \textbf{Resultado.} Cada sistema se describe con pocos estados; en Lucene, el transductor comparte prefijos y sufijos entre palabras y así reduce la memoria del índice.
\item
  \textbf{Tema de la Unidad I.} Autómatas finitos: definición formal del AFD y aplicaciones.
\item
  \textbf{Aportación.} Ofrece ejemplos de modelado que sirven de punto de partida para diseñar autómatas propios, como el del fantasma de Pac-Man que se amplía en la Práctica 2.
\end{enumerate}

\textbf{AFD y máquina de Mealy.} El AFD sólo produce un veredicto al final de la lectura: acepta o rechaza. La máquina de Mealy agrega un alfabeto de salida y una función que emite un símbolo en \emph{cada} transición, de modo que convierte una cadena de entrada en una de salida. El autocompletado de Lucene necesita esa salida: no basta con saber que un prefijo pertenece al diccionario; hay que recuperar información de cada palabra completada (en el ejemplo, un número que se obtiene acumulando las salidas del camino).

\textbf{Máquina expendedora con F = ∅.} Si ningún estado es de aceptación, δ̂(q₀, ω) ∉ F para toda ω, así que el lenguaje del autómata es ∅: ninguna cadena se acepta. El modelo sigue siendo útil para describir cómo cambia el estado de la máquina con cada moneda, pero la noción de «cadena aceptada» deja de aportar información; se está usando el autómata como un sistema de transiciones y no como un reconocedor.

\textbf{Afirmación por verificar.} El texto afirma que, gracias al transductor, Lucene mantiene en memoria el índice completo del autocompletado y responde mucho más rápido. No da cifras, versiones ni una referencia; antes de citarla debería confirmarse en una publicación arbitrada o en la documentación oficial de Lucene.

\hypertarget{texto-2.-luna-benoso-et-al.-2022}{%
\subsection{Texto 2. Luna-Benoso et al.~(2022)}\label{texto-2.-luna-benoso-et-al.-2022}}

\begin{enumerate}
\def\labelenumi{\arabic{enumi}.}
\tightlist
\item
  \textbf{Cita.} Luna-Benoso, B., Martínez-Perales, J. C., Cortés-Galicia, J., Flores-Carapia, R., y Silva-García, V. M. (2022). Melanoma detection in dermoscopic images using a cellular automata classifier. \emph{Computers}, \emph{11}(1), 8. \url{https://doi.org/10.3390/computers11010008}
\item
  \textbf{Problema.} Distinguir el melanoma de lesiones benignas en imágenes dermatoscópicas como apoyo al diagnóstico temprano.
\item
  \textbf{Método.} Segmentan la lesión (filtro de mediana, umbral de Otsu y operaciones morfológicas sobre la vecindad de Moore), extraen 11 descriptores estadísticos de textura y clasifican con un autómata celular asociativo de estados binarios.
\item
  \textbf{Resultado.} Con 190 imágenes de la base PH2 obtienen exactitud de 0.978, sensibilidad de 0.944, especificidad de 0.987 y área bajo la curva ROC de 0.965.
\item
  \textbf{Tema de la Unidad I.} Modelos de cómputo; diferencia entre autómata finito y autómata celular.
\item
  \textbf{Aportación.} Un modelo de transiciones locales, barato de calcular, compite con clasificadores mucho más costosos en una aplicación médica.
\end{enumerate}

\textbf{Autómata celular y autómata finito.} Se parecen en que ambos son discretos y evolucionan con una función de transición sobre un conjunto finito de estados. Difieren en dónde vive el estado y en qué se lee. En el autómata celular el \textbf{estado} es de cada celda de una retícula (en el artículo, 0 o 1); el \textbf{vecindario} es el conjunto de celdas que influyen en una celda (la vecindad de Moore: sus ocho vecinas); y la \textbf{regla local} calcula el nuevo estado de cada celda a partir de su vecindario, para todas las celdas a la vez. Un autómata finito tiene un único estado global que cambia al leer un símbolo de una cadena, uno por paso, y termina aceptando o rechazando.

\textbf{Problema y elección del modelo.} El melanoma es el cáncer de piel más agresivo y su pronóstico depende de detectarlo temprano. Los autores eligen autómatas celulares porque se implementan con facilidad y requieren pocos recursos, a diferencia de redes profundas como AlexNet o VGG16.

\textbf{Datos y métricas.} Usan la base PH2 del Hospital Pedro Hispano: 200 imágenes RGB de 768 × 560 píxeles, de las que descartan 10 por falta de nitidez. Reportan exactitud, sensibilidad, especificidad y AUC, con los valores del punto 4.

\hypertarget{texto-3.-turing-1936-trabajo-original}{%
\subsection{Texto 3. Turing (1936), trabajo original}\label{texto-3.-turing-1936-trabajo-original}}

\begin{enumerate}
\def\labelenumi{\arabic{enumi}.}
\tightlist
\item
  \textbf{Cita.} Turing, A. M. (1936).\footnote{El artículo se publicó en las partes 3 y 4 del volumen 42 (serie 2), aparecidas a finales de 1936, pero el volumen se cerró en 1937. El registro del DOI usa la fecha del volumen y la literatura, la de las partes.} On computable numbers, with an application to the Entscheidungsproblem. \emph{Proceedings of the London Mathematical Society}, \emph{s2-42}(1), 230--265. \url{https://doi.org/10.1112/plms/s2-42.1.230}
\item
  \textbf{Problema.} Precisar qué significa que un número sea calculable por un procedimiento mecánico y, con esa definición, responder el \emph{Entscheidungsproblem}.
\item
  \textbf{Método.} Introduce la \emph{máquina automática}: una cinta dividida en casillas, una cabeza que lee y escribe un símbolo a la vez y un número finito de configuraciones. Un número es computable si una máquina lo escribe cifra por cifra. Construye una máquina \emph{universal} que imita a cualquier otra a partir de su descripción.
\item
  \textbf{Resultado.} Ninguna máquina decide si una máquina arbitraria escribe infinitas cifras; de ahí se sigue que no existe un procedimiento general de decisión para la lógica de primer orden. En un apéndice muestra que su noción coincide con la λ-definibilidad de Church.
\item
  \textbf{Tema de la Unidad I.} Orígenes de la teoría, computabilidad y tesis de Church-Turing.
\item
  \textbf{Aportación.} Es el origen del modelo que ocupa el nivel más alto de la jerarquía de Chomsky estudiada en el curso.
\end{enumerate}

\textbf{Resultado y modelo.} Se demuestra la indecidibilidad del \emph{Entscheidungsproblem} y se introduce la máquina de Turing, incluida la universal.

\textbf{Tesis de Church-Turing.} Toda función efectivamente calculable, en el sentido intuitivo, es calculable por una máquina de Turing. Es una tesis porque «efectivamente calculable» no es un concepto formal; sólo se puede argumentar a su favor, como hace Turing al analizar lo que una persona hace al calcular con lápiz y papel, y con la equivalencia entre modelos independientes.

\textbf{Dificultades de lectura.} La notación es la mayor barrera: letras góticas para las configuraciones, máquinas descritas como tablas de «configuración, símbolo, operaciones, configuración final» y «tablas esqueleto» que funcionan como subrutinas. Se resolvió leyendo en paralelo el capítulo 3 de \textcite{sipser2013} y reescribiendo cada tabla como δ(q, a) = (p, b, D), la notación del curso.

\hypertarget{texto-4.-mavridou-y-laszka-2018}{%
\subsection{Texto 4. Mavridou y Laszka (2018)}\label{texto-4.-mavridou-y-laszka-2018}}

\begin{enumerate}
\def\labelenumi{\arabic{enumi}.}
\tightlist
\item
  \textbf{Cita.} Mavridou, A., y Laszka, A. (2018). Designing secure Ethereum smart contracts: A finite state machine based approach. En S. Meiklejohn y K. Sako (Eds.), \emph{Financial Cryptography and Data Security} (pp.~523--540). Springer. \url{https://doi.org/10.1007/978-3-662-58387-6_28}
\item
  \textbf{Problema.} Los contratos inteligentes de Ethereum administran activos valiosos y no pueden corregirse una vez desplegados, pero muchos se escriben con errores que se han explotado para robar fondos.
\item
  \textbf{Método.} Proponen FSolidM: el desarrollador describe el contrato como una máquina de estados finitos (estados, transiciones y las acciones de cada transición) en una herramienta gráfica, y el marco genera el código Solidity. La semántica de la máquina se define formalmente y la seguridad se refuerza con patrones de diseño que se agregan como complementos.
\item
  \textbf{Resultado.} Muestran con casos como una subasta a ciegas que el modelo de estados evita por construcción errores frecuentes (llamadas en un estado indebido, reentrada) y que los complementos añaden protecciones sin reescribir la lógica.
\item
  \textbf{Tema de la Unidad I.} Autómatas finitos aplicados fuera de la teoría (seguridad y verificación de software).
\item
  \textbf{Aportación.} Ejemplo de que fijar el conjunto de estados y la función de transición antes de programar es una técnica de ingeniería, no sólo de teoría; es la misma separación entre modelo e implementación que sigue el software de esta práctica.
\end{enumerate}

\hypertarget{texto-5.-davis-et-al.-2019}{%
\subsection{Texto 5. Davis et al.~(2019)}\label{texto-5.-davis-et-al.-2019}}

\begin{enumerate}
\def\labelenumi{\arabic{enumi}.}
\tightlist
\item
  \textbf{Cita.} Davis, J. C., Michael IV, L. G., Coghlan, C. A., Servant, F., y Lee, D. (2019). Why aren't regular expressions a lingua franca? An empirical study on the re-use and portability of regular expressions. En \emph{Proceedings of the 27th ACM Joint European Software Engineering Conference and Symposium on the Foundations of Software Engineering} (pp.~443--454). ACM. \url{https://doi.org/10.1145/3338906.3338909}
\item
  \textbf{Problema.} Los programadores copian expresiones regulares de un lenguaje a otro suponiendo que significan lo mismo; nadie había medido si esa suposición es cierta.
\item
  \textbf{Método.} Encuesta a 158 desarrolladores y análisis de 537 806 expresiones únicas extraídas de 193 524 proyectos en ocho lenguajes; cada expresión se ejecuta en todos los motores con unas 10 000 entradas generadas y se buscan diferencias de resultado y de tiempo.
\item
  \textbf{Resultado.} El 94 \% de los encuestados reutiliza expresiones; alrededor del 15 \% de las expresiones produce resultados distintos en algún lenguaje y alrededor del 10 \% tiene un costo en tiempo distinto. Go y Rust no presentaron expresiones de tiempo exponencial; Java, JavaScript, Python y Ruby, más de mil cada uno.
\item
  \textbf{Tema de la Unidad I.} Expresiones regulares y su relación con los autómatas finitos.
\item
  \textbf{Aportación.} Explica por qué el curso distingue la expresión regular teórica de la de los lenguajes de programación y por qué conviene validar con un autómata cuyo tiempo es lineal.
\end{enumerate}

\hypertarget{comparaciuxf3n}{%
\subsection{Comparación}\label{comparaciuxf3n}}

\begin{longtable}[]{@{}
  >{\raggedright\arraybackslash}p{(\columnwidth - 8\tabcolsep) * \real{0.2000}}
  >{\centering\arraybackslash}p{(\columnwidth - 8\tabcolsep) * \real{0.2000}}
  >{\centering\arraybackslash}p{(\columnwidth - 8\tabcolsep) * \real{0.2000}}
  >{\raggedright\arraybackslash}p{(\columnwidth - 8\tabcolsep) * \real{0.2000}}
  >{\raggedright\arraybackslash}p{(\columnwidth - 8\tabcolsep) * \real{0.2000}}@{}}
\caption{Resumen de los cinco textos.}\tabularnewline
\toprule\noalign{}
\begin{minipage}[b]{\linewidth}\raggedright
Texto
\end{minipage} & \begin{minipage}[b]{\linewidth}\centering
Arbitrado
\end{minipage} & \begin{minipage}[b]{\linewidth}\centering
Año
\end{minipage} & \begin{minipage}[b]{\linewidth}\raggedright
Modelo que usa
\end{minipage} & \begin{minipage}[b]{\linewidth}\raggedright
Campo de aplicación
\end{minipage} \\
\midrule\noalign{}
\endfirsthead
\toprule\noalign{}
\begin{minipage}[b]{\linewidth}\raggedright
Texto
\end{minipage} & \begin{minipage}[b]{\linewidth}\centering
Arbitrado
\end{minipage} & \begin{minipage}[b]{\linewidth}\centering
Año
\end{minipage} & \begin{minipage}[b]{\linewidth}\raggedright
Modelo que usa
\end{minipage} & \begin{minipage}[b]{\linewidth}\raggedright
Campo de aplicación
\end{minipage} \\
\midrule\noalign{}
\endhead
\bottomrule\noalign{}
\endlastfoot
1. Gribkoff & No & 2013 & AFD y máquina de Mealy & Docencia: máquinas, juegos, redes, búsqueda \\
2. Luna-Benoso et al. & Sí & 2022 & Autómata celular & Diagnóstico médico por imagen \\
3. Turing (original) & Sí & 1936 & Máquina de Turing & Fundamentos de la matemática \\
4. Mavridou y Laszka & Sí & 2018 & Máquina de estados finitos & Seguridad de contratos inteligentes \\
5. Davis et al. & Sí & 2019 & Expresiones regulares y sus motores & Ingeniería de software \\
\end{longtable}

Los cinco describen un sistema mediante estados y reglas que dicen cómo cambiar de uno a otro; la diferencia está en el rigor y en el propósito. Turing demuestra un teorema de fundamentos con una definición nueva; Gribkoff, sin revisión ni bibliografía, divulga; Luna-Benoso et al.~y Mavridou y Laszka construyen un sistema y lo evalúan; Davis et al.~miden cómo se usa la teoría en la práctica. Los textos 4 y 5 se complementan: uno muestra el beneficio de fijar el modelo de estados antes de programar y el otro, el costo de no hacerlo, pues las expresiones de los lenguajes de programación ya no corresponden a un autómata finito. De ahí surge un problema abierto: cómo garantizar que una especificación escrita como autómata o expresión regular conserve su significado y su costo lineal cuando se traduce a la implementación concreta de cada lenguaje.

\hypertarget{ejercicio-4.-autuxf3matas-finitos-en-jflap}{%
\section{Ejercicio 4. Autómatas finitos en JFLAP}\label{ejercicio-4.-autuxf3matas-finitos-en-jflap}}

\hypertarget{muxe9todo-de-diseuxf1o}{%
\subsection{Método de diseño}\label{muxe9todo-de-diseuxf1o}}

Todos los autómatas se diseñaron con la misma pregunta: \emph{¿qué es lo mínimo que hay que recordar del prefijo leído para decidir el resto?} La respuesta define los estados; la función de transición dice cómo cambia ese recuerdo al leer un símbolo, y los estados de aceptación son aquellos cuyo recuerdo ya garantiza que la cadena pertenece al lenguaje. Cuatro familias de enunciados cubren casi todas las listas:

\begin{longtable}[]{@{}
  >{\raggedright\arraybackslash}p{(\columnwidth - 4\tabcolsep) * \real{0.1667}}
  >{\raggedright\arraybackslash}p{(\columnwidth - 4\tabcolsep) * \real{0.6667}}
  >{\raggedright\arraybackslash}p{(\columnwidth - 4\tabcolsep) * \real{0.1667}}@{}}
\caption{Familias de diseño.}\tabularnewline
\toprule\noalign{}
\begin{minipage}[b]{\linewidth}\raggedright
Familia
\end{minipage} & \begin{minipage}[b]{\linewidth}\raggedright
Qué guarda cada estado
\end{minipage} & \begin{minipage}[b]{\linewidth}\raggedright
Ejemplos
\end{minipage} \\
\midrule\noalign{}
\endfirsthead
\toprule\noalign{}
\begin{minipage}[b]{\linewidth}\raggedright
Familia
\end{minipage} & \begin{minipage}[b]{\linewidth}\raggedright
Qué guarda cada estado
\end{minipage} & \begin{minipage}[b]{\linewidth}\raggedright
Ejemplos
\end{minipage} \\
\midrule\noalign{}
\endhead
\bottomrule\noalign{}
\endlastfoot
Conteo módulo k & El residuo de una longitud, de un número de apariciones o del valor numérico leído & 097, 139, 145, 147 \\
Búsqueda de patrones & El sufijo más largo leído que es prefijo del patrón (construcción de Knuth-Morris-Pratt) & 099, 102, 103, 106 \\
Producto de autómatas & Un par de recuerdos, uno por condición, cuando el enunciado pide «y» o «o» & 107, 109, 133 \\
Prefijos y extremos & Cuánto del prefijo exigido se ha leído, o el primer símbolo y el último & 120, 135, 137 \\
\end{longtable}

Cuando una condición ya no puede cumplirse (por ejemplo, la cadena no empieza como se exige) se usa un estado \textbf{sumidero}: no es de aceptación y todas sus transiciones vuelven a él. En los enunciados de la Lista 3 que piden un AFN, el AFD se obtuvo convirtiendo ese AFN y minimizando el resultado; su ficha indica qué estados del AFN representa cada estado del AFD.

\hypertarget{implementaciuxf3n-y-validaciuxf3n}{%
\subsection{Implementación y validación}\label{implementaciuxf3n-y-validaciuxf3n}}

Cada autómata se capturó en JFLAP 7.1 y se exportó a \texttt{automatas/lista2/} y \texttt{automatas/lista3/}; los archivos \texttt{.jff} se leen también con el software del Ejercicio 5 de la Práctica 2. Se resolvieron los 77 ejercicios con asterisco que piden un autómata; los ejercicios 116 (lenguaje no regular), 143 (figuras ausentes en la lista), 191 y 201--235 no piden un AFD y se resuelven en las tareas.

Cada ficha del Anexo A tiene el diagrama en JFLAP, la tabla de transiciones ---en todos los autómatas, no sólo en los cinco que pide el enunciado---, tres cadenas válidas, tres inválidas y los casos frontera: la cadena vacía, la aceptada más corta y, cuando existe, una cadena justo por debajo de esa longitud. Las cadenas se eligieron de longitudes crecientes y se comprobaron dos veces: con \emph{Input → Multiple Run} en JFLAP y con un programa que compara cada autómata contra la definición del lenguaje sobre todas las cadenas de longitud ≤ 14 (≤ 9 con tres símbolos). El archivo \texttt{automatas/pruebas.json} reúne las cadenas de todas las fichas.

\hypertarget{simulaciuxf3n-paso-a-paso}{%
\subsection{Simulación paso a paso}\label{simulaciuxf3n-paso-a-paso}}

La simulación paso a paso (\emph{Input → Step with Closure}) se usó en los cinco autómatas de la tabla; cada captura muestra la configuración (estado y resto de la entrada) después de algunos pasos y el resultado final.

\begin{longtable}[]{@{}
  >{\raggedright\arraybackslash}p{(\columnwidth - 6\tabcolsep) * \real{0.1579}}
  >{\raggedright\arraybackslash}p{(\columnwidth - 6\tabcolsep) * \real{0.5263}}
  >{\raggedright\arraybackslash}p{(\columnwidth - 6\tabcolsep) * \real{0.1579}}
  >{\raggedright\arraybackslash}p{(\columnwidth - 6\tabcolsep) * \real{0.1579}}@{}}
\caption{Autómatas simulados paso a paso.}\tabularnewline
\toprule\noalign{}
\begin{minipage}[b]{\linewidth}\raggedright
Ejercicio
\end{minipage} & \begin{minipage}[b]{\linewidth}\raggedright
Lenguaje
\end{minipage} & \begin{minipage}[b]{\linewidth}\raggedright
Cadena simulada
\end{minipage} & \begin{minipage}[b]{\linewidth}\raggedright
Resultado
\end{minipage} \\
\midrule\noalign{}
\endfirsthead
\toprule\noalign{}
\begin{minipage}[b]{\linewidth}\raggedright
Ejercicio
\end{minipage} & \begin{minipage}[b]{\linewidth}\raggedright
Lenguaje
\end{minipage} & \begin{minipage}[b]{\linewidth}\raggedright
Cadena simulada
\end{minipage} & \begin{minipage}[b]{\linewidth}\raggedright
Resultado
\end{minipage} \\
\midrule\noalign{}
\endhead
\bottomrule\noalign{}
\endlastfoot
097 & Longitud múltiplo de 3 & \texttt{102211} & aceptada en r0 \\
113 & Dos apariciones de 010, con traslape & \texttt{01010} & aceptada \\
120 & Primer y último símbolo distintos & \texttt{2101} & aceptada \\
135 & bᵐabⁿ con m, n \textgreater{} 0 & \texttt{bbab} & aceptada \\
147 & Binario (izquierda a derecha) múltiplo de 6 & \texttt{1100} & aceptada (12 = 2·6) \\
\end{longtable}

\evidencia[0.95]{../evidencias/jflap/lista2-097-paso-a-paso.png}{JFLAP, simulación paso a paso de \texttt{102211} en lista2-097.}

\evidencia[0.95]{../evidencias/jflap/lista2-113-paso-a-paso.png}{JFLAP, simulación paso a paso de \texttt{01010} en lista2-113.}

\evidencia[0.95]{../evidencias/jflap/lista2-120-paso-a-paso.png}{JFLAP, simulación paso a paso de \texttt{2101} en lista2-120.}

\evidencia[0.95]{../evidencias/jflap/lista2-135-paso-a-paso.png}{JFLAP, simulación paso a paso de \texttt{bbab} en lista2-135.}

\evidencia[0.95]{../evidencias/jflap/lista2-147-paso-a-paso.png}{JFLAP, simulación paso a paso de \texttt{1100} en lista2-147.}

\hypertarget{resumen-de-los-autuxf3matas}{%
\subsection{Resumen de los autómatas}\label{resumen-de-los-autuxf3matas}}

\begin{longtable}[]{@{}lllcl@{}}
\caption{Autómatas de los ejercicios con asterisco.}\tabularnewline
\toprule\noalign{}
Ejercicio & Σ & Diseño original & Estados del AFD & Archivo \\
\midrule\noalign{}
\endfirsthead
\toprule\noalign{}
Ejercicio & Σ & Diseño original & Estados del AFD & Archivo \\
\midrule\noalign{}
\endhead
\bottomrule\noalign{}
\endlastfoot
L2 · 097 & \{0, 1, 2\} & AFD & 3 & \texttt{lista2-097.jff} \\
L2 · 099 & \{0, 1, 2\} & AFD & 7 & \texttt{lista2-099.jff} \\
L2 · 102 & \{a, b\} & AFD & 7 & \texttt{lista2-102.jff} \\
L2 · 103 & \{a, b\} & AFD & 7 & \texttt{lista2-103.jff} \\
L2 · 106 & \{0, 1, 2\} & AFD & 8 & \texttt{lista2-106.jff} \\
L2 · 107 & \{0, 1, 2\} & AFD & 13 & \texttt{lista2-107.jff} \\
L2 · 109 & \{0, 1, a, b, c\} & AFD & 22 & \texttt{lista2-109.jff} \\
L2 · 111 & \{0, 1, 2\} & AFD & 21 & \texttt{lista2-111.jff} \\
L2 · 112 & \{0, 1, 2\} & AFD & 7 & \texttt{lista2-112.jff} \\
L2 · 113 & \{0, 1, 2\} & AFD & 7 & \texttt{lista2-113.jff} \\
L2 · 118 & \{a, b\} & AFD & 6 & \texttt{lista2-118.jff} \\
L2 · 120 & \{0, 1, 2\} & AFD & 7 & \texttt{lista2-120.jff} \\
L2 · 121 & \{0, 1, 2\} & AFD & 6 & \texttt{lista2-121.jff} \\
L2 · 122 & \{0, 1, 2\} & AFD & 4 & \texttt{lista2-122.jff} \\
L2 · 123 & \{a, b\} & AFD & 15 & \texttt{lista2-123.jff} \\
L2 · 124 & \{0, 1\} & AFD & 13 & \texttt{lista2-124.jff} \\
L2 · 125 & \{0, 1\} & AFD & 2 & \texttt{lista2-125.jff} \\
L2 · 126 & \{0, 1\} & AFD & 2 & \texttt{lista2-126.jff} \\
L2 · 127 & \{0, 1\} & AFD & 2 & \texttt{lista2-127.jff} \\
L2 · 128 & \{0, 1\} & AFD & 6 & \texttt{lista2-128.jff} \\
L2 · 131 & \{0, 1\} & AFD & 2 & \texttt{lista2-131.jff} \\
L2 · 133 & \{a, b\} & AFD & 4 & \texttt{lista2-133.jff} \\
L2 · 135 & \{a, b\} & AFD & 5 & \texttt{lista2-135.jff} \\
L2 · 137 & \{a, b\} & AFD & 4 & \texttt{lista2-137.jff} \\
L2 · 139 & \{0, 1\} & AFD & 2 & \texttt{lista2-139.jff} \\
L2 · 141 & \{0, 1\} & AFD & 3 & \texttt{lista2-141.jff} \\
L2 · 145 & \{0, 1\} & AFD & 6 & \texttt{lista2-145.jff} \\
L2 · 147 & \{0, 1\} & AFD & 7 & \texttt{lista2-147.jff} \\
L2 · 148 & \{A, N, P, S\} & AFD & 13 & \texttt{lista2-148.jff} \\
L2 · 149 & \{0, 1\} & AFD & 3 & \texttt{lista2-149.jff} \\
L2 · 150 & \{0, 1, a\} & AFD & 4 & \texttt{lista2-150.jff} \\
L3 · 151 & \{0, 1\} & AFD & 10 & \texttt{lista3-151.jff} \\
L3 · 152 & \{0, 1, 2\} & AFD & 6 & \texttt{lista3-152.jff} \\
L3 · 153 & \{0, 1\} & AFD & 7 & \texttt{lista3-153.jff} \\
L3 · 154 & \{0, 1\} & AFD & 11 & \texttt{lista3-154.jff} \\
L3 · 155 & \{0, 1\} & AFD & 12 & \texttt{lista3-155.jff} \\
L3 · 156 & \{0, 1\} & AFD & 5 & \texttt{lista3-156.jff} \\
L3 · 157 & \{0, 1\} & AFD & 6 & \texttt{lista3-157.jff} \\
L3 · 158 & \{0, 1\} & AFD & 3 & \texttt{lista3-158.jff} \\
L3 · 159 & \{0, 1\} & AFD & 3 & \texttt{lista3-159.jff} \\
L3 · 160 & \{0, 1\} & AFD & 3 & \texttt{lista3-160.jff} \\
L3 · 161 & \{0, 1\} & AFD & 3 & \texttt{lista3-161.jff} \\
L3 · AFN del 097 & \{0, 1, 2\} & AFD & 3 & \texttt{lista3-afn-097.jff} \\
L3 · AFN del 099 & \{0, 1, 2\} & AFN & 7 & \texttt{lista3-afn-099.jff} \\
L3 · AFN del 102 & \{a, b\} & AFN & 7 & \texttt{lista3-afn-102.jff} \\
L3 · AFN del 103 & \{a, b\} & AFN & 7 & \texttt{lista3-afn-103.jff} \\
L3 · AFN del 106 & \{0, 1, 2\} & AFN & 8 & \texttt{lista3-afn-106.jff} \\
L3 · AFN del 107 & \{0, 1, 2\} & AFN & 13 & \texttt{lista3-afn-107.jff} \\
L3 · AFN del 109 & \{0, 1, a, b, c\} & AFN & 22 & \texttt{lista3-afn-109.jff} \\
L3 · AFN del 111 & \{0, 1, 2\} & AFD & 20 & \texttt{lista3-afn-111.jff} \\
L3 · AFN del 112 & \{0, 1, 2\} & AFN & 7 & \texttt{lista3-afn-112.jff} \\
L3 · AFN del 113 & \{0, 1, 2\} & AFN & 7 & \texttt{lista3-afn-113.jff} \\
L3 · AFN del 118 & \{a, b\} & AFN & 5 & \texttt{lista3-afn-118.jff} \\
L3 · AFN del 120 & \{0, 1, 2\} & AFN & 7 & \texttt{lista3-afn-120.jff} \\
L3 · AFN del 121 & \{0, 1, 2\} & AFD & 5 & \texttt{lista3-afn-121.jff} \\
L3 · AFN del 122 & \{0, 1, 2\} & AFD & 3 & \texttt{lista3-afn-122.jff} \\
L3 · AFN del 123 & \{a, b\} & AFN & 15 & \texttt{lista3-afn-123.jff} \\
L3 · 189 & \{0, 1\} & AFN & 7 & \texttt{lista3-189.jff} \\
L3 · 190 & \{0, 1\} & AFN & 19 & \texttt{lista3-190.jff} \\
L3 · 192 & \{0, 1\} & AFN & 5 & \texttt{lista3-192.jff} \\
L3 · 193 & \{0, 1\} & AFN & 6 & \texttt{lista3-193.jff} \\
L3 · 194 & \{0, 1, 2\} & AFN & 2 & \texttt{lista3-194.jff} \\
L3 · 195 & \{0, 1\} & AFN & 7 & \texttt{lista3-195.jff} \\
L3 · 236 & \{0, 1\} & AFN-λ & 4 & \texttt{lista3-236.jff} \\
L3 · 237 & \{0, 1\} & AFN-λ & 4 & \texttt{lista3-237.jff} \\
L3 · 238 & \{0, 1\} & AFN-λ & 6 & \texttt{lista3-238.jff} \\
L3 · 239 & \{0, 1\} & AFN-λ & 6 & \texttt{lista3-239.jff} \\
L3 · 240 & \{0, 1\} & AFN-λ & 6 & \texttt{lista3-240.jff} \\
L3 · 254 & \{0, 1\} & AFD & 2 & \texttt{lista3-254.jff} \\
L3 · 255 & \{0, 1\} & AFD & 2 & \texttt{lista3-255.jff} \\
L3 · 256 & \{0, 1, 2\} & AFD & 2 & \texttt{lista3-256.jff} \\
L3 · 257 & \{0, 1\} & AFD & 3 & \texttt{lista3-257.jff} \\
L3 · 258 & \{0, 1\} & AFD & 3 & \texttt{lista3-258.jff} \\
L3 · 260 & \{0, 1\} & AFN & 6 & \texttt{lista3-260.jff} \\
L3 · 261 & \{0, 1\} & AFN & 3 & \texttt{lista3-261.jff} \\
L3 · 262 & \{a, b\} & AFD & 5 & \texttt{lista3-262.jff} \\
L3 · 263 & \{0, 1\} & AFN & 3 & \texttt{lista3-263.jff} \\
\end{longtable}

\hypertarget{ejercicio-5.-aplicaciuxf3n-con-interfaz-gruxe1fica}{%
\section{Ejercicio 5. Aplicación con interfaz gráfica}\label{ejercicio-5.-aplicaciuxf3n-con-interfaz-gruxe1fica}}

\hypertarget{arquitectura}{%
\subsection{Arquitectura}\label{arquitectura}}

La aplicación separa el \textbf{núcleo}, que no importa Flet, de la \textbf{interfaz}, que sólo traduce clics en llamadas al núcleo. Así las pruebas ejercitan toda la lógica sin abrir una ventana.

\begin{Shaded}
\begin{Highlighting}[]
\NormalTok{src/}
\NormalTok{├── app.py                  \# punto de entrada: lista de secciones y arranque}
\NormalTok{├── nucleo/}
\NormalTok{│   ├── errores.py          \# ErrorTC y sus subclases (mensajes para el usuario)}
\NormalTok{│   └── lenguajes.py        \# Alfabeto, Palabra, Lenguaje, Enumerador, Informe}
\NormalTok{└── interfaz/}
\NormalTok{    ├── aplicacion.py       \# ventana: menú lateral, barra superior, área de trabajo}
\NormalTok{    ├── componentes.py      \# tarjetas, avisos, pastillas, métricas, ayuda numerada}
\NormalTok{    ├── tema.py · sesion.py \# paleta institucional y estado compartido}
\NormalTok{    └── vistas/             \# Inicio, Palabras, Cerraduras, Lenguajes}
\end{Highlighting}
\end{Shaded}

\begin{longtable}[]{@{}
  >{\raggedright\arraybackslash}p{(\columnwidth - 2\tabcolsep) * \real{0.2000}}
  >{\raggedright\arraybackslash}p{(\columnwidth - 2\tabcolsep) * \real{0.8000}}@{}}
\caption{Clases principales.}\tabularnewline
\toprule\noalign{}
\begin{minipage}[b]{\linewidth}\raggedright
Clase
\end{minipage} & \begin{minipage}[b]{\linewidth}\raggedright
Responsabilidad
\end{minipage} \\
\midrule\noalign{}
\endfirsthead
\toprule\noalign{}
\begin{minipage}[b]{\linewidth}\raggedright
Clase
\end{minipage} & \begin{minipage}[b]{\linewidth}\raggedright
Responsabilidad
\end{minipage} \\
\midrule\noalign{}
\endhead
\bottomrule\noalign{}
\endlastfoot
\texttt{Alfabeto} & Conjunto finito y no vacío de símbolos de un carácter; rechaza λ, comas, llaves y espacios. Responde si una palabra está en Σ\emph{, cuenta Σⁿ sin generarlo y construye Σⁿ, Σ} y Σ⁺ acotados. \\
\texttt{Palabra} & Cadena inmutable: longitud, concatenación (\texttt{+}), potencia, reflexión, prefijos, sufijos y subcadenas. \\
\texttt{Lenguaje} & Conjunto finito de palabras con \texttt{\textbar{}}, \texttt{\&}, \texttt{-}, concatenación, potencia, reflexión y cerraduras acotadas. \\
\texttt{Enumerador} & Recorre Σᵐⁱⁿ ∪ \ldots{} ∪ Σᵐᵃˣ en orden canónico y verifica el límite \emph{antes} de generar. \\
\texttt{Informe} & Texto exportable con encabezado, datos de la operación, total y resultados. \\
\texttt{Vista} (interfaz) & Base de cada sección: encabezado, ayuda «¿Cómo se usa?», avisos y manejo uniforme de errores. \\
\end{longtable}

Todas las operaciones del núcleo devuelven objetos nuevos; ninguna modifica sus argumentos. Los errores de dominio (símbolo inválido, palabra fuera de Σ, límite excedido) son subclases de \texttt{ErrorTC} y la interfaz los muestra en una franja roja con el mensaje, sin interrumpir el programa.

\hypertarget{diseuxf1o-de-la-interfaz}{%
\subsection{Diseño de la interfaz}\label{diseuxf1o-de-la-interfaz}}

La ventana tiene un menú lateral con las secciones y un área de trabajo. Cada sección empieza con un encabezado que dice qué hace y un botón \textbf{¿Cómo se usa?} que despliega los pasos numerados. Los resultados se calculan \textbf{mientras se escribe}: no hay que pulsar un botón para ver los prefijos de una palabra, y un símbolo que no pertenece a Σ se señala en el propio campo. Los resultados se muestran como pastillas (la palabra vacía, resaltada) acompañadas de métricas: la longitud, el número de resultados y el máximo teórico.

\hypertarget{operaciuxf3n-1.-prefijos-sufijos-y-subcadenas}{%
\subsection{Operación 1. Prefijos, sufijos y subcadenas}\label{operaciuxf3n-1.-prefijos-sufijos-y-subcadenas}}

\texttt{Palabra.prefijos()} devuelve \texttt{{[}ω{[}:0{]},\ ω{[}:1{]},\ …,\ ω{[}:n{]}{]}}; \texttt{sufijos()}, los n + 1 finales; y \texttt{subcadenas()} recorre todas las longitudes y posiciones y conserva cada subcadena la primera vez que aparece. La sección \textbf{Palabras} permite además la reflexión y la potencia, y guarda cualquier resultado en \texttt{salidas/palabras-\textless{}operación\textgreater{}.txt}.

\textbf{Conteos.} Una palabra de longitud n tiene exactamente \textbf{n + 1 prefijos} (uno por cada punto de corte, de 0 a n) y n + 1 sufijos. Las subcadenas no vacías se determinan por un par de posiciones 0 ≤ i \textless{} j ≤ n, de modo que hay a lo sumo n(n + 1)/2; con λ, el máximo es \textbf{n(n + 1)/2 + 1}, y se alcanza sólo cuando todos los símbolos son distintos. Para \texttt{abcb} hay 5 prefijos y 10 subcadenas distintas, una menos que la cota 11 porque \texttt{b} se repite.

\textbf{Convención sobre λ.} El programa incluye λ y la propia palabra entre los prefijos y los sufijos. Se sigue la definición: x es prefijo de ω si existe y tal que xy = ω; con y = ω se obtiene x = λ, y con y = λ, x = ω. Excluirlas haría que una palabra de longitud n tuviera n − 1 prefijos \emph{propios}, que es otro concepto; el núcleo ofrece esa variante con \texttt{prefijos(propios=True)}.

\hypertarget{operaciuxf3n-2.-cerradura-de-kleene-y-cerradura-positiva}{%
\subsection{Operación 2. Cerradura de Kleene y cerradura positiva}\label{operaciuxf3n-2.-cerradura-de-kleene-y-cerradura-positiva}}

La sección \textbf{Cerraduras} genera Σ*, Σ⁺ o Σⁿ hasta la longitud que el usuario elige con un control deslizante. Antes de generar, \texttt{Alfabeto.cuantas} calcula el total con la fórmula

\[|\Sigma^0 \cup \Sigma^1 \cup \dots \cup \Sigma^n| = \sum_{i=0}^{n} |\Sigma|^i = \frac{|\Sigma|^{n+1} - 1}{|\Sigma| - 1} \quad (|\Sigma| > 1),\]

y la interfaz muestra ese número y una barra respecto del límite \textbf{mientras se mueve el control}. Si el total supera 200 000, el botón \emph{Generar} se desactiva y el mensaje indica la longitud máxima permitida para ese alfabeto; el núcleo repite la verificación y lanza \texttt{LimiteExcedido} aunque se le llame desde otro programa.

\begin{Shaded}
\begin{Highlighting}[]
\VariableTok{self}\NormalTok{.total }\OperatorTok{=}\NormalTok{ alfabeto.cuantas(maximo, minimo)}
\ControlFlowTok{if} \VariableTok{self}\NormalTok{.total }\OperatorTok{\textgreater{}}\NormalTok{ limite:}
    \ControlFlowTok{raise}\NormalTok{ LimiteExcedido(}\SpecialStringTok{f"... produce }\SpecialCharTok{\{}\VariableTok{self}\SpecialCharTok{.}\NormalTok{total}\SpecialCharTok{\}}\SpecialStringTok{ cadenas; el límite es }\SpecialCharTok{\{}\NormalTok{limite}\SpecialCharTok{\}}\SpecialStringTok{."}\NormalTok{)}
\end{Highlighting}
\end{Shaded}

\textbf{Relación con la unidad.} Σ* es infinito para todo alfabeto, así que ningún programa puede «calcularlo»; lo que se genera es un truncamiento finito. Cada nivel tiene \textbar Σⁿ\textbar{} = \textbar Σ\textbar ⁿ cadenas, de modo que el truncamiento crece exponencialmente en la longitud: con \textbar Σ\textbar{} = 2 se llega a n = 16 (131 071 cadenas), pero con \textbar Σ\textbar{} = 10 sólo a n = 5. El límite convierte esa propiedad teórica en una regla que el usuario ve antes de esperar. Σ⁺ con longitud 0 es vacío: la única cadena de longitud 0 es λ, y Σ⁺ la excluye.

\evidencia[0.95]{../evidencias/app/01-palabras.png}{Sección Palabras: subcadenas de \texttt{abcb} con las métricas y la cota
teórica.}

\evidencia[0.95]{../evidencias/app/02-cerraduras.png}{Sección Cerraduras: Σ* con Σ = \{a, b\} y n = 3, agrupada por longitud.}

\evidencia[0.95]{../evidencias/app/03-limite.png}{Con Σ = \{0, 1\} y n = 17 el contador supera 200 000 y el botón Generar
se desactiva.}

\hypertarget{pruebas-automatizadas}{%
\subsection{Pruebas automatizadas}\label{pruebas-automatizadas}}

\texttt{tests/test\_lenguajes.py} (34 pruebas) cubre lo que pide el enunciado y algo más: la palabra vacía (λ = «» = «ε», sus prefijos, sufijos y subcadenas), un alfabeto de un solo símbolo, los prefijos y sufijos de una palabra de longitud 1, la diferencia Σ* − Σ⁺ = \{λ\} y Σ⁺ = ∅ para longitud 0, los conteos n + 1 y n(n + 1)/2 + 1, el orden canónico, el límite exacto (n = 16 se permite y n = 17 no, con \textbar Σ\textbar{} = 2) y las operaciones entre lenguajes. \texttt{tests/test\_interfaz.py} construye cada sección sin ventana y ejecuta sus acciones principales, por ejemplo escribir una palabra con un símbolo ajeno y comprobar que el campo muestra el error.

\begin{longtable}[]{@{}lll@{}}
\caption{Resultados de la suite en los tres entornos.}\tabularnewline
\toprule\noalign{}
Servicio & Intérprete & Resultado de \texttt{pytest\ -q} \\
\midrule\noalign{}
\endfirsthead
\toprule\noalign{}
Servicio & Intérprete & Resultado de \texttt{pytest\ -q} \\
\midrule\noalign{}
\endhead
\bottomrule\noalign{}
\endlastfoot
\texttt{py311} & Python 3.11 & 41 aprobadas, 3 omitidas \\
\texttt{py312} & Python 3.12 & 41 aprobadas, 3 omitidas \\
\texttt{py313} & Python 3.13 & 41 aprobadas, 3 omitidas \\
\end{longtable}

Las tres pruebas omitidas pertenecen a secciones de prácticas posteriores (editor, algoritmos y expresiones) y se saltan porque esta práctica no incluye esos módulos. No se observaron diferencias de comportamiento entre versiones.

\evidencia[0.95]{../evidencias/docker/pytest-tres-versiones.png}{\texttt{docker\ compose\ run\ -\/-rm\ py311\ pytest\ -q} y sus
equivalentes en 3.12 y 3.13.}

\evidencia[0.95]{../evidencias/app/04-contenedor.png}{Interfaz servida por el contenedor \texttt{py312} en
http://localhost:8550.}

\hypertarget{conclusiones}{%
\section{Conclusiones}\label{conclusiones}}

El entorno de trabajo quedó resuelto antes de escribir la primera línea de lógica: un repositorio con historia, ramas y Pull Requests, y tres contenedores construidos con un único \texttt{Dockerfile} que difieren sólo en la versión del intérprete. Fijar la imagen y las dependencias tuvo un efecto concreto: la incompatibilidad entre \texttt{flet-web} y las versiones nuevas de FastAPI apareció una sola vez, se resolvió con una línea en \texttt{requirements.txt} y no volvió a afectar a ninguno de los tres entornos.

La investigación mostró que las definiciones de la unidad no son convenciones arbitrarias. Que Σ⁰ sea \{λ\} y no ∅ es lo que hace funcionar las leyes de las potencias, y la diferencia entre Σ* y Σ⁺ se reduce a λ. La lectura de Turing explicó por qué la tesis de Church-Turing no puede demostrarse, y los artículos recientes, por qué la distinción entre la expresión regular de la teoría y la de los lenguajes de programación tiene consecuencias de seguridad.

En JFLAP, preguntarse qué hay que recordar del prefijo leído redujo casi todos los enunciados a cuatro familias de diseño (conteo modular, búsqueda de patrones, producto y extremos), y comparar cada autómata contra la definición del lenguaje sobre todas las cadenas cortas detectó errores que tres cadenas de prueba no habrían mostrado.

En la aplicación, separar el núcleo de la interfaz permitió probar toda la lógica con 34 pruebas en tres versiones de Python sin abrir una ventana, y calcular el número de cadenas con la fórmula antes de generarlas convirtió el crecimiento exponencial de \textbar Σ\textbar ⁿ en un aviso que el usuario ve mientras elige la longitud.


\clearpage
\printbibliography[title={Referencias}]

\clearpage
\appendix
\addtocontents{toc}{\protect\setcounter{tocdepth}{1}}
\renewcommand{\thesection}{\Alph{section}}
\begingroup
\singlespacing\small
\setlength{\LTpre}{0.3em}\setlength{\LTpost}{0.3em}
\hypertarget{anexo-a.-fichas-de-diseuxf1o-de-los-autuxf3matas}{%
\section{Anexo A. Fichas de diseño de los autómatas}\label{anexo-a.-fichas-de-diseuxf1o-de-los-autuxf3matas}}

Cada ficha reúne el enunciado, la idea del diseño, lo que recuerda cada estado, la tabla de transiciones, las cadenas de prueba (tres válidas, tres inválidas y los casos frontera) y la captura de JFLAP con el \emph{Multiple Run} de las tres válidas y las tres inválidas. Cuando el diseño natural del enunciado es no determinista, la ficha presenta el AFD que se obtuvo con la construcción de subconjuntos y la minimización, e indica qué estados del AFN representa cada estado del AFD. El archivo de cada autómata está en \texttt{automatas/lista2/} o \texttt{automatas/lista3/} con el nombre de la ficha.

\hypertarget{lista-2}{%
\subsection{Lista 2}\label{lista-2}}

\hypertarget{lista-2-ejercicio-097}{%
\subsubsection{Lista 2 · Ejercicio 097}\label{lista-2-ejercicio-097}}

\textbf{Enunciado.} Las palabras de Σ1 cuya longitud es un múltiplo de 3.

\textbf{Diseño.} Basta recordar la longitud módulo 3; λ tiene longitud 0, que es múltiplo de 3.

\textbf{Archivo.} \texttt{automatas/lista2/lista2-097.jff} · Σ = \{0, 1, 2\} · \textbar Q\textbar{} = 3

\begin{longtable}[]{@{}lllll@{}}
\caption{Tabla de transiciones (→ inicial, * aceptación, --- sin transición).}\tabularnewline
\toprule\noalign{}
δ & 0 & 1 & 2 & Qué recuerda del prefijo leído \\
\midrule\noalign{}
\endfirsthead
\toprule\noalign{}
δ & 0 & 1 & 2 & Qué recuerda del prefijo leído \\
\midrule\noalign{}
\endhead
\bottomrule\noalign{}
\endlastfoot
→* r0 & r1 & r1 & r1 & la longitud leída es ≡ 0 (mod 3) \\
r1 & r2 & r2 & r2 & la longitud leída es ≡ 1 (mod 3) \\
r2 & r0 & r0 & r0 & la longitud leída es ≡ 2 (mod 3) \\
\end{longtable}

\begin{longtable}[]{@{}lclll@{}}
\caption{Cadenas de prueba.}\tabularnewline
\toprule\noalign{}
Cadena & Longitud & Esperado & Obtenido & Tipo de caso \\
\midrule\noalign{}
\endfirsthead
\toprule\noalign{}
Cadena & Longitud & Esperado & Obtenido & Tipo de caso \\
\midrule\noalign{}
\endhead
\bottomrule\noalign{}
\endlastfoot
λ & 0 & acepta & acepta & válida \\
\texttt{212200} & 6 & acepta & acepta & válida \\
\texttt{100022002} & 9 & acepta & acepta & válida \\
\texttt{1} & 1 & rechaza & rechaza & inválida \\
\texttt{20} & 2 & rechaza & rechaza & inválida \\
\texttt{0100} & 4 & rechaza & rechaza & inválida \\
\end{longtable}

\evidenciaaqui[0.95]{../evidencias/jflap/fichas/lista2-097.png}{lista2-097 en JFLAP: Multiple Run con las tres cadenas válidas y las
tres inválidas de la tabla anterior.}

\hypertarget{lista-2-ejercicio-099}{%
\subsubsection{Lista 2 · Ejercicio 099}\label{lista-2-ejercicio-099}}

\textbf{Enunciado.} Las palabras de Σ1 que terminan con la cadena 001001.

\textbf{Diseño.} Autómata de búsqueda de patrones (Knuth-Morris-Pratt): cada estado guarda el sufijo más largo leído que es prefijo de 001001; se acepta cuando ese sufijo es el patrón completo.

\textbf{Archivo.} \texttt{automatas/lista2/lista2-099.jff} · Σ = \{0, 1, 2\} · \textbar Q\textbar{} = 7

\begin{longtable}[]{@{}
  >{\raggedright\arraybackslash}p{(\columnwidth - 8\tabcolsep) * \real{0.1364}}
  >{\raggedright\arraybackslash}p{(\columnwidth - 8\tabcolsep) * \real{0.1364}}
  >{\raggedright\arraybackslash}p{(\columnwidth - 8\tabcolsep) * \real{0.1364}}
  >{\raggedright\arraybackslash}p{(\columnwidth - 8\tabcolsep) * \real{0.1364}}
  >{\raggedright\arraybackslash}p{(\columnwidth - 8\tabcolsep) * \real{0.4545}}@{}}
\caption{Tabla de transiciones (→ inicial, * aceptación, --- sin transición).}\tabularnewline
\toprule\noalign{}
\begin{minipage}[b]{\linewidth}\raggedright
δ
\end{minipage} & \begin{minipage}[b]{\linewidth}\raggedright
0
\end{minipage} & \begin{minipage}[b]{\linewidth}\raggedright
1
\end{minipage} & \begin{minipage}[b]{\linewidth}\raggedright
2
\end{minipage} & \begin{minipage}[b]{\linewidth}\raggedright
Qué recuerda del prefijo leído
\end{minipage} \\
\midrule\noalign{}
\endfirsthead
\toprule\noalign{}
\begin{minipage}[b]{\linewidth}\raggedright
δ
\end{minipage} & \begin{minipage}[b]{\linewidth}\raggedright
0
\end{minipage} & \begin{minipage}[b]{\linewidth}\raggedright
1
\end{minipage} & \begin{minipage}[b]{\linewidth}\raggedright
2
\end{minipage} & \begin{minipage}[b]{\linewidth}\raggedright
Qué recuerda del prefijo leído
\end{minipage} \\
\midrule\noalign{}
\endhead
\bottomrule\noalign{}
\endlastfoot
→ λ & 0 & λ & λ & el sufijo más largo leído que es prefijo de 001001 es λ \\
0 & 00 & λ & λ & el sufijo más largo leído que es prefijo de 001001 es 0 \\
00 & 00 & 001 & λ & el sufijo más largo leído que es prefijo de 001001 es 00 \\
001 & 0010 & λ & λ & el sufijo más largo leído que es prefijo de 001001 es 001 \\
0010 & 00100 & λ & λ & el sufijo más largo leído que es prefijo de 001001 es 0010 \\
00100 & 00 & 001001 & λ & el sufijo más largo leído que es prefijo de 001001 es 00100 \\
* 001001 & 0010 & λ & λ & el sufijo más largo leído que es prefijo de 001001 es 001001 (la cadena termina en el patrón) \\
\end{longtable}

\begin{longtable}[]{@{}
  >{\raggedright\arraybackslash}p{(\columnwidth - 8\tabcolsep) * \real{0.2000}}
  >{\centering\arraybackslash}p{(\columnwidth - 8\tabcolsep) * \real{0.2000}}
  >{\raggedright\arraybackslash}p{(\columnwidth - 8\tabcolsep) * \real{0.2000}}
  >{\raggedright\arraybackslash}p{(\columnwidth - 8\tabcolsep) * \real{0.2000}}
  >{\raggedright\arraybackslash}p{(\columnwidth - 8\tabcolsep) * \real{0.2000}}@{}}
\caption{Cadenas de prueba.}\tabularnewline
\toprule\noalign{}
\begin{minipage}[b]{\linewidth}\raggedright
Cadena
\end{minipage} & \begin{minipage}[b]{\linewidth}\centering
Longitud
\end{minipage} & \begin{minipage}[b]{\linewidth}\raggedright
Esperado
\end{minipage} & \begin{minipage}[b]{\linewidth}\raggedright
Obtenido
\end{minipage} & \begin{minipage}[b]{\linewidth}\raggedright
Tipo de caso
\end{minipage} \\
\midrule\noalign{}
\endfirsthead
\toprule\noalign{}
\begin{minipage}[b]{\linewidth}\raggedright
Cadena
\end{minipage} & \begin{minipage}[b]{\linewidth}\centering
Longitud
\end{minipage} & \begin{minipage}[b]{\linewidth}\raggedright
Esperado
\end{minipage} & \begin{minipage}[b]{\linewidth}\raggedright
Obtenido
\end{minipage} & \begin{minipage}[b]{\linewidth}\raggedright
Tipo de caso
\end{minipage} \\
\midrule\noalign{}
\endhead
\bottomrule\noalign{}
\endlastfoot
\texttt{001001} & 6 & acepta & acepta & válida \\
\texttt{1001001} & 7 & acepta & acepta & válida \\
\texttt{120001001} & 9 & acepta & acepta & válida \\
λ & 0 & rechaza & rechaza & inválida \\
\texttt{1} & 1 & rechaza & rechaza & inválida \\
\texttt{002} & 3 & rechaza & rechaza & inválida \\
\texttt{11111} & 5 & rechaza & rechaza & frontera: justo por debajo de la longitud mínima \\
\end{longtable}

\evidenciaaqui[0.95]{../evidencias/jflap/fichas/lista2-099.png}{lista2-099 en JFLAP: Multiple Run con las tres cadenas válidas y las
tres inválidas de la tabla anterior.}

\hypertarget{lista-2-ejercicio-102}{%
\subsubsection{Lista 2 · Ejercicio 102}\label{lista-2-ejercicio-102}}

\textbf{Enunciado.} Las palabras de Σ3 que contienen la cadena abbbab.

\textbf{Diseño.} Búsqueda de abbbab con el autómata de Knuth-Morris-Pratt; al completar el patrón el estado es absorbente.

\textbf{Archivo.} \texttt{automatas/lista2/lista2-102.jff} · Σ = \{a, b\} · \textbar Q\textbar{} = 7

\begin{longtable}[]{@{}
  >{\raggedright\arraybackslash}p{(\columnwidth - 6\tabcolsep) * \real{0.1579}}
  >{\raggedright\arraybackslash}p{(\columnwidth - 6\tabcolsep) * \real{0.1579}}
  >{\raggedright\arraybackslash}p{(\columnwidth - 6\tabcolsep) * \real{0.1579}}
  >{\raggedright\arraybackslash}p{(\columnwidth - 6\tabcolsep) * \real{0.5263}}@{}}
\caption{Tabla de transiciones (→ inicial, * aceptación, --- sin transición).}\tabularnewline
\toprule\noalign{}
\begin{minipage}[b]{\linewidth}\raggedright
δ
\end{minipage} & \begin{minipage}[b]{\linewidth}\raggedright
a
\end{minipage} & \begin{minipage}[b]{\linewidth}\raggedright
b
\end{minipage} & \begin{minipage}[b]{\linewidth}\raggedright
Qué recuerda del prefijo leído
\end{minipage} \\
\midrule\noalign{}
\endfirsthead
\toprule\noalign{}
\begin{minipage}[b]{\linewidth}\raggedright
δ
\end{minipage} & \begin{minipage}[b]{\linewidth}\raggedright
a
\end{minipage} & \begin{minipage}[b]{\linewidth}\raggedright
b
\end{minipage} & \begin{minipage}[b]{\linewidth}\raggedright
Qué recuerda del prefijo leído
\end{minipage} \\
\midrule\noalign{}
\endhead
\bottomrule\noalign{}
\endlastfoot
→ λ & a & λ & aún no aparece abbbab; el sufijo más largo leído que es prefijo de abbbab es λ \\
a & a & ab & aún no aparece abbbab; el sufijo más largo leído que es prefijo de abbbab es a \\
ab & a & abb & aún no aparece abbbab; el sufijo más largo leído que es prefijo de abbbab es ab \\
abb & a & abbb & aún no aparece abbbab; el sufijo más largo leído que es prefijo de abbbab es abb \\
abbb & abbba & λ & aún no aparece abbbab; el sufijo más largo leído que es prefijo de abbbab es abbb \\
abbba & a & abbbab & aún no aparece abbbab; el sufijo más largo leído que es prefijo de abbbab es abbba \\
* abbbab & abbbab & abbbab & ya apareció abbbab \\
\end{longtable}

\begin{longtable}[]{@{}
  >{\raggedright\arraybackslash}p{(\columnwidth - 8\tabcolsep) * \real{0.2000}}
  >{\centering\arraybackslash}p{(\columnwidth - 8\tabcolsep) * \real{0.2000}}
  >{\raggedright\arraybackslash}p{(\columnwidth - 8\tabcolsep) * \real{0.2000}}
  >{\raggedright\arraybackslash}p{(\columnwidth - 8\tabcolsep) * \real{0.2000}}
  >{\raggedright\arraybackslash}p{(\columnwidth - 8\tabcolsep) * \real{0.2000}}@{}}
\caption{Cadenas de prueba.}\tabularnewline
\toprule\noalign{}
\begin{minipage}[b]{\linewidth}\raggedright
Cadena
\end{minipage} & \begin{minipage}[b]{\linewidth}\centering
Longitud
\end{minipage} & \begin{minipage}[b]{\linewidth}\raggedright
Esperado
\end{minipage} & \begin{minipage}[b]{\linewidth}\raggedright
Obtenido
\end{minipage} & \begin{minipage}[b]{\linewidth}\raggedright
Tipo de caso
\end{minipage} \\
\midrule\noalign{}
\endfirsthead
\toprule\noalign{}
\begin{minipage}[b]{\linewidth}\raggedright
Cadena
\end{minipage} & \begin{minipage}[b]{\linewidth}\centering
Longitud
\end{minipage} & \begin{minipage}[b]{\linewidth}\raggedright
Esperado
\end{minipage} & \begin{minipage}[b]{\linewidth}\raggedright
Obtenido
\end{minipage} & \begin{minipage}[b]{\linewidth}\raggedright
Tipo de caso
\end{minipage} \\
\midrule\noalign{}
\endhead
\bottomrule\noalign{}
\endlastfoot
\texttt{abbbab} & 6 & acepta & acepta & válida \\
\texttt{abbbaba} & 7 & acepta & acepta & válida \\
\texttt{aabbbabab} & 9 & acepta & acepta & válida \\
λ & 0 & rechaza & rechaza & inválida \\
\texttt{b} & 1 & rechaza & rechaza & inválida \\
\texttt{aab} & 3 & rechaza & rechaza & inválida \\
\texttt{baaab} & 5 & rechaza & rechaza & frontera: justo por debajo de la longitud mínima \\
\end{longtable}

\evidenciaaqui[0.95]{../evidencias/jflap/fichas/lista2-102.png}{lista2-102 en JFLAP: Multiple Run con las tres cadenas válidas y las
tres inválidas de la tabla anterior.}

\hypertarget{lista-2-ejercicio-103}{%
\subsubsection{Lista 2 · Ejercicio 103}\label{lista-2-ejercicio-103}}

\textbf{Enunciado.} Las palabras de Σ3 que contienen la cadena aabbba.

\textbf{Diseño.} Búsqueda de aabbba con el autómata de Knuth-Morris-Pratt; al completar el patrón el estado es absorbente.

\textbf{Archivo.} \texttt{automatas/lista2/lista2-103.jff} · Σ = \{a, b\} · \textbar Q\textbar{} = 7

\begin{longtable}[]{@{}
  >{\raggedright\arraybackslash}p{(\columnwidth - 6\tabcolsep) * \real{0.1579}}
  >{\raggedright\arraybackslash}p{(\columnwidth - 6\tabcolsep) * \real{0.1579}}
  >{\raggedright\arraybackslash}p{(\columnwidth - 6\tabcolsep) * \real{0.1579}}
  >{\raggedright\arraybackslash}p{(\columnwidth - 6\tabcolsep) * \real{0.5263}}@{}}
\caption{Tabla de transiciones (→ inicial, * aceptación, --- sin transición).}\tabularnewline
\toprule\noalign{}
\begin{minipage}[b]{\linewidth}\raggedright
δ
\end{minipage} & \begin{minipage}[b]{\linewidth}\raggedright
a
\end{minipage} & \begin{minipage}[b]{\linewidth}\raggedright
b
\end{minipage} & \begin{minipage}[b]{\linewidth}\raggedright
Qué recuerda del prefijo leído
\end{minipage} \\
\midrule\noalign{}
\endfirsthead
\toprule\noalign{}
\begin{minipage}[b]{\linewidth}\raggedright
δ
\end{minipage} & \begin{minipage}[b]{\linewidth}\raggedright
a
\end{minipage} & \begin{minipage}[b]{\linewidth}\raggedright
b
\end{minipage} & \begin{minipage}[b]{\linewidth}\raggedright
Qué recuerda del prefijo leído
\end{minipage} \\
\midrule\noalign{}
\endhead
\bottomrule\noalign{}
\endlastfoot
→ λ & a & λ & aún no aparece aabbba; el sufijo más largo leído que es prefijo de aabbba es λ \\
a & aa & λ & aún no aparece aabbba; el sufijo más largo leído que es prefijo de aabbba es a \\
aa & aa & aab & aún no aparece aabbba; el sufijo más largo leído que es prefijo de aabbba es aa \\
aab & a & aabb & aún no aparece aabbba; el sufijo más largo leído que es prefijo de aabbba es aab \\
aabb & a & aabbb & aún no aparece aabbba; el sufijo más largo leído que es prefijo de aabbba es aabb \\
aabbb & aabbba & λ & aún no aparece aabbba; el sufijo más largo leído que es prefijo de aabbba es aabbb \\
* aabbba & aabbba & aabbba & ya apareció aabbba \\
\end{longtable}

\begin{longtable}[]{@{}
  >{\raggedright\arraybackslash}p{(\columnwidth - 8\tabcolsep) * \real{0.2000}}
  >{\centering\arraybackslash}p{(\columnwidth - 8\tabcolsep) * \real{0.2000}}
  >{\raggedright\arraybackslash}p{(\columnwidth - 8\tabcolsep) * \real{0.2000}}
  >{\raggedright\arraybackslash}p{(\columnwidth - 8\tabcolsep) * \real{0.2000}}
  >{\raggedright\arraybackslash}p{(\columnwidth - 8\tabcolsep) * \real{0.2000}}@{}}
\caption{Cadenas de prueba.}\tabularnewline
\toprule\noalign{}
\begin{minipage}[b]{\linewidth}\raggedright
Cadena
\end{minipage} & \begin{minipage}[b]{\linewidth}\centering
Longitud
\end{minipage} & \begin{minipage}[b]{\linewidth}\raggedright
Esperado
\end{minipage} & \begin{minipage}[b]{\linewidth}\raggedright
Obtenido
\end{minipage} & \begin{minipage}[b]{\linewidth}\raggedright
Tipo de caso
\end{minipage} \\
\midrule\noalign{}
\endfirsthead
\toprule\noalign{}
\begin{minipage}[b]{\linewidth}\raggedright
Cadena
\end{minipage} & \begin{minipage}[b]{\linewidth}\centering
Longitud
\end{minipage} & \begin{minipage}[b]{\linewidth}\raggedright
Esperado
\end{minipage} & \begin{minipage}[b]{\linewidth}\raggedright
Obtenido
\end{minipage} & \begin{minipage}[b]{\linewidth}\raggedright
Tipo de caso
\end{minipage} \\
\midrule\noalign{}
\endhead
\bottomrule\noalign{}
\endlastfoot
\texttt{aabbba} & 6 & acepta & acepta & válida \\
\texttt{aabbbab} & 7 & acepta & acepta & válida \\
\texttt{aaaabbbab} & 9 & acepta & acepta & válida \\
λ & 0 & rechaza & rechaza & inválida \\
\texttt{a} & 1 & rechaza & rechaza & inválida \\
\texttt{aab} & 3 & rechaza & rechaza & inválida \\
\texttt{aabab} & 5 & rechaza & rechaza & frontera: justo por debajo de la longitud mínima \\
\end{longtable}

\evidenciaaqui[0.95]{../evidencias/jflap/fichas/lista2-103.png}{lista2-103 en JFLAP: Multiple Run con las tres cadenas válidas y las
tres inválidas de la tabla anterior.}

\hypertarget{lista-2-ejercicio-106}{%
\subsubsection{Lista 2 · Ejercicio 106}\label{lista-2-ejercicio-106}}

\textbf{Enunciado.} Las palabras de Σ1 que incluyen la subcadena 1101001.

\textbf{Diseño.} Búsqueda de 1101001 con el autómata de Knuth-Morris-Pratt; al completar el patrón el estado es absorbente.

\textbf{Archivo.} \texttt{automatas/lista2/lista2-106.jff} · Σ = \{0, 1, 2\} · \textbar Q\textbar{} = 8

\begin{longtable}[]{@{}
  >{\raggedright\arraybackslash}p{(\columnwidth - 8\tabcolsep) * \real{0.1364}}
  >{\raggedright\arraybackslash}p{(\columnwidth - 8\tabcolsep) * \real{0.1364}}
  >{\raggedright\arraybackslash}p{(\columnwidth - 8\tabcolsep) * \real{0.1364}}
  >{\raggedright\arraybackslash}p{(\columnwidth - 8\tabcolsep) * \real{0.1364}}
  >{\raggedright\arraybackslash}p{(\columnwidth - 8\tabcolsep) * \real{0.4545}}@{}}
\caption{Tabla de transiciones (→ inicial, * aceptación, --- sin transición).}\tabularnewline
\toprule\noalign{}
\begin{minipage}[b]{\linewidth}\raggedright
δ
\end{minipage} & \begin{minipage}[b]{\linewidth}\raggedright
0
\end{minipage} & \begin{minipage}[b]{\linewidth}\raggedright
1
\end{minipage} & \begin{minipage}[b]{\linewidth}\raggedright
2
\end{minipage} & \begin{minipage}[b]{\linewidth}\raggedright
Qué recuerda del prefijo leído
\end{minipage} \\
\midrule\noalign{}
\endfirsthead
\toprule\noalign{}
\begin{minipage}[b]{\linewidth}\raggedright
δ
\end{minipage} & \begin{minipage}[b]{\linewidth}\raggedright
0
\end{minipage} & \begin{minipage}[b]{\linewidth}\raggedright
1
\end{minipage} & \begin{minipage}[b]{\linewidth}\raggedright
2
\end{minipage} & \begin{minipage}[b]{\linewidth}\raggedright
Qué recuerda del prefijo leído
\end{minipage} \\
\midrule\noalign{}
\endhead
\bottomrule\noalign{}
\endlastfoot
→ λ & λ & 1 & λ & aún no aparece 1101001; el sufijo más largo leído que es prefijo de 1101001 es λ \\
1 & λ & 11 & λ & aún no aparece 1101001; el sufijo más largo leído que es prefijo de 1101001 es 1 \\
11 & 110 & 11 & λ & aún no aparece 1101001; el sufijo más largo leído que es prefijo de 1101001 es 11 \\
110 & λ & 1101 & λ & aún no aparece 1101001; el sufijo más largo leído que es prefijo de 1101001 es 110 \\
1101 & 11010 & 11 & λ & aún no aparece 1101001; el sufijo más largo leído que es prefijo de 1101001 es 1101 \\
11010 & 110100 & 1 & λ & aún no aparece 1101001; el sufijo más largo leído que es prefijo de 1101001 es 11010 \\
110100 & λ & 1101001 & λ & aún no aparece 1101001; el sufijo más largo leído que es prefijo de 1101001 es 110100 \\
* 1101001 & 1101001 & 1101001 & 1101001 & ya apareció 1101001 \\
\end{longtable}

\begin{longtable}[]{@{}
  >{\raggedright\arraybackslash}p{(\columnwidth - 8\tabcolsep) * \real{0.2000}}
  >{\centering\arraybackslash}p{(\columnwidth - 8\tabcolsep) * \real{0.2000}}
  >{\raggedright\arraybackslash}p{(\columnwidth - 8\tabcolsep) * \real{0.2000}}
  >{\raggedright\arraybackslash}p{(\columnwidth - 8\tabcolsep) * \real{0.2000}}
  >{\raggedright\arraybackslash}p{(\columnwidth - 8\tabcolsep) * \real{0.2000}}@{}}
\caption{Cadenas de prueba.}\tabularnewline
\toprule\noalign{}
\begin{minipage}[b]{\linewidth}\raggedright
Cadena
\end{minipage} & \begin{minipage}[b]{\linewidth}\centering
Longitud
\end{minipage} & \begin{minipage}[b]{\linewidth}\raggedright
Esperado
\end{minipage} & \begin{minipage}[b]{\linewidth}\raggedright
Obtenido
\end{minipage} & \begin{minipage}[b]{\linewidth}\raggedright
Tipo de caso
\end{minipage} \\
\midrule\noalign{}
\endfirsthead
\toprule\noalign{}
\begin{minipage}[b]{\linewidth}\raggedright
Cadena
\end{minipage} & \begin{minipage}[b]{\linewidth}\centering
Longitud
\end{minipage} & \begin{minipage}[b]{\linewidth}\raggedright
Esperado
\end{minipage} & \begin{minipage}[b]{\linewidth}\raggedright
Obtenido
\end{minipage} & \begin{minipage}[b]{\linewidth}\raggedright
Tipo de caso
\end{minipage} \\
\midrule\noalign{}
\endhead
\bottomrule\noalign{}
\endlastfoot
\texttt{1101001} & 7 & acepta & acepta & válida \\
\texttt{11010010} & 8 & acepta & acepta & válida \\
\texttt{0021101001} & 10 & acepta & acepta & válida \\
λ & 0 & rechaza & rechaza & inválida \\
\texttt{0} & 1 & rechaza & rechaza & inválida \\
\texttt{122} & 3 & rechaza & rechaza & inválida \\
\texttt{101012} & 6 & rechaza & rechaza & frontera: justo por debajo de la longitud mínima \\
\end{longtable}

\evidenciaaqui[0.95]{../evidencias/jflap/fichas/lista2-106.png}{lista2-106 en JFLAP: Multiple Run con las tres cadenas válidas y las
tres inválidas de la tabla anterior.}

\hypertarget{lista-2-ejercicio-107}{%
\subsubsection{Lista 2 · Ejercicio 107}\label{lista-2-ejercicio-107}}

\textbf{Enunciado.} Las palabras de Σ1 que incluyen al menos una vez 0101 y al menos una vez 021.

\textbf{Diseño.} Producto de los autómatas que buscan 0101 y 021: cada estado es un par (avance en 0101, avance en 021).

\textbf{Archivo.} \texttt{automatas/lista2/lista2-107.jff} · Σ = \{0, 1, 2\} · \textbar Q\textbar{} = 13

\begin{longtable}[]{@{}
  >{\raggedright\arraybackslash}p{(\columnwidth - 8\tabcolsep) * \real{0.1364}}
  >{\raggedright\arraybackslash}p{(\columnwidth - 8\tabcolsep) * \real{0.1364}}
  >{\raggedright\arraybackslash}p{(\columnwidth - 8\tabcolsep) * \real{0.1364}}
  >{\raggedright\arraybackslash}p{(\columnwidth - 8\tabcolsep) * \real{0.1364}}
  >{\raggedright\arraybackslash}p{(\columnwidth - 8\tabcolsep) * \real{0.4545}}@{}}
\caption{Tabla de transiciones (→ inicial, * aceptación, --- sin transición).}\tabularnewline
\toprule\noalign{}
\begin{minipage}[b]{\linewidth}\raggedright
δ
\end{minipage} & \begin{minipage}[b]{\linewidth}\raggedright
0
\end{minipage} & \begin{minipage}[b]{\linewidth}\raggedright
1
\end{minipage} & \begin{minipage}[b]{\linewidth}\raggedright
2
\end{minipage} & \begin{minipage}[b]{\linewidth}\raggedright
Qué recuerda del prefijo leído
\end{minipage} \\
\midrule\noalign{}
\endfirsthead
\toprule\noalign{}
\begin{minipage}[b]{\linewidth}\raggedright
δ
\end{minipage} & \begin{minipage}[b]{\linewidth}\raggedright
0
\end{minipage} & \begin{minipage}[b]{\linewidth}\raggedright
1
\end{minipage} & \begin{minipage}[b]{\linewidth}\raggedright
2
\end{minipage} & \begin{minipage}[b]{\linewidth}\raggedright
Qué recuerda del prefijo leído
\end{minipage} \\
\midrule\noalign{}
\endhead
\bottomrule\noalign{}
\endlastfoot
→ (λ, λ) & (0, 0) & (λ, λ) & (λ, λ) & en el primer autómata λ, en el segundo λ \\
(0, 0) & (0, 0) & (01, λ) & (λ, 02) & en el primer autómata 0, en el segundo 0 \\
(01, λ) & (010, 0) & (λ, λ) & (λ, λ) & en el primer autómata 01, en el segundo λ \\
(λ, 02) & (0, 0) & (λ, 021) & (λ, λ) & en el primer autómata λ, en el segundo 02 \\
(010, 0) & (0, 0) & (0101, λ) & (λ, 02) & en el primer autómata 010, en el segundo 0 \\
(λ, 021) & (0, 021) & (λ, 021) & (λ, 021) & en el primer autómata λ, en el segundo 021 \\
(0101, λ) & (0101, 0) & (0101, λ) & (0101, λ) & en el primer autómata 0101, en el segundo λ \\
(0, 021) & (0, 021) & (01, 021) & (λ, 021) & en el primer autómata 0, en el segundo 021 \\
(0101, 0) & (0101, 0) & (0101, λ) & (0101, 02) & en el primer autómata 0101, en el segundo 0 \\
(01, 021) & (010, 021) & (λ, 021) & (λ, 021) & en el primer autómata 01, en el segundo 021 \\
(0101, 02) & (0101, 0) & (0101, 021) & (0101, λ) & en el primer autómata 0101, en el segundo 02 \\
(010, 021) & (0, 021) & (0101, 021) & (λ, 021) & en el primer autómata 010, en el segundo 021 \\
* (0101, 021) & (0101, 021) & (0101, 021) & (0101, 021) & en el primer autómata 0101, en el segundo 021 \\
\end{longtable}

\begin{longtable}[]{@{}
  >{\raggedright\arraybackslash}p{(\columnwidth - 8\tabcolsep) * \real{0.2000}}
  >{\centering\arraybackslash}p{(\columnwidth - 8\tabcolsep) * \real{0.2000}}
  >{\raggedright\arraybackslash}p{(\columnwidth - 8\tabcolsep) * \real{0.2000}}
  >{\raggedright\arraybackslash}p{(\columnwidth - 8\tabcolsep) * \real{0.2000}}
  >{\raggedright\arraybackslash}p{(\columnwidth - 8\tabcolsep) * \real{0.2000}}@{}}
\caption{Cadenas de prueba.}\tabularnewline
\toprule\noalign{}
\begin{minipage}[b]{\linewidth}\raggedright
Cadena
\end{minipage} & \begin{minipage}[b]{\linewidth}\centering
Longitud
\end{minipage} & \begin{minipage}[b]{\linewidth}\raggedright
Esperado
\end{minipage} & \begin{minipage}[b]{\linewidth}\raggedright
Obtenido
\end{minipage} & \begin{minipage}[b]{\linewidth}\raggedright
Tipo de caso
\end{minipage} \\
\midrule\noalign{}
\endfirsthead
\toprule\noalign{}
\begin{minipage}[b]{\linewidth}\raggedright
Cadena
\end{minipage} & \begin{minipage}[b]{\linewidth}\centering
Longitud
\end{minipage} & \begin{minipage}[b]{\linewidth}\raggedright
Esperado
\end{minipage} & \begin{minipage}[b]{\linewidth}\raggedright
Obtenido
\end{minipage} & \begin{minipage}[b]{\linewidth}\raggedright
Tipo de caso
\end{minipage} \\
\midrule\noalign{}
\endhead
\bottomrule\noalign{}
\endlastfoot
\texttt{0101021} & 7 & acepta & acepta & válida \\
\texttt{01012021} & 8 & acepta & acepta & válida \\
\texttt{0101002120} & 10 & acepta & acepta & válida \\
λ & 0 & rechaza & rechaza & inválida \\
\texttt{1} & 1 & rechaza & rechaza & inválida \\
\texttt{010} & 3 & rechaza & rechaza & inválida \\
\texttt{0210101} & 7 & acepta & acepta & frontera: aceptada más corta \\
\texttt{111100} & 6 & rechaza & rechaza & frontera: justo por debajo de la longitud mínima \\
\end{longtable}

\evidenciaaqui[0.95]{../evidencias/jflap/fichas/lista2-107.png}{lista2-107 en JFLAP: Multiple Run con las tres cadenas válidas y las
tres inválidas de la tabla anterior.}

\hypertarget{lista-2-ejercicio-109}{%
\subsubsection{Lista 2 · Ejercicio 109}\label{lista-2-ejercicio-109}}

\textbf{Enunciado.} Las palabras de Σ2 que incluyen abbcb y cbaabb, que sí se pueden traslapar.

\textbf{Diseño.} Producto de los dos buscadores de patrones: como pueden traslaparse, cada patrón se busca de forma independiente.

\textbf{Archivo.} \texttt{automatas/lista2/lista2-109.jff} · Σ = \{0, 1, a, b, c\} · \textbar Q\textbar{} = 22

\begin{longtable}[]{@{}
  >{\raggedright\arraybackslash}p{(\columnwidth - 12\tabcolsep) * \real{0.1071}}
  >{\raggedright\arraybackslash}p{(\columnwidth - 12\tabcolsep) * \real{0.1071}}
  >{\raggedright\arraybackslash}p{(\columnwidth - 12\tabcolsep) * \real{0.1071}}
  >{\raggedright\arraybackslash}p{(\columnwidth - 12\tabcolsep) * \real{0.1071}}
  >{\raggedright\arraybackslash}p{(\columnwidth - 12\tabcolsep) * \real{0.1071}}
  >{\raggedright\arraybackslash}p{(\columnwidth - 12\tabcolsep) * \real{0.1071}}
  >{\raggedright\arraybackslash}p{(\columnwidth - 12\tabcolsep) * \real{0.3571}}@{}}
\caption{Tabla de transiciones (→ inicial, * aceptación, --- sin transición).}\tabularnewline
\toprule\noalign{}
\begin{minipage}[b]{\linewidth}\raggedright
δ
\end{minipage} & \begin{minipage}[b]{\linewidth}\raggedright
0
\end{minipage} & \begin{minipage}[b]{\linewidth}\raggedright
1
\end{minipage} & \begin{minipage}[b]{\linewidth}\raggedright
a
\end{minipage} & \begin{minipage}[b]{\linewidth}\raggedright
b
\end{minipage} & \begin{minipage}[b]{\linewidth}\raggedright
c
\end{minipage} & \begin{minipage}[b]{\linewidth}\raggedright
Qué recuerda del prefijo leído
\end{minipage} \\
\midrule\noalign{}
\endfirsthead
\toprule\noalign{}
\begin{minipage}[b]{\linewidth}\raggedright
δ
\end{minipage} & \begin{minipage}[b]{\linewidth}\raggedright
0
\end{minipage} & \begin{minipage}[b]{\linewidth}\raggedright
1
\end{minipage} & \begin{minipage}[b]{\linewidth}\raggedright
a
\end{minipage} & \begin{minipage}[b]{\linewidth}\raggedright
b
\end{minipage} & \begin{minipage}[b]{\linewidth}\raggedright
c
\end{minipage} & \begin{minipage}[b]{\linewidth}\raggedright
Qué recuerda del prefijo leído
\end{minipage} \\
\midrule\noalign{}
\endhead
\bottomrule\noalign{}
\endlastfoot
→ (λ, λ) & (λ, λ) & (λ, λ) & (a, λ) & (λ, λ) & (λ, c) & en el primer autómata λ, en el segundo λ \\
(a, λ) & (λ, λ) & (λ, λ) & (a, λ) & (ab, λ) & (λ, c) & en el primer autómata a, en el segundo λ \\
(λ, c) & (λ, λ) & (λ, λ) & (a, λ) & (λ, cb) & (λ, c) & en el primer autómata λ, en el segundo c \\
(ab, λ) & (λ, λ) & (λ, λ) & (a, λ) & (abb, λ) & (λ, c) & en el primer autómata ab, en el segundo λ \\
(λ, cb) & (λ, λ) & (λ, λ) & (a, cba) & (λ, λ) & (λ, c) & en el primer autómata λ, en el segundo cb \\
(abb, λ) & (λ, λ) & (λ, λ) & (a, λ) & (λ, λ) & (abbc, c) & en el primer autómata abb, en el segundo λ \\
(a, cba) & (λ, λ) & (λ, λ) & (a, cbaa) & (ab, λ) & (λ, c) & en el primer autómata a, en el segundo cba \\
(abbc, c) & (λ, λ) & (λ, λ) & (a, λ) & (abbcb, cb) & (λ, c) & en el primer autómata abbc, en el segundo c \\
(a, cbaa) & (λ, λ) & (λ, λ) & (a, λ) & (ab, cbaab) & (λ, c) & en el primer autómata a, en el segundo cbaa \\
(abbcb, cb) & (abbcb, λ) & (abbcb, λ) & (abbcb, cba) & (abbcb, λ) & (abbcb, c) & en el primer autómata abbcb, en el segundo cb \\
(ab, cbaab) & (λ, λ) & (λ, λ) & (a, λ) & (abb, cbaabb) & (λ, c) & en el primer autómata ab, en el segundo cbaab \\
(abbcb, cba) & (abbcb, λ) & (abbcb, λ) & (abbcb, cbaa) & (abbcb, λ) & (abbcb, c) & en el primer autómata abbcb, en el segundo cba \\
(abbcb, λ) & (abbcb, λ) & (abbcb, λ) & (abbcb, λ) & (abbcb, λ) & (abbcb, c) & en el primer autómata abbcb, en el segundo λ \\
(abbcb, c) & (abbcb, λ) & (abbcb, λ) & (abbcb, λ) & (abbcb, cb) & (abbcb, c) & en el primer autómata abbcb, en el segundo c \\
(abb, cbaabb) & (λ, cbaabb) & (λ, cbaabb) & (a, cbaabb) & (λ, cbaabb) & (abbc, cbaabb) & en el primer autómata abb, en el segundo cbaabb \\
(abbcb, cbaa) & (abbcb, λ) & (abbcb, λ) & (abbcb, λ) & (abbcb, cbaab) & (abbcb, c) & en el primer autómata abbcb, en el segundo cbaa \\
(a, cbaabb) & (λ, cbaabb) & (λ, cbaabb) & (a, cbaabb) & (ab, cbaabb) & (λ, cbaabb) & en el primer autómata a, en el segundo cbaabb \\
(λ, cbaabb) & (λ, cbaabb) & (λ, cbaabb) & (a, cbaabb) & (λ, cbaabb) & (λ, cbaabb) & en el primer autómata λ, en el segundo cbaabb \\
(abbc, cbaabb) & (λ, cbaabb) & (λ, cbaabb) & (a, cbaabb) & (abbcb, cbaabb) & (λ, cbaabb) & en el primer autómata abbc, en el segundo cbaabb \\
(abbcb, cbaab) & (abbcb, λ) & (abbcb, λ) & (abbcb, λ) & (abbcb, cbaabb) & (abbcb, c) & en el primer autómata abbcb, en el segundo cbaab \\
(ab, cbaabb) & (λ, cbaabb) & (λ, cbaabb) & (a, cbaabb) & (abb, cbaabb) & (λ, cbaabb) & en el primer autómata ab, en el segundo cbaabb \\
* (abbcb, cbaabb) & (abbcb, cbaabb) & (abbcb, cbaabb) & (abbcb, cbaabb) & (abbcb, cbaabb) & (abbcb, cbaabb) & en el primer autómata abbcb, en el segundo cbaabb \\
\end{longtable}

\begin{longtable}[]{@{}
  >{\raggedright\arraybackslash}p{(\columnwidth - 8\tabcolsep) * \real{0.2000}}
  >{\centering\arraybackslash}p{(\columnwidth - 8\tabcolsep) * \real{0.2000}}
  >{\raggedright\arraybackslash}p{(\columnwidth - 8\tabcolsep) * \real{0.2000}}
  >{\raggedright\arraybackslash}p{(\columnwidth - 8\tabcolsep) * \real{0.2000}}
  >{\raggedright\arraybackslash}p{(\columnwidth - 8\tabcolsep) * \real{0.2000}}@{}}
\caption{Cadenas de prueba.}\tabularnewline
\toprule\noalign{}
\begin{minipage}[b]{\linewidth}\raggedright
Cadena
\end{minipage} & \begin{minipage}[b]{\linewidth}\centering
Longitud
\end{minipage} & \begin{minipage}[b]{\linewidth}\raggedright
Esperado
\end{minipage} & \begin{minipage}[b]{\linewidth}\raggedright
Obtenido
\end{minipage} & \begin{minipage}[b]{\linewidth}\raggedright
Tipo de caso
\end{minipage} \\
\midrule\noalign{}
\endfirsthead
\toprule\noalign{}
\begin{minipage}[b]{\linewidth}\raggedright
Cadena
\end{minipage} & \begin{minipage}[b]{\linewidth}\centering
Longitud
\end{minipage} & \begin{minipage}[b]{\linewidth}\raggedright
Esperado
\end{minipage} & \begin{minipage}[b]{\linewidth}\raggedright
Obtenido
\end{minipage} & \begin{minipage}[b]{\linewidth}\raggedright
Tipo de caso
\end{minipage} \\
\midrule\noalign{}
\endhead
\bottomrule\noalign{}
\endlastfoot
\texttt{cbaabbcb} & 8 & acepta & acepta & válida \\
\texttt{0cbaabbcb} & 9 & acepta & acepta & válida \\
\texttt{acacbaabbcb} & 11 & acepta & acepta & válida \\
λ & 0 & rechaza & rechaza & inválida \\
\texttt{1} & 1 & rechaza & rechaza & inválida \\
\texttt{b0c} & 3 & rechaza & rechaza & inválida \\
\texttt{cbbcccb} & 7 & rechaza & rechaza & frontera: justo por debajo de la longitud mínima \\
\end{longtable}

\evidenciaaqui[0.95]{../evidencias/jflap/fichas/lista2-109.png}{lista2-109 en JFLAP: Multiple Run con las tres cadenas válidas y las
tres inválidas de la tabla anterior.}

\hypertarget{lista-2-ejercicio-111}{%
\subsubsection{Lista 2 · Ejercicio 111}\label{lista-2-ejercicio-111}}

\textbf{Enunciado.} Las palabras de Σ1 que contienen exactamente tres veces la subcadena 01002.

\textbf{Diseño.} Cada estado combina el número de apariciones de 01002 (0 a 3, y muerto si pasa de 3) con el avance en el patrón; 01002 no puede traslaparse consigo mismo.

\textbf{Archivo.} \texttt{automatas/lista2/lista2-111.jff} · Σ = \{0, 1, 2\} · \textbar Q\textbar{} = 21

\begin{longtable}[]{@{}
  >{\raggedright\arraybackslash}p{(\columnwidth - 8\tabcolsep) * \real{0.1364}}
  >{\raggedright\arraybackslash}p{(\columnwidth - 8\tabcolsep) * \real{0.1364}}
  >{\raggedright\arraybackslash}p{(\columnwidth - 8\tabcolsep) * \real{0.1364}}
  >{\raggedright\arraybackslash}p{(\columnwidth - 8\tabcolsep) * \real{0.1364}}
  >{\raggedright\arraybackslash}p{(\columnwidth - 8\tabcolsep) * \real{0.4545}}@{}}
\caption{Tabla de transiciones (→ inicial, * aceptación, --- sin transición).}\tabularnewline
\toprule\noalign{}
\begin{minipage}[b]{\linewidth}\raggedright
δ
\end{minipage} & \begin{minipage}[b]{\linewidth}\raggedright
0
\end{minipage} & \begin{minipage}[b]{\linewidth}\raggedright
1
\end{minipage} & \begin{minipage}[b]{\linewidth}\raggedright
2
\end{minipage} & \begin{minipage}[b]{\linewidth}\raggedright
Qué recuerda del prefijo leído
\end{minipage} \\
\midrule\noalign{}
\endfirsthead
\toprule\noalign{}
\begin{minipage}[b]{\linewidth}\raggedright
δ
\end{minipage} & \begin{minipage}[b]{\linewidth}\raggedright
0
\end{minipage} & \begin{minipage}[b]{\linewidth}\raggedright
1
\end{minipage} & \begin{minipage}[b]{\linewidth}\raggedright
2
\end{minipage} & \begin{minipage}[b]{\linewidth}\raggedright
Qué recuerda del prefijo leído
\end{minipage} \\
\midrule\noalign{}
\endhead
\bottomrule\noalign{}
\endlastfoot
→ 0·λ & 0·0 & 0·λ & 0·λ & van 0 apariciones; el sufijo útil leído es λ \\
0·0 & 0·0 & 0·01 & 0·λ & van 0 apariciones; el sufijo útil leído es 0 \\
0·01 & 0·010 & 0·λ & 0·λ & van 0 apariciones; el sufijo útil leído es 01 \\
0·010 & 0·0100 & 0·01 & 0·λ & van 0 apariciones; el sufijo útil leído es 010 \\
0·0100 & 0·0 & 0·01 & 1·λ & van 0 apariciones; el sufijo útil leído es 0100 \\
1·λ & 1·0 & 1·λ & 1·λ & van 1 apariciones; el sufijo útil leído es λ \\
1·0 & 1·0 & 1·01 & 1·λ & van 1 apariciones; el sufijo útil leído es 0 \\
1·01 & 1·010 & 1·λ & 1·λ & van 1 apariciones; el sufijo útil leído es 01 \\
1·010 & 1·0100 & 1·01 & 1·λ & van 1 apariciones; el sufijo útil leído es 010 \\
1·0100 & 1·0 & 1·01 & 2·λ & van 1 apariciones; el sufijo útil leído es 0100 \\
2·λ & 2·0 & 2·λ & 2·λ & van 2 apariciones; el sufijo útil leído es λ \\
2·0 & 2·0 & 2·01 & 2·λ & van 2 apariciones; el sufijo útil leído es 0 \\
2·01 & 2·010 & 2·λ & 2·λ & van 2 apariciones; el sufijo útil leído es 01 \\
2·010 & 2·0100 & 2·01 & 2·λ & van 2 apariciones; el sufijo útil leído es 010 \\
2·0100 & 2·0 & 2·01 & 3·λ & van 2 apariciones; el sufijo útil leído es 0100 \\
* 3·λ & 3·0 & 3·λ & 3·λ & van 3 apariciones; el sufijo útil leído es λ \\
* 3·0 & 3·0 & 3·01 & 3·λ & van 3 apariciones; el sufijo útil leído es 0 \\
* 3·01 & 3·010 & 3·λ & 3·λ & van 3 apariciones; el sufijo útil leído es 01 \\
* 3·010 & 3·0100 & 3·01 & 3·λ & van 3 apariciones; el sufijo útil leído es 010 \\
* 3·0100 & 3·0 & 3·01 & muerto & van 3 apariciones; el sufijo útil leído es 0100 \\
muerto & muerto & muerto & muerto & ya hay más de 3 apariciones de 01002 \\
\end{longtable}

\begin{longtable}[]{@{}
  >{\raggedright\arraybackslash}p{(\columnwidth - 8\tabcolsep) * \real{0.2000}}
  >{\centering\arraybackslash}p{(\columnwidth - 8\tabcolsep) * \real{0.2000}}
  >{\raggedright\arraybackslash}p{(\columnwidth - 8\tabcolsep) * \real{0.2000}}
  >{\raggedright\arraybackslash}p{(\columnwidth - 8\tabcolsep) * \real{0.2000}}
  >{\raggedright\arraybackslash}p{(\columnwidth - 8\tabcolsep) * \real{0.2000}}@{}}
\caption{Cadenas de prueba.}\tabularnewline
\toprule\noalign{}
\begin{minipage}[b]{\linewidth}\raggedright
Cadena
\end{minipage} & \begin{minipage}[b]{\linewidth}\centering
Longitud
\end{minipage} & \begin{minipage}[b]{\linewidth}\raggedright
Esperado
\end{minipage} & \begin{minipage}[b]{\linewidth}\raggedright
Obtenido
\end{minipage} & \begin{minipage}[b]{\linewidth}\raggedright
Tipo de caso
\end{minipage} \\
\midrule\noalign{}
\endfirsthead
\toprule\noalign{}
\begin{minipage}[b]{\linewidth}\raggedright
Cadena
\end{minipage} & \begin{minipage}[b]{\linewidth}\centering
Longitud
\end{minipage} & \begin{minipage}[b]{\linewidth}\raggedright
Esperado
\end{minipage} & \begin{minipage}[b]{\linewidth}\raggedright
Obtenido
\end{minipage} & \begin{minipage}[b]{\linewidth}\raggedright
Tipo de caso
\end{minipage} \\
\midrule\noalign{}
\endhead
\bottomrule\noalign{}
\endlastfoot
\texttt{010020100201002} & 15 & acepta & acepta & válida \\
\texttt{0100201002001002} & 16 & acepta & acepta & válida \\
\texttt{010020010020101002} & 18 & acepta & acepta & válida \\
λ & 0 & rechaza & rechaza & inválida \\
\texttt{0} & 1 & rechaza & rechaza & inválida \\
\texttt{111} & 3 & rechaza & rechaza & inválida \\
\texttt{01122210022010} & 14 & rechaza & rechaza & frontera: justo por debajo de la longitud mínima \\
\end{longtable}

\evidenciaaqui[0.95]{../evidencias/jflap/fichas/lista2-111.png}{lista2-111 en JFLAP: Multiple Run con las tres cadenas válidas y las
tres inválidas de la tabla anterior.}

\hypertarget{lista-2-ejercicio-112}{%
\subsubsection{Lista 2 · Ejercicio 112}\label{lista-2-ejercicio-112}}

\textbf{Enunciado.} Las palabras de Σ1 que contienen dos veces 010 sin traslaparse (01010 no pertenece, 010010 sí).

\textbf{Diseño.} Se interpreta «al menos dos apariciones disjuntas». Tras completar 010 la búsqueda reinicia desde λ, de modo que la siguiente aparición no reutiliza el último 0.

\textbf{Archivo.} \texttt{automatas/lista2/lista2-112.jff} · Σ = \{0, 1, 2\} · \textbar Q\textbar{} = 7

\begin{longtable}[]{@{}
  >{\raggedright\arraybackslash}p{(\columnwidth - 8\tabcolsep) * \real{0.1364}}
  >{\raggedright\arraybackslash}p{(\columnwidth - 8\tabcolsep) * \real{0.1364}}
  >{\raggedright\arraybackslash}p{(\columnwidth - 8\tabcolsep) * \real{0.1364}}
  >{\raggedright\arraybackslash}p{(\columnwidth - 8\tabcolsep) * \real{0.1364}}
  >{\raggedright\arraybackslash}p{(\columnwidth - 8\tabcolsep) * \real{0.4545}}@{}}
\caption{Tabla de transiciones (→ inicial, * aceptación, --- sin transición).}\tabularnewline
\toprule\noalign{}
\begin{minipage}[b]{\linewidth}\raggedright
δ
\end{minipage} & \begin{minipage}[b]{\linewidth}\raggedright
0
\end{minipage} & \begin{minipage}[b]{\linewidth}\raggedright
1
\end{minipage} & \begin{minipage}[b]{\linewidth}\raggedright
2
\end{minipage} & \begin{minipage}[b]{\linewidth}\raggedright
Qué recuerda del prefijo leído
\end{minipage} \\
\midrule\noalign{}
\endfirsthead
\toprule\noalign{}
\begin{minipage}[b]{\linewidth}\raggedright
δ
\end{minipage} & \begin{minipage}[b]{\linewidth}\raggedright
0
\end{minipage} & \begin{minipage}[b]{\linewidth}\raggedright
1
\end{minipage} & \begin{minipage}[b]{\linewidth}\raggedright
2
\end{minipage} & \begin{minipage}[b]{\linewidth}\raggedright
Qué recuerda del prefijo leído
\end{minipage} \\
\midrule\noalign{}
\endhead
\bottomrule\noalign{}
\endlastfoot
→ 0·λ & 0·0 & 0·λ & 0·λ & van 0 apariciones; el sufijo útil leído es λ \\
0·0 & 0·0 & 0·01 & 0·λ & van 0 apariciones; el sufijo útil leído es 0 \\
0·01 & 1·λ & 0·λ & 0·λ & van 0 apariciones; el sufijo útil leído es 01 \\
1·λ & 1·0 & 1·λ & 1·λ & van 1 apariciones; el sufijo útil leído es λ \\
1·0 & 1·0 & 1·01 & 1·λ & van 1 apariciones; el sufijo útil leído es 0 \\
1·01 & 2+ & 1·λ & 1·λ & van 1 apariciones; el sufijo útil leído es 01 \\
* 2+ & 2+ & 2+ & 2+ & ya hay 2 apariciones de 010 \\
\end{longtable}

\begin{longtable}[]{@{}
  >{\raggedright\arraybackslash}p{(\columnwidth - 8\tabcolsep) * \real{0.2000}}
  >{\centering\arraybackslash}p{(\columnwidth - 8\tabcolsep) * \real{0.2000}}
  >{\raggedright\arraybackslash}p{(\columnwidth - 8\tabcolsep) * \real{0.2000}}
  >{\raggedright\arraybackslash}p{(\columnwidth - 8\tabcolsep) * \real{0.2000}}
  >{\raggedright\arraybackslash}p{(\columnwidth - 8\tabcolsep) * \real{0.2000}}@{}}
\caption{Cadenas de prueba.}\tabularnewline
\toprule\noalign{}
\begin{minipage}[b]{\linewidth}\raggedright
Cadena
\end{minipage} & \begin{minipage}[b]{\linewidth}\centering
Longitud
\end{minipage} & \begin{minipage}[b]{\linewidth}\raggedright
Esperado
\end{minipage} & \begin{minipage}[b]{\linewidth}\raggedright
Obtenido
\end{minipage} & \begin{minipage}[b]{\linewidth}\raggedright
Tipo de caso
\end{minipage} \\
\midrule\noalign{}
\endfirsthead
\toprule\noalign{}
\begin{minipage}[b]{\linewidth}\raggedright
Cadena
\end{minipage} & \begin{minipage}[b]{\linewidth}\centering
Longitud
\end{minipage} & \begin{minipage}[b]{\linewidth}\raggedright
Esperado
\end{minipage} & \begin{minipage}[b]{\linewidth}\raggedright
Obtenido
\end{minipage} & \begin{minipage}[b]{\linewidth}\raggedright
Tipo de caso
\end{minipage} \\
\midrule\noalign{}
\endhead
\bottomrule\noalign{}
\endlastfoot
\texttt{010010} & 6 & acepta & acepta & válida \\
\texttt{0100010} & 7 & acepta & acepta & válida \\
\texttt{001001012} & 9 & acepta & acepta & válida \\
λ & 0 & rechaza & rechaza & inválida \\
\texttt{2} & 1 & rechaza & rechaza & inválida \\
\texttt{122} & 3 & rechaza & rechaza & inválida \\
\texttt{20101} & 5 & rechaza & rechaza & frontera: justo por debajo de la longitud mínima \\
\end{longtable}

\evidenciaaqui[0.95]{../evidencias/jflap/fichas/lista2-112.png}{lista2-112 en JFLAP: Multiple Run con las tres cadenas válidas y las
tres inválidas de la tabla anterior.}

\hypertarget{lista-2-ejercicio-113}{%
\subsubsection{Lista 2 · Ejercicio 113}\label{lista-2-ejercicio-113}}

\textbf{Enunciado.} Las palabras de Σ1 que contienen dos veces 010, que sí se pueden traslapar.

\textbf{Diseño.} Igual que 112, pero tras completar 010 el autómata conserva el borde 0, que puede iniciar la siguiente aparición (01010).

\textbf{Archivo.} \texttt{automatas/lista2/lista2-113.jff} · Σ = \{0, 1, 2\} · \textbar Q\textbar{} = 7

\begin{longtable}[]{@{}
  >{\raggedright\arraybackslash}p{(\columnwidth - 8\tabcolsep) * \real{0.1364}}
  >{\raggedright\arraybackslash}p{(\columnwidth - 8\tabcolsep) * \real{0.1364}}
  >{\raggedright\arraybackslash}p{(\columnwidth - 8\tabcolsep) * \real{0.1364}}
  >{\raggedright\arraybackslash}p{(\columnwidth - 8\tabcolsep) * \real{0.1364}}
  >{\raggedright\arraybackslash}p{(\columnwidth - 8\tabcolsep) * \real{0.4545}}@{}}
\caption{Tabla de transiciones (→ inicial, * aceptación, --- sin transición).}\tabularnewline
\toprule\noalign{}
\begin{minipage}[b]{\linewidth}\raggedright
δ
\end{minipage} & \begin{minipage}[b]{\linewidth}\raggedright
0
\end{minipage} & \begin{minipage}[b]{\linewidth}\raggedright
1
\end{minipage} & \begin{minipage}[b]{\linewidth}\raggedright
2
\end{minipage} & \begin{minipage}[b]{\linewidth}\raggedright
Qué recuerda del prefijo leído
\end{minipage} \\
\midrule\noalign{}
\endfirsthead
\toprule\noalign{}
\begin{minipage}[b]{\linewidth}\raggedright
δ
\end{minipage} & \begin{minipage}[b]{\linewidth}\raggedright
0
\end{minipage} & \begin{minipage}[b]{\linewidth}\raggedright
1
\end{minipage} & \begin{minipage}[b]{\linewidth}\raggedright
2
\end{minipage} & \begin{minipage}[b]{\linewidth}\raggedright
Qué recuerda del prefijo leído
\end{minipage} \\
\midrule\noalign{}
\endhead
\bottomrule\noalign{}
\endlastfoot
→ 0·λ & 0·0 & 0·λ & 0·λ & van 0 apariciones; el sufijo útil leído es λ \\
0·0 & 0·0 & 0·01 & 0·λ & van 0 apariciones; el sufijo útil leído es 0 \\
0·01 & 1·0 & 0·λ & 0·λ & van 0 apariciones; el sufijo útil leído es 01 \\
1·0 & 1·0 & 1·01 & 1·λ & van 1 apariciones; el sufijo útil leído es 0 \\
1·01 & 2+ & 1·λ & 1·λ & van 1 apariciones; el sufijo útil leído es 01 \\
1·λ & 1·0 & 1·λ & 1·λ & van 1 apariciones; el sufijo útil leído es λ \\
* 2+ & 2+ & 2+ & 2+ & ya hay 2 apariciones de 010 \\
\end{longtable}

\begin{longtable}[]{@{}
  >{\raggedright\arraybackslash}p{(\columnwidth - 8\tabcolsep) * \real{0.2000}}
  >{\centering\arraybackslash}p{(\columnwidth - 8\tabcolsep) * \real{0.2000}}
  >{\raggedright\arraybackslash}p{(\columnwidth - 8\tabcolsep) * \real{0.2000}}
  >{\raggedright\arraybackslash}p{(\columnwidth - 8\tabcolsep) * \real{0.2000}}
  >{\raggedright\arraybackslash}p{(\columnwidth - 8\tabcolsep) * \real{0.2000}}@{}}
\caption{Cadenas de prueba.}\tabularnewline
\toprule\noalign{}
\begin{minipage}[b]{\linewidth}\raggedright
Cadena
\end{minipage} & \begin{minipage}[b]{\linewidth}\centering
Longitud
\end{minipage} & \begin{minipage}[b]{\linewidth}\raggedright
Esperado
\end{minipage} & \begin{minipage}[b]{\linewidth}\raggedright
Obtenido
\end{minipage} & \begin{minipage}[b]{\linewidth}\raggedright
Tipo de caso
\end{minipage} \\
\midrule\noalign{}
\endfirsthead
\toprule\noalign{}
\begin{minipage}[b]{\linewidth}\raggedright
Cadena
\end{minipage} & \begin{minipage}[b]{\linewidth}\centering
Longitud
\end{minipage} & \begin{minipage}[b]{\linewidth}\raggedright
Esperado
\end{minipage} & \begin{minipage}[b]{\linewidth}\raggedright
Obtenido
\end{minipage} & \begin{minipage}[b]{\linewidth}\raggedright
Tipo de caso
\end{minipage} \\
\midrule\noalign{}
\endhead
\bottomrule\noalign{}
\endlastfoot
\texttt{01010} & 5 & acepta & acepta & válida \\
\texttt{010010} & 6 & acepta & acepta & válida \\
\texttt{01010101} & 8 & acepta & acepta & válida \\
λ & 0 & rechaza & rechaza & inválida \\
\texttt{1} & 1 & rechaza & rechaza & inválida \\
\texttt{012} & 3 & rechaza & rechaza & inválida \\
\texttt{0100} & 4 & rechaza & rechaza & frontera: justo por debajo de la longitud mínima \\
\end{longtable}

\evidenciaaqui[0.95]{../evidencias/jflap/fichas/lista2-113.png}{lista2-113 en JFLAP: Multiple Run con las tres cadenas válidas y las
tres inválidas de la tabla anterior.}

\hypertarget{lista-2-ejercicio-116}{%
\subsubsection{Lista 2 · Ejercicio 116}\label{lista-2-ejercicio-116}}

\textbf{Enunciado.} L = \{w ∈ Σ2* : w = aⁿbᵐcⁿ⁺¹, m, n ∈ N\}.

El lenguaje no es regular (lema de bombeo); no existe AFD.

\hypertarget{lista-2-ejercicio-118}{%
\subsubsection{Lista 2 · Ejercicio 118}\label{lista-2-ejercicio-118}}

\textbf{Enunciado.} L = \{w ∈ Σ3* : w = axabya, x, y ∈ Σ3*\}.

\textbf{Diseño.} La primera a no puede ser la del ab interior (debe haber una a inicial propia), y la a final debe ir después de la b.

\textbf{Archivo.} \texttt{automatas/lista2/lista2-118.jff} · Σ = \{a, b\} · \textbar Q\textbar{} = 6

\begin{longtable}[]{@{}
  >{\raggedright\arraybackslash}p{(\columnwidth - 6\tabcolsep) * \real{0.1579}}
  >{\raggedright\arraybackslash}p{(\columnwidth - 6\tabcolsep) * \real{0.1579}}
  >{\raggedright\arraybackslash}p{(\columnwidth - 6\tabcolsep) * \real{0.1579}}
  >{\raggedright\arraybackslash}p{(\columnwidth - 6\tabcolsep) * \real{0.5263}}@{}}
\caption{Tabla de transiciones (→ inicial, * aceptación, --- sin transición).}\tabularnewline
\toprule\noalign{}
\begin{minipage}[b]{\linewidth}\raggedright
δ
\end{minipage} & \begin{minipage}[b]{\linewidth}\raggedright
a
\end{minipage} & \begin{minipage}[b]{\linewidth}\raggedright
b
\end{minipage} & \begin{minipage}[b]{\linewidth}\raggedright
Qué recuerda del prefijo leído
\end{minipage} \\
\midrule\noalign{}
\endfirsthead
\toprule\noalign{}
\begin{minipage}[b]{\linewidth}\raggedright
δ
\end{minipage} & \begin{minipage}[b]{\linewidth}\raggedright
a
\end{minipage} & \begin{minipage}[b]{\linewidth}\raggedright
b
\end{minipage} & \begin{minipage}[b]{\linewidth}\raggedright
Qué recuerda del prefijo leído
\end{minipage} \\
\midrule\noalign{}
\endhead
\bottomrule\noalign{}
\endlastfoot
→ ini & a & muerto & no se ha leído nada \\
a & a·a & a & empieza con a; aún no hay una a utilizable para el ab interior \\
muerto & muerto & muerto & la cadena no empieza con a \\
a·a & a·a & ab & el último símbolo es una a que puede iniciar el ab interior \\
ab & ab·a & ab & ya apareció el ab interior; el último símbolo no es a \\
* ab·a & ab·a & ab & ya apareció el ab interior y el último símbolo es a \\
\end{longtable}

\begin{longtable}[]{@{}
  >{\raggedright\arraybackslash}p{(\columnwidth - 8\tabcolsep) * \real{0.2000}}
  >{\centering\arraybackslash}p{(\columnwidth - 8\tabcolsep) * \real{0.2000}}
  >{\raggedright\arraybackslash}p{(\columnwidth - 8\tabcolsep) * \real{0.2000}}
  >{\raggedright\arraybackslash}p{(\columnwidth - 8\tabcolsep) * \real{0.2000}}
  >{\raggedright\arraybackslash}p{(\columnwidth - 8\tabcolsep) * \real{0.2000}}@{}}
\caption{Cadenas de prueba.}\tabularnewline
\toprule\noalign{}
\begin{minipage}[b]{\linewidth}\raggedright
Cadena
\end{minipage} & \begin{minipage}[b]{\linewidth}\centering
Longitud
\end{minipage} & \begin{minipage}[b]{\linewidth}\raggedright
Esperado
\end{minipage} & \begin{minipage}[b]{\linewidth}\raggedright
Obtenido
\end{minipage} & \begin{minipage}[b]{\linewidth}\raggedright
Tipo de caso
\end{minipage} \\
\midrule\noalign{}
\endfirsthead
\toprule\noalign{}
\begin{minipage}[b]{\linewidth}\raggedright
Cadena
\end{minipage} & \begin{minipage}[b]{\linewidth}\centering
Longitud
\end{minipage} & \begin{minipage}[b]{\linewidth}\raggedright
Esperado
\end{minipage} & \begin{minipage}[b]{\linewidth}\raggedright
Obtenido
\end{minipage} & \begin{minipage}[b]{\linewidth}\raggedright
Tipo de caso
\end{minipage} \\
\midrule\noalign{}
\endhead
\bottomrule\noalign{}
\endlastfoot
\texttt{aaba} & 4 & acepta & acepta & válida \\
\texttt{aabba} & 5 & acepta & acepta & válida \\
\texttt{aabbaaa} & 7 & acepta & acepta & válida \\
λ & 0 & rechaza & rechaza & inválida \\
\texttt{a} & 1 & rechaza & rechaza & inválida \\
\texttt{bbb} & 3 & rechaza & rechaza & inválida \\
\texttt{bab} & 3 & rechaza & rechaza & frontera: justo por debajo de la longitud mínima \\
\end{longtable}

\evidenciaaqui[0.95]{../evidencias/jflap/fichas/lista2-118.png}{lista2-118 en JFLAP: Multiple Run con las tres cadenas válidas y las
tres inválidas de la tabla anterior.}

\hypertarget{lista-2-ejercicio-120}{%
\subsubsection{Lista 2 · Ejercicio 120}\label{lista-2-ejercicio-120}}

\textbf{Enunciado.} Las palabras de Σ1 tales que el primer carácter y el último son distintos.

\textbf{Diseño.} Se recuerda el primer símbolo y si el último coincide con él.

\textbf{Archivo.} \texttt{automatas/lista2/lista2-120.jff} · Σ = \{0, 1, 2\} · \textbar Q\textbar{} = 7

\begin{longtable}[]{@{}
  >{\raggedright\arraybackslash}p{(\columnwidth - 8\tabcolsep) * \real{0.1364}}
  >{\raggedright\arraybackslash}p{(\columnwidth - 8\tabcolsep) * \real{0.1364}}
  >{\raggedright\arraybackslash}p{(\columnwidth - 8\tabcolsep) * \real{0.1364}}
  >{\raggedright\arraybackslash}p{(\columnwidth - 8\tabcolsep) * \real{0.1364}}
  >{\raggedright\arraybackslash}p{(\columnwidth - 8\tabcolsep) * \real{0.4545}}@{}}
\caption{Tabla de transiciones (→ inicial, * aceptación, --- sin transición).}\tabularnewline
\toprule\noalign{}
\begin{minipage}[b]{\linewidth}\raggedright
δ
\end{minipage} & \begin{minipage}[b]{\linewidth}\raggedright
0
\end{minipage} & \begin{minipage}[b]{\linewidth}\raggedright
1
\end{minipage} & \begin{minipage}[b]{\linewidth}\raggedright
2
\end{minipage} & \begin{minipage}[b]{\linewidth}\raggedright
Qué recuerda del prefijo leído
\end{minipage} \\
\midrule\noalign{}
\endfirsthead
\toprule\noalign{}
\begin{minipage}[b]{\linewidth}\raggedright
δ
\end{minipage} & \begin{minipage}[b]{\linewidth}\raggedright
0
\end{minipage} & \begin{minipage}[b]{\linewidth}\raggedright
1
\end{minipage} & \begin{minipage}[b]{\linewidth}\raggedright
2
\end{minipage} & \begin{minipage}[b]{\linewidth}\raggedright
Qué recuerda del prefijo leído
\end{minipage} \\
\midrule\noalign{}
\endhead
\bottomrule\noalign{}
\endlastfoot
→ ini & 0\ldots= & 1\ldots= & 2\ldots= & no se ha leído nada \\
0\ldots= & 0\ldots= & 0\ldots≠ & 0\ldots≠ & empieza con 0 y el último símbolo es igual \\
1\ldots= & 1\ldots≠ & 1\ldots= & 1\ldots≠ & empieza con 1 y el último símbolo es igual \\
2\ldots= & 2\ldots≠ & 2\ldots≠ & 2\ldots= & empieza con 2 y el último símbolo es igual \\
* 0\ldots≠ & 0\ldots= & 0\ldots≠ & 0\ldots≠ & empieza con 0 y el último símbolo es distinto \\
* 1\ldots≠ & 1\ldots≠ & 1\ldots= & 1\ldots≠ & empieza con 1 y el último símbolo es distinto \\
* 2\ldots≠ & 2\ldots≠ & 2\ldots≠ & 2\ldots= & empieza con 2 y el último símbolo es distinto \\
\end{longtable}

\begin{longtable}[]{@{}
  >{\raggedright\arraybackslash}p{(\columnwidth - 8\tabcolsep) * \real{0.2000}}
  >{\centering\arraybackslash}p{(\columnwidth - 8\tabcolsep) * \real{0.2000}}
  >{\raggedright\arraybackslash}p{(\columnwidth - 8\tabcolsep) * \real{0.2000}}
  >{\raggedright\arraybackslash}p{(\columnwidth - 8\tabcolsep) * \real{0.2000}}
  >{\raggedright\arraybackslash}p{(\columnwidth - 8\tabcolsep) * \real{0.2000}}@{}}
\caption{Cadenas de prueba.}\tabularnewline
\toprule\noalign{}
\begin{minipage}[b]{\linewidth}\raggedright
Cadena
\end{minipage} & \begin{minipage}[b]{\linewidth}\centering
Longitud
\end{minipage} & \begin{minipage}[b]{\linewidth}\raggedright
Esperado
\end{minipage} & \begin{minipage}[b]{\linewidth}\raggedright
Obtenido
\end{minipage} & \begin{minipage}[b]{\linewidth}\raggedright
Tipo de caso
\end{minipage} \\
\midrule\noalign{}
\endfirsthead
\toprule\noalign{}
\begin{minipage}[b]{\linewidth}\raggedright
Cadena
\end{minipage} & \begin{minipage}[b]{\linewidth}\centering
Longitud
\end{minipage} & \begin{minipage}[b]{\linewidth}\raggedright
Esperado
\end{minipage} & \begin{minipage}[b]{\linewidth}\raggedright
Obtenido
\end{minipage} & \begin{minipage}[b]{\linewidth}\raggedright
Tipo de caso
\end{minipage} \\
\midrule\noalign{}
\endhead
\bottomrule\noalign{}
\endlastfoot
\texttt{21} & 2 & acepta & acepta & válida \\
\texttt{012} & 3 & acepta & acepta & válida \\
\texttt{20001} & 5 & acepta & acepta & válida \\
λ & 0 & rechaza & rechaza & inválida \\
\texttt{0} & 1 & rechaza & rechaza & inválida \\
\texttt{222} & 3 & rechaza & rechaza & inválida \\
\texttt{01} & 2 & acepta & acepta & frontera: aceptada más corta \\
\texttt{1} & 1 & rechaza & rechaza & frontera: justo por debajo de la longitud mínima \\
\end{longtable}

\evidenciaaqui[0.95]{../evidencias/jflap/fichas/lista2-120.png}{lista2-120 en JFLAP: Multiple Run con las tres cadenas válidas y las
tres inválidas de la tabla anterior.}

\hypertarget{lista-2-ejercicio-121}{%
\subsubsection{Lista 2 · Ejercicio 121}\label{lista-2-ejercicio-121}}

\textbf{Enunciado.} Las palabras de Σ1 de longitud 3n tales que cada bloque de longitud 3 tiene al menos un 0.

\textbf{Diseño.} Se lleva la posición dentro del bloque y si el bloque ya contiene un 0. λ (n = 0) se acepta.

\textbf{Archivo.} \texttt{automatas/lista2/lista2-121.jff} · Σ = \{0, 1, 2\} · \textbar Q\textbar{} = 6

\begin{longtable}[]{@{}
  >{\raggedright\arraybackslash}p{(\columnwidth - 8\tabcolsep) * \real{0.1364}}
  >{\raggedright\arraybackslash}p{(\columnwidth - 8\tabcolsep) * \real{0.1364}}
  >{\raggedright\arraybackslash}p{(\columnwidth - 8\tabcolsep) * \real{0.1364}}
  >{\raggedright\arraybackslash}p{(\columnwidth - 8\tabcolsep) * \real{0.1364}}
  >{\raggedright\arraybackslash}p{(\columnwidth - 8\tabcolsep) * \real{0.4545}}@{}}
\caption{Tabla de transiciones (→ inicial, * aceptación, --- sin transición).}\tabularnewline
\toprule\noalign{}
\begin{minipage}[b]{\linewidth}\raggedright
δ
\end{minipage} & \begin{minipage}[b]{\linewidth}\raggedright
0
\end{minipage} & \begin{minipage}[b]{\linewidth}\raggedright
1
\end{minipage} & \begin{minipage}[b]{\linewidth}\raggedright
2
\end{minipage} & \begin{minipage}[b]{\linewidth}\raggedright
Qué recuerda del prefijo leído
\end{minipage} \\
\midrule\noalign{}
\endfirsthead
\toprule\noalign{}
\begin{minipage}[b]{\linewidth}\raggedright
δ
\end{minipage} & \begin{minipage}[b]{\linewidth}\raggedright
0
\end{minipage} & \begin{minipage}[b]{\linewidth}\raggedright
1
\end{minipage} & \begin{minipage}[b]{\linewidth}\raggedright
2
\end{minipage} & \begin{minipage}[b]{\linewidth}\raggedright
Qué recuerda del prefijo leído
\end{minipage} \\
\midrule\noalign{}
\endhead
\bottomrule\noalign{}
\endlastfoot
→* b0 & b1+0 & b1 & b1 & posición 0 dentro del bloque; aún no hay un 0 en el bloque \\
b1+0 & b2+0 & b2+0 & b2+0 & posición 1 dentro del bloque; ya hay un 0 en el bloque \\
b1 & b2+0 & b2 & b2 & posición 1 dentro del bloque; aún no hay un 0 en el bloque \\
b2+0 & b0 & b0 & b0 & posición 2 dentro del bloque; ya hay un 0 en el bloque \\
b2 & b0 & muerto & muerto & posición 2 dentro del bloque; aún no hay un 0 en el bloque \\
muerto & muerto & muerto & muerto & un bloque terminó sin 0 \\
\end{longtable}

\begin{longtable}[]{@{}lclll@{}}
\caption{Cadenas de prueba.}\tabularnewline
\toprule\noalign{}
Cadena & Longitud & Esperado & Obtenido & Tipo de caso \\
\midrule\noalign{}
\endfirsthead
\toprule\noalign{}
Cadena & Longitud & Esperado & Obtenido & Tipo de caso \\
\midrule\noalign{}
\endhead
\bottomrule\noalign{}
\endlastfoot
λ & 0 & acepta & acepta & válida \\
\texttt{011} & 3 & acepta & acepta & válida \\
\texttt{012220} & 6 & acepta & acepta & válida \\
\texttt{2} & 1 & rechaza & rechaza & inválida \\
\texttt{02} & 2 & rechaza & rechaza & inválida \\
\texttt{0202} & 4 & rechaza & rechaza & inválida \\
\end{longtable}

\evidenciaaqui[0.95]{../evidencias/jflap/fichas/lista2-121.png}{lista2-121 en JFLAP: Multiple Run con las tres cadenas válidas y las
tres inválidas de la tabla anterior.}

\hypertarget{lista-2-ejercicio-122}{%
\subsubsection{Lista 2 · Ejercicio 122}\label{lista-2-ejercicio-122}}

\textbf{Enunciado.} Las palabras de Σ1 tales que cualquier cadena de tres símbolos consecutivos tiene un 0.

\textbf{Diseño.} Basta contar cuántos símbolos distintos de 0 hay al final; tres seguidos invalidan la cadena.

\textbf{Archivo.} \texttt{automatas/lista2/lista2-122.jff} · Σ = \{0, 1, 2\} · \textbar Q\textbar{} = 4

\begin{longtable}[]{@{}
  >{\raggedright\arraybackslash}p{(\columnwidth - 8\tabcolsep) * \real{0.1364}}
  >{\raggedright\arraybackslash}p{(\columnwidth - 8\tabcolsep) * \real{0.1364}}
  >{\raggedright\arraybackslash}p{(\columnwidth - 8\tabcolsep) * \real{0.1364}}
  >{\raggedright\arraybackslash}p{(\columnwidth - 8\tabcolsep) * \real{0.1364}}
  >{\raggedright\arraybackslash}p{(\columnwidth - 8\tabcolsep) * \real{0.4545}}@{}}
\caption{Tabla de transiciones (→ inicial, * aceptación, --- sin transición).}\tabularnewline
\toprule\noalign{}
\begin{minipage}[b]{\linewidth}\raggedright
δ
\end{minipage} & \begin{minipage}[b]{\linewidth}\raggedright
0
\end{minipage} & \begin{minipage}[b]{\linewidth}\raggedright
1
\end{minipage} & \begin{minipage}[b]{\linewidth}\raggedright
2
\end{minipage} & \begin{minipage}[b]{\linewidth}\raggedright
Qué recuerda del prefijo leído
\end{minipage} \\
\midrule\noalign{}
\endfirsthead
\toprule\noalign{}
\begin{minipage}[b]{\linewidth}\raggedright
δ
\end{minipage} & \begin{minipage}[b]{\linewidth}\raggedright
0
\end{minipage} & \begin{minipage}[b]{\linewidth}\raggedright
1
\end{minipage} & \begin{minipage}[b]{\linewidth}\raggedright
2
\end{minipage} & \begin{minipage}[b]{\linewidth}\raggedright
Qué recuerda del prefijo leído
\end{minipage} \\
\midrule\noalign{}
\endhead
\bottomrule\noalign{}
\endlastfoot
→* 0 sin 0 & 0 sin 0 & 1 sin 0 & 1 sin 0 & los últimos 0 símbolos son distintos de 0 \\
* 1 sin 0 & 0 sin 0 & 2 sin 0 & 2 sin 0 & los últimos 1 símbolos son distintos de 0 \\
* 2 sin 0 & 0 sin 0 & muerto & muerto & los últimos 2 símbolos son distintos de 0 \\
muerto & muerto & muerto & muerto & hubo tres símbolos seguidos sin 0 \\
\end{longtable}

\begin{longtable}[]{@{}lclll@{}}
\caption{Cadenas de prueba.}\tabularnewline
\toprule\noalign{}
Cadena & Longitud & Esperado & Obtenido & Tipo de caso \\
\midrule\noalign{}
\endfirsthead
\toprule\noalign{}
Cadena & Longitud & Esperado & Obtenido & Tipo de caso \\
\midrule\noalign{}
\endhead
\bottomrule\noalign{}
\endlastfoot
λ & 0 & acepta & acepta & válida \\
\texttt{1} & 1 & acepta & acepta & válida \\
\texttt{100} & 3 & acepta & acepta & válida \\
\texttt{112} & 3 & rechaza & rechaza & inválida \\
\texttt{0122} & 4 & rechaza & rechaza & inválida \\
\texttt{122010} & 6 & rechaza & rechaza & inválida \\
\end{longtable}

\evidenciaaqui[0.95]{../evidencias/jflap/fichas/lista2-122.png}{lista2-122 en JFLAP: Multiple Run con las tres cadenas válidas y las
tres inválidas de la tabla anterior.}

\hypertarget{lista-2-ejercicio-123}{%
\subsubsection{Lista 2 · Ejercicio 123}\label{lista-2-ejercicio-123}}

\textbf{Enunciado.} Las palabras de Σ3 de la forma wxw donde w ∈ Σ3² y x ∈ Σ3*.

\textbf{Diseño.} Se memoriza w (los dos primeros símbolos) y, para cada w, tres situaciones: la cadena termina en w, el último símbolo es la primera letra de w, u otra. Tras los dos primeros símbolos el estado es «otra» aunque la cadena termine en w, porque las dos copias de w no pueden traslaparse (se exige longitud ≥ 4).

\textbf{Archivo.} \texttt{automatas/lista2/lista2-123.jff} · Σ = \{a, b\} · \textbar Q\textbar{} = 15

\begin{longtable}[]{@{}
  >{\raggedright\arraybackslash}p{(\columnwidth - 6\tabcolsep) * \real{0.1579}}
  >{\raggedright\arraybackslash}p{(\columnwidth - 6\tabcolsep) * \real{0.1579}}
  >{\raggedright\arraybackslash}p{(\columnwidth - 6\tabcolsep) * \real{0.1579}}
  >{\raggedright\arraybackslash}p{(\columnwidth - 6\tabcolsep) * \real{0.5263}}@{}}
\caption{Tabla de transiciones (→ inicial, * aceptación, --- sin transición).}\tabularnewline
\toprule\noalign{}
\begin{minipage}[b]{\linewidth}\raggedright
δ
\end{minipage} & \begin{minipage}[b]{\linewidth}\raggedright
a
\end{minipage} & \begin{minipage}[b]{\linewidth}\raggedright
b
\end{minipage} & \begin{minipage}[b]{\linewidth}\raggedright
Qué recuerda del prefijo leído
\end{minipage} \\
\midrule\noalign{}
\endfirsthead
\toprule\noalign{}
\begin{minipage}[b]{\linewidth}\raggedright
δ
\end{minipage} & \begin{minipage}[b]{\linewidth}\raggedright
a
\end{minipage} & \begin{minipage}[b]{\linewidth}\raggedright
b
\end{minipage} & \begin{minipage}[b]{\linewidth}\raggedright
Qué recuerda del prefijo leído
\end{minipage} \\
\midrule\noalign{}
\endhead
\bottomrule\noalign{}
\endlastfoot
→ ini & a & b & no se ha leído nada \\
a & aa:∅ & ab:∅ & el primer símbolo es a \\
b & ba:∅ & bb:∅ & el primer símbolo es b \\
aa:∅ & aa:a\ldots{} & aa:∅ & w = aa; el final leído no sirve para la copia final de w (o la longitud aún es 2) \\
ab:∅ & ab:a\ldots{} & ab:∅ & w = ab; el final leído no sirve para la copia final de w (o la longitud aún es 2) \\
ba:∅ & ba:∅ & ba:b\ldots{} & w = ba; el final leído no sirve para la copia final de w (o la longitud aún es 2) \\
bb:∅ & bb:∅ & bb:b\ldots{} & w = bb; el final leído no sirve para la copia final de w (o la longitud aún es 2) \\
aa:a\ldots{} & aa:aa & aa:∅ & w = aa; el último símbolo es a, que puede iniciar la copia final de w \\
ab:a\ldots{} & ab:a\ldots{} & ab:ab & w = ab; el último símbolo es a, que puede iniciar la copia final de w \\
ba:b\ldots{} & ba:ba & ba:b\ldots{} & w = ba; el último símbolo es b, que puede iniciar la copia final de w \\
bb:b\ldots{} & bb:∅ & bb:bb & w = bb; el último símbolo es b, que puede iniciar la copia final de w \\
* aa:aa & aa:aa & aa:∅ & w = aa y la cadena (de longitud ≥ 4) termina en aa \\
* ab:ab & ab:a\ldots{} & ab:∅ & w = ab y la cadena (de longitud ≥ 4) termina en ab \\
* ba:ba & ba:∅ & ba:b\ldots{} & w = ba y la cadena (de longitud ≥ 4) termina en ba \\
* bb:bb & bb:∅ & bb:bb & w = bb y la cadena (de longitud ≥ 4) termina en bb \\
\end{longtable}

\begin{longtable}[]{@{}
  >{\raggedright\arraybackslash}p{(\columnwidth - 8\tabcolsep) * \real{0.2000}}
  >{\centering\arraybackslash}p{(\columnwidth - 8\tabcolsep) * \real{0.2000}}
  >{\raggedright\arraybackslash}p{(\columnwidth - 8\tabcolsep) * \real{0.2000}}
  >{\raggedright\arraybackslash}p{(\columnwidth - 8\tabcolsep) * \real{0.2000}}
  >{\raggedright\arraybackslash}p{(\columnwidth - 8\tabcolsep) * \real{0.2000}}@{}}
\caption{Cadenas de prueba.}\tabularnewline
\toprule\noalign{}
\begin{minipage}[b]{\linewidth}\raggedright
Cadena
\end{minipage} & \begin{minipage}[b]{\linewidth}\centering
Longitud
\end{minipage} & \begin{minipage}[b]{\linewidth}\raggedright
Esperado
\end{minipage} & \begin{minipage}[b]{\linewidth}\raggedright
Obtenido
\end{minipage} & \begin{minipage}[b]{\linewidth}\raggedright
Tipo de caso
\end{minipage} \\
\midrule\noalign{}
\endfirsthead
\toprule\noalign{}
\begin{minipage}[b]{\linewidth}\raggedright
Cadena
\end{minipage} & \begin{minipage}[b]{\linewidth}\centering
Longitud
\end{minipage} & \begin{minipage}[b]{\linewidth}\raggedright
Esperado
\end{minipage} & \begin{minipage}[b]{\linewidth}\raggedright
Obtenido
\end{minipage} & \begin{minipage}[b]{\linewidth}\raggedright
Tipo de caso
\end{minipage} \\
\midrule\noalign{}
\endhead
\bottomrule\noalign{}
\endlastfoot
\texttt{bbbb} & 4 & acepta & acepta & válida \\
\texttt{abbab} & 5 & acepta & acepta & válida \\
\texttt{bbbbabb} & 7 & acepta & acepta & válida \\
λ & 0 & rechaza & rechaza & inválida \\
\texttt{a} & 1 & rechaza & rechaza & inválida \\
\texttt{bab} & 3 & rechaza & rechaza & inválida \\
\texttt{abab} & 4 & acepta & acepta & frontera: aceptada más corta \\
\texttt{aaa} & 3 & rechaza & rechaza & frontera: justo por debajo de la longitud mínima \\
\end{longtable}

\evidenciaaqui[0.95]{../evidencias/jflap/fichas/lista2-123.png}{lista2-123 en JFLAP: Multiple Run con las tres cadenas válidas y las
tres inválidas de la tabla anterior.}

\hypertarget{lista-2-ejercicio-124}{%
\subsubsection{Lista 2 · Ejercicio 124}\label{lista-2-ejercicio-124}}

\textbf{Enunciado.} Mecanismo de tres palancas con canicas (Σ = \{0, 1\}); se acepta si la última canica sale por A.

\textbf{Diseño.} Ver el análisis del mecanismo.

\textbf{Archivo.} \texttt{automatas/lista2/lista2-124.jff} · Σ = \{0, 1\} · \textbar Q\textbar{} = 13

\begin{longtable}[]{@{}
  >{\raggedright\arraybackslash}p{(\columnwidth - 6\tabcolsep) * \real{0.1579}}
  >{\raggedright\arraybackslash}p{(\columnwidth - 6\tabcolsep) * \real{0.1579}}
  >{\raggedright\arraybackslash}p{(\columnwidth - 6\tabcolsep) * \real{0.1579}}
  >{\raggedright\arraybackslash}p{(\columnwidth - 6\tabcolsep) * \real{0.5263}}@{}}
\caption{Tabla de transiciones (→ inicial, * aceptación, --- sin transición).}\tabularnewline
\toprule\noalign{}
\begin{minipage}[b]{\linewidth}\raggedright
δ
\end{minipage} & \begin{minipage}[b]{\linewidth}\raggedright
0
\end{minipage} & \begin{minipage}[b]{\linewidth}\raggedright
1
\end{minipage} & \begin{minipage}[b]{\linewidth}\raggedright
Qué recuerda del prefijo leído
\end{minipage} \\
\midrule\noalign{}
\endfirsthead
\toprule\noalign{}
\begin{minipage}[b]{\linewidth}\raggedright
δ
\end{minipage} & \begin{minipage}[b]{\linewidth}\raggedright
0
\end{minipage} & \begin{minipage}[b]{\linewidth}\raggedright
1
\end{minipage} & \begin{minipage}[b]{\linewidth}\raggedright
Qué recuerda del prefijo leído
\end{minipage} \\
\midrule\noalign{}
\endhead
\bottomrule\noalign{}
\endlastfoot
→ III & DIIa & IDDa & palancas izquierda=I, derecha=I, central=I (I: desvía a la izquierda, D: a la derecha); aún no cae ninguna canica \\
* DIIa & IIDa & DDDa & palancas izquierda=D, derecha=I, central=I (I: desvía a la izquierda, D: a la derecha); la última canica salió por A \\
* IDDa & DDDa & IIDr & palancas izquierda=I, derecha=D, central=D (I: desvía a la izquierda, D: a la derecha); la última canica salió por A \\
* IIDa & DIDa & IDIr & palancas izquierda=I, derecha=I, central=D (I: desvía a la izquierda, D: a la derecha); la última canica salió por A \\
* DDDa & IDIr & DIDr & palancas izquierda=D, derecha=D, central=D (I: desvía a la izquierda, D: a la derecha); la última canica salió por A \\
IIDr & DIDa & IDIr & palancas izquierda=I, derecha=I, central=D (I: desvía a la izquierda, D: a la derecha); la última canica salió por R \\
* DIDa & IIIr & DDIr & palancas izquierda=D, derecha=I, central=D (I: desvía a la izquierda, D: a la derecha); la última canica salió por A \\
IDIr & DDIa & IIIr & palancas izquierda=I, derecha=D, central=I (I: desvía a la izquierda, D: a la derecha); la última canica salió por R \\
DIDr & IIIr & DDIr & palancas izquierda=D, derecha=I, central=D (I: desvía a la izquierda, D: a la derecha); la última canica salió por R \\
IIIr & DIIa & IDDa & palancas izquierda=I, derecha=I, central=I (I: desvía a la izquierda, D: a la derecha); la última canica salió por R \\
DDIr & IDDa & DIIr & palancas izquierda=D, derecha=D, central=I (I: desvía a la izquierda, D: a la derecha); la última canica salió por R \\
* DDIa & IDDa & DIIr & palancas izquierda=D, derecha=D, central=I (I: desvía a la izquierda, D: a la derecha); la última canica salió por A \\
DIIr & IIDa & DDDa & palancas izquierda=D, derecha=I, central=I (I: desvía a la izquierda, D: a la derecha); la última canica salió por R \\
\end{longtable}

\begin{longtable}[]{@{}lclll@{}}
\caption{Cadenas de prueba.}\tabularnewline
\toprule\noalign{}
Cadena & Longitud & Esperado & Obtenido & Tipo de caso \\
\midrule\noalign{}
\endfirsthead
\toprule\noalign{}
Cadena & Longitud & Esperado & Obtenido & Tipo de caso \\
\midrule\noalign{}
\endhead
\bottomrule\noalign{}
\endlastfoot
\texttt{0} & 1 & acepta & acepta & válida \\
\texttt{01} & 2 & acepta & acepta & válida \\
\texttt{0010} & 4 & acepta & acepta & válida \\
λ & 0 & rechaza & rechaza & inválida \\
\texttt{1010} & 4 & rechaza & rechaza & inválida \\
\texttt{00101} & 5 & rechaza & rechaza & inválida \\
\end{longtable}

\evidenciaaqui[0.95]{../evidencias/jflap/fichas/lista2-124.png}{lista2-124 en JFLAP: Multiple Run con las tres cadenas válidas y las
tres inválidas de la tabla anterior.}

\hypertarget{lista-2-ejercicio-125}{%
\subsubsection{Lista 2 · Ejercicio 125}\label{lista-2-ejercicio-125}}

\textbf{Enunciado.} Sea A = (Q, Σ, δ, A0, F) un AFD. ¿Cuál es el lenguaje de A' = (Q, Σ, δ, A0, Q ~F)?

\textbf{Diseño.} Como ejemplo se complementa el AFD de 139 (número par de ceros).

\textbf{Archivo.} \texttt{automatas/lista2/lista2-125.jff} · Σ = \{0, 1\} · \textbar Q\textbar{} = 2

\begin{longtable}[]{@{}llll@{}}
\caption{Tabla de transiciones (→ inicial, * aceptación, --- sin transición).}\tabularnewline
\toprule\noalign{}
δ & 0 & 1 & Qué recuerda del prefijo leído \\
\midrule\noalign{}
\endfirsthead
\toprule\noalign{}
δ & 0 & 1 & Qué recuerda del prefijo leído \\
\midrule\noalign{}
\endhead
\bottomrule\noalign{}
\endlastfoot
→ par & impar & par & se ha leído un número par de ceros \\
* impar & par & impar & se ha leído un número impar de ceros \\
\end{longtable}

\begin{longtable}[]{@{}lclll@{}}
\caption{Cadenas de prueba.}\tabularnewline
\toprule\noalign{}
Cadena & Longitud & Esperado & Obtenido & Tipo de caso \\
\midrule\noalign{}
\endfirsthead
\toprule\noalign{}
Cadena & Longitud & Esperado & Obtenido & Tipo de caso \\
\midrule\noalign{}
\endhead
\bottomrule\noalign{}
\endlastfoot
\texttt{0} & 1 & acepta & acepta & válida \\
\texttt{01} & 2 & acepta & acepta & válida \\
\texttt{1101} & 4 & acepta & acepta & válida \\
λ & 0 & rechaza & rechaza & inválida \\
\texttt{1} & 1 & rechaza & rechaza & inválida \\
\texttt{100} & 3 & rechaza & rechaza & inválida \\
\end{longtable}

\evidenciaaqui[0.95]{../evidencias/jflap/fichas/lista2-125.png}{lista2-125 en JFLAP: Multiple Run con las tres cadenas válidas y las
tres inválidas de la tabla anterior.}

\hypertarget{lista-2-ejercicio-126}{%
\subsubsection{Lista 2 · Ejercicio 126}\label{lista-2-ejercicio-126}}

\textbf{Enunciado.} Sea A un AFD. ¿Cuál es el lenguaje del autómata si F = ∅?

\textbf{Diseño.} Ejemplo: el AFD de 139 con F = ∅.

\textbf{Archivo.} \texttt{automatas/lista2/lista2-126.jff} · Σ = \{0, 1\} · \textbar Q\textbar{} = 2

\begin{longtable}[]{@{}llll@{}}
\caption{Tabla de transiciones (→ inicial, * aceptación, --- sin transición).}\tabularnewline
\toprule\noalign{}
δ & 0 & 1 & Qué recuerda del prefijo leído \\
\midrule\noalign{}
\endfirsthead
\toprule\noalign{}
δ & 0 & 1 & Qué recuerda del prefijo leído \\
\midrule\noalign{}
\endhead
\bottomrule\noalign{}
\endlastfoot
→ par & impar & par & número par de ceros \\
impar & par & impar & número impar de ceros \\
\end{longtable}

\begin{longtable}[]{@{}lclll@{}}
\caption{Cadenas de prueba.}\tabularnewline
\toprule\noalign{}
Cadena & Longitud & Esperado & Obtenido & Tipo de caso \\
\midrule\noalign{}
\endfirsthead
\toprule\noalign{}
Cadena & Longitud & Esperado & Obtenido & Tipo de caso \\
\midrule\noalign{}
\endhead
\bottomrule\noalign{}
\endlastfoot
λ & 0 & rechaza & rechaza & inválida \\
\texttt{1} & 1 & rechaza & rechaza & inválida \\
\texttt{001} & 3 & rechaza & rechaza & inválida \\
\end{longtable}

\evidenciaaqui[0.95]{../evidencias/jflap/fichas/lista2-126.png}{lista2-126 en JFLAP: Multiple Run con las tres cadenas válidas y las
tres inválidas de la tabla anterior.}

\hypertarget{lista-2-ejercicio-127}{%
\subsubsection{Lista 2 · Ejercicio 127}\label{lista-2-ejercicio-127}}

\textbf{Enunciado.} Sea A un AFD. ¿Cuál es el lenguaje del autómata si F = Q?

\textbf{Diseño.} Ejemplo: el AFD de 139 con F = Q.

\textbf{Archivo.} \texttt{automatas/lista2/lista2-127.jff} · Σ = \{0, 1\} · \textbar Q\textbar{} = 2

\begin{longtable}[]{@{}llll@{}}
\caption{Tabla de transiciones (→ inicial, * aceptación, --- sin transición).}\tabularnewline
\toprule\noalign{}
δ & 0 & 1 & Qué recuerda del prefijo leído \\
\midrule\noalign{}
\endfirsthead
\toprule\noalign{}
δ & 0 & 1 & Qué recuerda del prefijo leído \\
\midrule\noalign{}
\endhead
\bottomrule\noalign{}
\endlastfoot
→* par & impar & par & número par de ceros \\
* impar & par & impar & número impar de ceros \\
\end{longtable}

\begin{longtable}[]{@{}lclll@{}}
\caption{Cadenas de prueba.}\tabularnewline
\toprule\noalign{}
Cadena & Longitud & Esperado & Obtenido & Tipo de caso \\
\midrule\noalign{}
\endfirsthead
\toprule\noalign{}
Cadena & Longitud & Esperado & Obtenido & Tipo de caso \\
\midrule\noalign{}
\endhead
\bottomrule\noalign{}
\endlastfoot
λ & 0 & acepta & acepta & válida \\
\texttt{1} & 1 & acepta & acepta & válida \\
\texttt{100} & 3 & acepta & acepta & válida \\
\end{longtable}

\evidenciaaqui[0.95]{../evidencias/jflap/fichas/lista2-127.png}{lista2-127 en JFLAP: Multiple Run con las tres cadenas válidas y las
tres inválidas de la tabla anterior.}

\hypertarget{lista-2-ejercicio-128}{%
\subsubsection{Lista 2 · Ejercicio 128}\label{lista-2-ejercicio-128}}

\textbf{Enunciado.} Sea L el lenguaje del autómata del diagrama. Encuentre un AFD para L2 = \{w : wa ∈ L para algún símbolo a\}.

\textbf{Diseño.} Se conservan los estados y las transiciones; un estado q es de aceptación en el nuevo autómata si algún símbolo lo lleva a F. Como δ(D, 0) = F y δ(E, 1) = F, F' = \{D, E\}.

\textbf{Archivo.} \texttt{automatas/lista2/lista2-128.jff} · Σ = \{0, 1\} · \textbar Q\textbar{} = 6

\begin{longtable}[]{@{}llll@{}}
\caption{Tabla de transiciones (→ inicial, * aceptación, --- sin transición).}\tabularnewline
\toprule\noalign{}
δ & 0 & 1 & Qué recuerda del prefijo leído \\
\midrule\noalign{}
\endfirsthead
\toprule\noalign{}
δ & 0 & 1 & Qué recuerda del prefijo leído \\
\midrule\noalign{}
\endhead
\bottomrule\noalign{}
\endlastfoot
→ A & B & C & mismo estado que en el autómata original \\
B & C & D & mismo estado que en el autómata original \\
C & D & E & mismo estado que en el autómata original \\
* D & F & E & mismo estado que en el autómata original \\
* E & E & F & mismo estado que en el autómata original \\
F & D & D & mismo estado que en el autómata original \\
\end{longtable}

\begin{longtable}[]{@{}lclll@{}}
\caption{Cadenas de prueba.}\tabularnewline
\toprule\noalign{}
Cadena & Longitud & Esperado & Obtenido & Tipo de caso \\
\midrule\noalign{}
\endfirsthead
\toprule\noalign{}
Cadena & Longitud & Esperado & Obtenido & Tipo de caso \\
\midrule\noalign{}
\endhead
\bottomrule\noalign{}
\endlastfoot
\texttt{10} & 2 & acepta & acepta & válida \\
\texttt{001} & 3 & acepta & acepta & válida \\
\texttt{10100} & 5 & acepta & acepta & válida \\
λ & 0 & rechaza & rechaza & inválida \\
\texttt{1} & 1 & rechaza & rechaza & inválida \\
\texttt{111} & 3 & rechaza & rechaza & inválida \\
\end{longtable}

\evidenciaaqui[0.95]{../evidencias/jflap/fichas/lista2-128.png}{lista2-128 en JFLAP: Multiple Run con las tres cadenas válidas y las
tres inválidas de la tabla anterior.}

\hypertarget{lista-2-ejercicio-131}{%
\subsubsection{Lista 2 · Ejercicio 131}\label{lista-2-ejercicio-131}}

\textbf{Enunciado.} AFD dado M = (\{q1, q2\}, \{0, 1\}, δ, q1, \{q2\}).

\textbf{Diseño.} q2 significa que el último símbolo leído es 1.

\textbf{Archivo.} \texttt{automatas/lista2/lista2-131.jff} · Σ = \{0, 1\} · \textbar Q\textbar{} = 2

\begin{longtable}[]{@{}lll@{}}
\caption{Tabla de transiciones (→ inicial, * aceptación, --- sin transición).}\tabularnewline
\toprule\noalign{}
δ & 0 & 1 \\
\midrule\noalign{}
\endfirsthead
\toprule\noalign{}
δ & 0 & 1 \\
\midrule\noalign{}
\endhead
\bottomrule\noalign{}
\endlastfoot
→ q1 & q1 & q2 \\
* q2 & q1 & q2 \\
\end{longtable}

\begin{longtable}[]{@{}lclll@{}}
\caption{Cadenas de prueba.}\tabularnewline
\toprule\noalign{}
Cadena & Longitud & Esperado & Obtenido & Tipo de caso \\
\midrule\noalign{}
\endfirsthead
\toprule\noalign{}
Cadena & Longitud & Esperado & Obtenido & Tipo de caso \\
\midrule\noalign{}
\endhead
\bottomrule\noalign{}
\endlastfoot
\texttt{1} & 1 & acepta & acepta & válida \\
\texttt{11} & 2 & acepta & acepta & válida \\
\texttt{0111} & 4 & acepta & acepta & válida \\
λ & 0 & rechaza & rechaza & inválida \\
\texttt{0} & 1 & rechaza & rechaza & inválida \\
\texttt{010} & 3 & rechaza & rechaza & inválida \\
\end{longtable}

\evidenciaaqui[0.95]{../evidencias/jflap/fichas/lista2-131.png}{lista2-131 en JFLAP: Multiple Run con las tres cadenas válidas y las
tres inválidas de la tabla anterior.}

\hypertarget{lista-2-ejercicio-133}{%
\subsubsection{Lista 2 · Ejercicio 133}\label{lista-2-ejercicio-133}}

\textbf{Enunciado.} AFD dado con Q = \{q0, q1, q2, q3\}, Σ = \{a, b\}, F = \{q0\}.

\textbf{Diseño.} q0 = (a par, b par), q1 = (a par, b impar), q2 = (a impar, b par), q3 = (a impar, b impar).

\textbf{Archivo.} \texttt{automatas/lista2/lista2-133.jff} · Σ = \{a, b\} · \textbar Q\textbar{} = 4

\begin{longtable}[]{@{}lll@{}}
\caption{Tabla de transiciones (→ inicial, * aceptación, --- sin transición).}\tabularnewline
\toprule\noalign{}
δ & a & b \\
\midrule\noalign{}
\endfirsthead
\toprule\noalign{}
δ & a & b \\
\midrule\noalign{}
\endhead
\bottomrule\noalign{}
\endlastfoot
→* q0 & q2 & q1 \\
q1 & q3 & q0 \\
q2 & q0 & q3 \\
q3 & q1 & q2 \\
\end{longtable}

\begin{longtable}[]{@{}lclll@{}}
\caption{Cadenas de prueba.}\tabularnewline
\toprule\noalign{}
Cadena & Longitud & Esperado & Obtenido & Tipo de caso \\
\midrule\noalign{}
\endfirsthead
\toprule\noalign{}
Cadena & Longitud & Esperado & Obtenido & Tipo de caso \\
\midrule\noalign{}
\endhead
\bottomrule\noalign{}
\endlastfoot
λ & 0 & acepta & acepta & válida \\
\texttt{baab} & 4 & acepta & acepta & válida \\
\texttt{bbaa} & 4 & acepta & acepta & válida \\
\texttt{a} & 1 & rechaza & rechaza & inválida \\
\texttt{ab} & 2 & rechaza & rechaza & inválida \\
\texttt{baaa} & 4 & rechaza & rechaza & inválida \\
\end{longtable}

\evidenciaaqui[0.95]{../evidencias/jflap/fichas/lista2-133.png}{lista2-133 en JFLAP: Multiple Run con las tres cadenas válidas y las
tres inválidas de la tabla anterior.}

\hypertarget{lista-2-ejercicio-135}{%
\subsubsection{Lista 2 · Ejercicio 135}\label{lista-2-ejercicio-135}}

\textbf{Enunciado.} Dado Σ = \{a, b\}, AFD para L = \{bᵐabⁿ : m, n \textgreater{} 0\}.

\textbf{Diseño.} Cuatro fases; m \textgreater{} 0 y n \textgreater{} 0 exigen al menos una b en cada lado.

\textbf{Archivo.} \texttt{automatas/lista2/lista2-135.jff} · Σ = \{a, b\} · \textbar Q\textbar{} = 5

\begin{longtable}[]{@{}llll@{}}
\caption{Tabla de transiciones (→ inicial, * aceptación, --- sin transición).}\tabularnewline
\toprule\noalign{}
δ & a & b & Qué recuerda del prefijo leído \\
\midrule\noalign{}
\endfirsthead
\toprule\noalign{}
δ & a & b & Qué recuerda del prefijo leído \\
\midrule\noalign{}
\endhead
\bottomrule\noalign{}
\endlastfoot
→ ini & muerto & b+ & no se ha leído nada \\
muerto & muerto & muerto & la cadena no tiene la forma \\
b+ & b+a & b+ & una o más bes iniciales \\
b+a & muerto & b+ab+ & bes y la a \\
* b+ab+ & muerto & b+ab+ & bes, la a y una o más bes \\
\end{longtable}

\begin{longtable}[]{@{}
  >{\raggedright\arraybackslash}p{(\columnwidth - 8\tabcolsep) * \real{0.2000}}
  >{\centering\arraybackslash}p{(\columnwidth - 8\tabcolsep) * \real{0.2000}}
  >{\raggedright\arraybackslash}p{(\columnwidth - 8\tabcolsep) * \real{0.2000}}
  >{\raggedright\arraybackslash}p{(\columnwidth - 8\tabcolsep) * \real{0.2000}}
  >{\raggedright\arraybackslash}p{(\columnwidth - 8\tabcolsep) * \real{0.2000}}@{}}
\caption{Cadenas de prueba.}\tabularnewline
\toprule\noalign{}
\begin{minipage}[b]{\linewidth}\raggedright
Cadena
\end{minipage} & \begin{minipage}[b]{\linewidth}\centering
Longitud
\end{minipage} & \begin{minipage}[b]{\linewidth}\raggedright
Esperado
\end{minipage} & \begin{minipage}[b]{\linewidth}\raggedright
Obtenido
\end{minipage} & \begin{minipage}[b]{\linewidth}\raggedright
Tipo de caso
\end{minipage} \\
\midrule\noalign{}
\endfirsthead
\toprule\noalign{}
\begin{minipage}[b]{\linewidth}\raggedright
Cadena
\end{minipage} & \begin{minipage}[b]{\linewidth}\centering
Longitud
\end{minipage} & \begin{minipage}[b]{\linewidth}\raggedright
Esperado
\end{minipage} & \begin{minipage}[b]{\linewidth}\raggedright
Obtenido
\end{minipage} & \begin{minipage}[b]{\linewidth}\raggedright
Tipo de caso
\end{minipage} \\
\midrule\noalign{}
\endhead
\bottomrule\noalign{}
\endlastfoot
\texttt{bab} & 3 & acepta & acepta & válida \\
\texttt{babb} & 4 & acepta & acepta & válida \\
\texttt{bbbbab} & 6 & acepta & acepta & válida \\
λ & 0 & rechaza & rechaza & inválida \\
\texttt{b} & 1 & rechaza & rechaza & inválida \\
\texttt{bba} & 3 & rechaza & rechaza & inválida \\
\texttt{aa} & 2 & rechaza & rechaza & frontera: justo por debajo de la longitud mínima \\
\end{longtable}

\evidenciaaqui[0.95]{../evidencias/jflap/fichas/lista2-135.png}{lista2-135 en JFLAP: Multiple Run con las tres cadenas válidas y las
tres inválidas de la tabla anterior.}

\hypertarget{lista-2-ejercicio-137}{%
\subsubsection{Lista 2 · Ejercicio 137}\label{lista-2-ejercicio-137}}

\textbf{Enunciado.} AFD que reconozca las cadenas sobre \{a, b\} que comienzan con el prefijo ab.

\textbf{Diseño.} Se verifica el prefijo y después se acepta todo.

\textbf{Archivo.} \texttt{automatas/lista2/lista2-137.jff} · Σ = \{a, b\} · \textbar Q\textbar{} = 4

\begin{longtable}[]{@{}llll@{}}
\caption{Tabla de transiciones (→ inicial, * aceptación, --- sin transición).}\tabularnewline
\toprule\noalign{}
δ & a & b & Qué recuerda del prefijo leído \\
\midrule\noalign{}
\endfirsthead
\toprule\noalign{}
δ & a & b & Qué recuerda del prefijo leído \\
\midrule\noalign{}
\endhead
\bottomrule\noalign{}
\endlastfoot
→ λ & a & muerto & se han leído los primeros 0 símbolos de ab \\
a & muerto & ab & se han leído los primeros 1 símbolos de ab \\
muerto & muerto & muerto & el prefijo leído no coincide con ab \\
* ab & ab & ab & ya se leyó el prefijo ab \\
\end{longtable}

\begin{longtable}[]{@{}
  >{\raggedright\arraybackslash}p{(\columnwidth - 8\tabcolsep) * \real{0.2000}}
  >{\centering\arraybackslash}p{(\columnwidth - 8\tabcolsep) * \real{0.2000}}
  >{\raggedright\arraybackslash}p{(\columnwidth - 8\tabcolsep) * \real{0.2000}}
  >{\raggedright\arraybackslash}p{(\columnwidth - 8\tabcolsep) * \real{0.2000}}
  >{\raggedright\arraybackslash}p{(\columnwidth - 8\tabcolsep) * \real{0.2000}}@{}}
\caption{Cadenas de prueba.}\tabularnewline
\toprule\noalign{}
\begin{minipage}[b]{\linewidth}\raggedright
Cadena
\end{minipage} & \begin{minipage}[b]{\linewidth}\centering
Longitud
\end{minipage} & \begin{minipage}[b]{\linewidth}\raggedright
Esperado
\end{minipage} & \begin{minipage}[b]{\linewidth}\raggedright
Obtenido
\end{minipage} & \begin{minipage}[b]{\linewidth}\raggedright
Tipo de caso
\end{minipage} \\
\midrule\noalign{}
\endfirsthead
\toprule\noalign{}
\begin{minipage}[b]{\linewidth}\raggedright
Cadena
\end{minipage} & \begin{minipage}[b]{\linewidth}\centering
Longitud
\end{minipage} & \begin{minipage}[b]{\linewidth}\raggedright
Esperado
\end{minipage} & \begin{minipage}[b]{\linewidth}\raggedright
Obtenido
\end{minipage} & \begin{minipage}[b]{\linewidth}\raggedright
Tipo de caso
\end{minipage} \\
\midrule\noalign{}
\endhead
\bottomrule\noalign{}
\endlastfoot
\texttt{ab} & 2 & acepta & acepta & válida \\
\texttt{aba} & 3 & acepta & acepta & válida \\
\texttt{abbbb} & 5 & acepta & acepta & válida \\
λ & 0 & rechaza & rechaza & inválida \\
\texttt{b} & 1 & rechaza & rechaza & inválida \\
\texttt{bab} & 3 & rechaza & rechaza & inválida \\
\texttt{a} & 1 & rechaza & rechaza & frontera: justo por debajo de la longitud mínima \\
\end{longtable}

\evidenciaaqui[0.95]{../evidencias/jflap/fichas/lista2-137.png}{lista2-137 en JFLAP: Multiple Run con las tres cadenas válidas y las
tres inválidas de la tabla anterior.}

\hypertarget{lista-2-ejercicio-139}{%
\subsubsection{Lista 2 · Ejercicio 139}\label{lista-2-ejercicio-139}}

\textbf{Enunciado.} Autómata que acepte las cadenas en \{0, 1\}* con un número par de ceros.

\textbf{Diseño.} Paridad de ceros.

\textbf{Archivo.} \texttt{automatas/lista2/lista2-139.jff} · Σ = \{0, 1\} · \textbar Q\textbar{} = 2

\begin{longtable}[]{@{}llll@{}}
\caption{Tabla de transiciones (→ inicial, * aceptación, --- sin transición).}\tabularnewline
\toprule\noalign{}
δ & 0 & 1 & Qué recuerda del prefijo leído \\
\midrule\noalign{}
\endfirsthead
\toprule\noalign{}
δ & 0 & 1 & Qué recuerda del prefijo leído \\
\midrule\noalign{}
\endhead
\bottomrule\noalign{}
\endlastfoot
→* par & impar & par & se ha leído un número par de ceros \\
impar & par & impar & se ha leído un número impar de ceros \\
\end{longtable}

\begin{longtable}[]{@{}lclll@{}}
\caption{Cadenas de prueba.}\tabularnewline
\toprule\noalign{}
Cadena & Longitud & Esperado & Obtenido & Tipo de caso \\
\midrule\noalign{}
\endfirsthead
\toprule\noalign{}
Cadena & Longitud & Esperado & Obtenido & Tipo de caso \\
\midrule\noalign{}
\endhead
\bottomrule\noalign{}
\endlastfoot
λ & 0 & acepta & acepta & válida \\
\texttt{1} & 1 & acepta & acepta & válida \\
\texttt{010} & 3 & acepta & acepta & válida \\
\texttt{0} & 1 & rechaza & rechaza & inválida \\
\texttt{01} & 2 & rechaza & rechaza & inválida \\
\texttt{0100} & 4 & rechaza & rechaza & inválida \\
\end{longtable}

\evidenciaaqui[0.95]{../evidencias/jflap/fichas/lista2-139.png}{lista2-139 en JFLAP: Multiple Run con las tres cadenas válidas y las
tres inválidas de la tabla anterior.}

\hypertarget{lista-2-ejercicio-141}{%
\subsubsection{Lista 2 · Ejercicio 141}\label{lista-2-ejercicio-141}}

\textbf{Enunciado.} L = \{w ∈ \{0, 1\}* : cada 0 en w tiene un 1 inmediatamente a su derecha\}.

\textbf{Diseño.} Un estado de espera después de cada 0.

\textbf{Archivo.} \texttt{automatas/lista2/lista2-141.jff} · Σ = \{0, 1\} · \textbar Q\textbar{} = 3

\begin{longtable}[]{@{}llll@{}}
\caption{Tabla de transiciones (→ inicial, * aceptación, --- sin transición).}\tabularnewline
\toprule\noalign{}
δ & 0 & 1 & Qué recuerda del prefijo leído \\
\midrule\noalign{}
\endfirsthead
\toprule\noalign{}
δ & 0 & 1 & Qué recuerda del prefijo leído \\
\midrule\noalign{}
\endhead
\bottomrule\noalign{}
\endlastfoot
→* ok & 0? & ok & todo 0 leído ya tiene su 1 \\
0? & muerto & ok & el último 0 aún espera su 1 \\
muerto & muerto & muerto & un 0 no fue seguido de 1 \\
\end{longtable}

\begin{longtable}[]{@{}lclll@{}}
\caption{Cadenas de prueba.}\tabularnewline
\toprule\noalign{}
Cadena & Longitud & Esperado & Obtenido & Tipo de caso \\
\midrule\noalign{}
\endfirsthead
\toprule\noalign{}
Cadena & Longitud & Esperado & Obtenido & Tipo de caso \\
\midrule\noalign{}
\endhead
\bottomrule\noalign{}
\endlastfoot
λ & 0 & acepta & acepta & válida \\
\texttt{1} & 1 & acepta & acepta & válida \\
\texttt{111} & 3 & acepta & acepta & válida \\
\texttt{0} & 1 & rechaza & rechaza & inválida \\
\texttt{00} & 2 & rechaza & rechaza & inválida \\
\texttt{0100} & 4 & rechaza & rechaza & inválida \\
\end{longtable}

\evidenciaaqui[0.95]{../evidencias/jflap/fichas/lista2-141.png}{lista2-141 en JFLAP: Multiple Run con las tres cadenas válidas y las
tres inválidas de la tabla anterior.}

\hypertarget{lista-2-ejercicio-143}{%
\subsubsection{Lista 2 · Ejercicio 143}\label{lista-2-ejercicio-143}}

\textbf{Enunciado.} Determine los lenguajes de los autómatas de las Figuras (a) y (b).

Las figuras del enunciado no aparecen en el PDF de la lista.

\hypertarget{lista-2-ejercicio-145}{%
\subsubsection{Lista 2 · Ejercicio 145}\label{lista-2-ejercicio-145}}

\textbf{Enunciado.} AFD que lee un número binario de derecha a izquierda e identifica los múltiplos de 5.

\textbf{Diseño.} El primer bit leído es el menos significativo. Si los bits restantes forman (de derecha a izquierda) el número V, leer el bit b hace que el valor total sea b + 2V; así, si se necesita que el total sea ≡ t, el resto debe valer V ≡ (t − b)·2⁻¹ = 3(t − b) (mod 5). El estado guarda ese valor t pendiente: se empieza pidiendo t = 0 y se acepta si, al terminar, lo pendiente es 0. La cadena vacía no es un número y se rechaza.

\textbf{Archivo.} \texttt{automatas/lista2/lista2-145.jff} · Σ = \{0, 1\} · \textbar Q\textbar{} = 6

\begin{longtable}[]{@{}
  >{\raggedright\arraybackslash}p{(\columnwidth - 6\tabcolsep) * \real{0.1579}}
  >{\raggedright\arraybackslash}p{(\columnwidth - 6\tabcolsep) * \real{0.1579}}
  >{\raggedright\arraybackslash}p{(\columnwidth - 6\tabcolsep) * \real{0.1579}}
  >{\raggedright\arraybackslash}p{(\columnwidth - 6\tabcolsep) * \real{0.5263}}@{}}
\caption{Tabla de transiciones (→ inicial, * aceptación, --- sin transición).}\tabularnewline
\toprule\noalign{}
\begin{minipage}[b]{\linewidth}\raggedright
δ
\end{minipage} & \begin{minipage}[b]{\linewidth}\raggedright
0
\end{minipage} & \begin{minipage}[b]{\linewidth}\raggedright
1
\end{minipage} & \begin{minipage}[b]{\linewidth}\raggedright
Qué recuerda del prefijo leído
\end{minipage} \\
\midrule\noalign{}
\endfirsthead
\toprule\noalign{}
\begin{minipage}[b]{\linewidth}\raggedright
δ
\end{minipage} & \begin{minipage}[b]{\linewidth}\raggedright
0
\end{minipage} & \begin{minipage}[b]{\linewidth}\raggedright
1
\end{minipage} & \begin{minipage}[b]{\linewidth}\raggedright
Qué recuerda del prefijo leído
\end{minipage} \\
\midrule\noalign{}
\endhead
\bottomrule\noalign{}
\endlastfoot
→ ini & t0 & t2 & no se ha leído ningún bit \\
* t0 & t0 & t2 & los bits que faltan, leídos como número desde su bit menos significativo, deben valer ≡ 0 (mod 5) \\
t2 & t1 & t3 & los bits que faltan, leídos como número desde su bit menos significativo, deben valer ≡ 2 (mod 5) \\
t1 & t3 & t0 & los bits que faltan, leídos como número desde su bit menos significativo, deben valer ≡ 1 (mod 5) \\
t3 & t4 & t1 & los bits que faltan, leídos como número desde su bit menos significativo, deben valer ≡ 3 (mod 5) \\
t4 & t2 & t4 & los bits que faltan, leídos como número desde su bit menos significativo, deben valer ≡ 4 (mod 5) \\
\end{longtable}

\begin{longtable}[]{@{}lclll@{}}
\caption{Cadenas de prueba.}\tabularnewline
\toprule\noalign{}
Cadena & Longitud & Esperado & Obtenido & Tipo de caso \\
\midrule\noalign{}
\endfirsthead
\toprule\noalign{}
Cadena & Longitud & Esperado & Obtenido & Tipo de caso \\
\midrule\noalign{}
\endhead
\bottomrule\noalign{}
\endlastfoot
\texttt{0} & 1 & acepta & acepta & válida \\
\texttt{00} & 2 & acepta & acepta & válida \\
\texttt{1111} & 4 & acepta & acepta & válida \\
λ & 0 & rechaza & rechaza & inválida \\
\texttt{1} & 1 & rechaza & rechaza & inválida \\
\texttt{111} & 3 & rechaza & rechaza & inválida \\
\end{longtable}

\evidenciaaqui[0.95]{../evidencias/jflap/fichas/lista2-145.png}{lista2-145 en JFLAP: Multiple Run con las tres cadenas válidas y las
tres inválidas de la tabla anterior.}

\hypertarget{lista-2-ejercicio-147}{%
\subsubsection{Lista 2 · Ejercicio 147}\label{lista-2-ejercicio-147}}

\textbf{Enunciado.} AFD que lee un número binario de izquierda a derecha y lo acepta si es múltiplo de 6.

\textbf{Diseño.} Leer un bit b convierte el valor v en 2v + b, así que basta el residuo mod 6. La cadena vacía no representa un número y se rechaza. Este AFD de 7 estados se puede minimizar a 4: r y r + 3 tienen el mismo sucesor (2(r + 3) ≡ 2r), de modo que \{r1, r4\}, \{r2, r5\} e \{ini, r3\} son equivalentes; es la caracterización «múltiplo de 3 y par».

\textbf{Archivo.} \texttt{automatas/lista2/lista2-147.jff} · Σ = \{0, 1\} · \textbar Q\textbar{} = 7

\begin{longtable}[]{@{}llll@{}}
\caption{Tabla de transiciones (→ inicial, * aceptación, --- sin transición).}\tabularnewline
\toprule\noalign{}
δ & 0 & 1 & Qué recuerda del prefijo leído \\
\midrule\noalign{}
\endfirsthead
\toprule\noalign{}
δ & 0 & 1 & Qué recuerda del prefijo leído \\
\midrule\noalign{}
\endhead
\bottomrule\noalign{}
\endlastfoot
→ ini & r0 & r1 & no se ha leído ningún bit \\
* r0 & r0 & r1 & el número leído es ≡ 0 (mod 6) \\
r1 & r2 & r3 & el número leído es ≡ 1 (mod 6) \\
r2 & r4 & r5 & el número leído es ≡ 2 (mod 6) \\
r3 & r0 & r1 & el número leído es ≡ 3 (mod 6) \\
r4 & r2 & r3 & el número leído es ≡ 4 (mod 6) \\
r5 & r4 & r5 & el número leído es ≡ 5 (mod 6) \\
\end{longtable}

\begin{longtable}[]{@{}lclll@{}}
\caption{Cadenas de prueba.}\tabularnewline
\toprule\noalign{}
Cadena & Longitud & Esperado & Obtenido & Tipo de caso \\
\midrule\noalign{}
\endfirsthead
\toprule\noalign{}
Cadena & Longitud & Esperado & Obtenido & Tipo de caso \\
\midrule\noalign{}
\endhead
\bottomrule\noalign{}
\endlastfoot
\texttt{0} & 1 & acepta & acepta & válida \\
\texttt{00} & 2 & acepta & acepta & válida \\
\texttt{0000} & 4 & acepta & acepta & válida \\
λ & 0 & rechaza & rechaza & inválida \\
\texttt{1} & 1 & rechaza & rechaza & inválida \\
\texttt{010} & 3 & rechaza & rechaza & inválida \\
\end{longtable}

\evidenciaaqui[0.95]{../evidencias/jflap/fichas/lista2-147.png}{lista2-147 en JFLAP: Multiple Run con las tres cadenas válidas y las
tres inválidas de la tabla anterior.}

\hypertarget{lista-2-ejercicio-148}{%
\subsubsection{Lista 2 · Ejercicio 148}\label{lista-2-ejercicio-148}}

\textbf{Enunciado.} AFD del progreso de un alumno en Teoría de la Computación (convocatorias de febrero, septiembre y diciembre; Σ = \{P, N, A, S\}; no más de dos presentaciones por año; termina al aprobar).

\textbf{Diseño.} Interpretación: en cada convocatoria el alumno no se presenta (N) o se presenta (P) y después aprueba (A) o suspende (S). Se cuentan las presentaciones del año; con dos, en la convocatoria restante sólo cabe N. Después de diciembre empieza un año nuevo con el contador en cero. APROBADO es el único estado de aceptación y cualquier símbolo posterior lleva a muerto.

\textbf{Archivo.} \texttt{automatas/lista2/lista2-148.jff} · Σ = \{A, N, P, S\} · \textbar Q\textbar{} = 13

\begin{longtable}[]{@{}
  >{\raggedright\arraybackslash}p{(\columnwidth - 10\tabcolsep) * \real{0.1200}}
  >{\raggedright\arraybackslash}p{(\columnwidth - 10\tabcolsep) * \real{0.1200}}
  >{\raggedright\arraybackslash}p{(\columnwidth - 10\tabcolsep) * \real{0.1200}}
  >{\raggedright\arraybackslash}p{(\columnwidth - 10\tabcolsep) * \real{0.1200}}
  >{\raggedright\arraybackslash}p{(\columnwidth - 10\tabcolsep) * \real{0.1200}}
  >{\raggedright\arraybackslash}p{(\columnwidth - 10\tabcolsep) * \real{0.4000}}@{}}
\caption{Tabla de transiciones (→ inicial, * aceptación, --- sin transición).}\tabularnewline
\toprule\noalign{}
\begin{minipage}[b]{\linewidth}\raggedright
δ
\end{minipage} & \begin{minipage}[b]{\linewidth}\raggedright
A
\end{minipage} & \begin{minipage}[b]{\linewidth}\raggedright
N
\end{minipage} & \begin{minipage}[b]{\linewidth}\raggedright
P
\end{minipage} & \begin{minipage}[b]{\linewidth}\raggedright
S
\end{minipage} & \begin{minipage}[b]{\linewidth}\raggedright
Qué recuerda del prefijo leído
\end{minipage} \\
\midrule\noalign{}
\endfirsthead
\toprule\noalign{}
\begin{minipage}[b]{\linewidth}\raggedright
δ
\end{minipage} & \begin{minipage}[b]{\linewidth}\raggedright
A
\end{minipage} & \begin{minipage}[b]{\linewidth}\raggedright
N
\end{minipage} & \begin{minipage}[b]{\linewidth}\raggedright
P
\end{minipage} & \begin{minipage}[b]{\linewidth}\raggedright
S
\end{minipage} & \begin{minipage}[b]{\linewidth}\raggedright
Qué recuerda del prefijo leído
\end{minipage} \\
\midrule\noalign{}
\endhead
\bottomrule\noalign{}
\endlastfoot
→ feb0 & muerto & sep0 & feb0? & muerto & convocatoria de feb, 0 presentaciones usadas en el año; decide si se presenta \\
muerto & muerto & muerto & muerto & muerto & secuencia inválida \\
sep0 & muerto & dic0 & sep0? & muerto & convocatoria de sep, 0 presentaciones usadas en el año; decide si se presenta \\
feb0? & APROBADO & muerto & muerto & sep1 & convocatoria de feb, 0 presentaciones usadas en el año; se presentó y espera el resultado \\
dic0 & muerto & feb0 & dic0? & muerto & convocatoria de dic, 0 presentaciones usadas en el año; decide si se presenta \\
sep0? & APROBADO & muerto & muerto & dic1 & convocatoria de sep, 0 presentaciones usadas en el año; se presentó y espera el resultado \\
* APROBADO & muerto & muerto & muerto & muerto & aprobó el curso \\
sep1 & muerto & dic1 & sep1? & muerto & convocatoria de sep, 1 presentaciones usadas en el año; decide si se presenta \\
dic0? & APROBADO & muerto & muerto & feb0 & convocatoria de dic, 0 presentaciones usadas en el año; se presentó y espera el resultado \\
dic1 & muerto & feb0 & dic1? & muerto & convocatoria de dic, 1 presentaciones usadas en el año; decide si se presenta \\
sep1? & APROBADO & muerto & muerto & dic2 & convocatoria de sep, 1 presentaciones usadas en el año; se presentó y espera el resultado \\
dic1? & APROBADO & muerto & muerto & feb0 & convocatoria de dic, 1 presentaciones usadas en el año; se presentó y espera el resultado \\
dic2 & muerto & feb0 & muerto & muerto & convocatoria de dic, 2 presentaciones usadas en el año; decide si se presenta \\
\end{longtable}

\begin{longtable}[]{@{}
  >{\raggedright\arraybackslash}p{(\columnwidth - 8\tabcolsep) * \real{0.2000}}
  >{\centering\arraybackslash}p{(\columnwidth - 8\tabcolsep) * \real{0.2000}}
  >{\raggedright\arraybackslash}p{(\columnwidth - 8\tabcolsep) * \real{0.2000}}
  >{\raggedright\arraybackslash}p{(\columnwidth - 8\tabcolsep) * \real{0.2000}}
  >{\raggedright\arraybackslash}p{(\columnwidth - 8\tabcolsep) * \real{0.2000}}@{}}
\caption{Cadenas de prueba.}\tabularnewline
\toprule\noalign{}
\begin{minipage}[b]{\linewidth}\raggedright
Cadena
\end{minipage} & \begin{minipage}[b]{\linewidth}\centering
Longitud
\end{minipage} & \begin{minipage}[b]{\linewidth}\raggedright
Esperado
\end{minipage} & \begin{minipage}[b]{\linewidth}\raggedright
Obtenido
\end{minipage} & \begin{minipage}[b]{\linewidth}\raggedright
Tipo de caso
\end{minipage} \\
\midrule\noalign{}
\endfirsthead
\toprule\noalign{}
\begin{minipage}[b]{\linewidth}\raggedright
Cadena
\end{minipage} & \begin{minipage}[b]{\linewidth}\centering
Longitud
\end{minipage} & \begin{minipage}[b]{\linewidth}\raggedright
Esperado
\end{minipage} & \begin{minipage}[b]{\linewidth}\raggedright
Obtenido
\end{minipage} & \begin{minipage}[b]{\linewidth}\raggedright
Tipo de caso
\end{minipage} \\
\midrule\noalign{}
\endhead
\bottomrule\noalign{}
\endlastfoot
\texttt{PA} & 2 & acepta & acepta & válida \\
\texttt{NPA} & 3 & acepta & acepta & válida \\
\texttt{NPSPA} & 5 & acepta & acepta & válida \\
λ & 0 & rechaza & rechaza & inválida \\
\texttt{N} & 1 & rechaza & rechaza & inválida \\
\texttt{PAS} & 3 & rechaza & rechaza & inválida \\
\texttt{P} & 1 & rechaza & rechaza & frontera: justo por debajo de la longitud mínima \\
\end{longtable}

\evidenciaaqui[0.95]{../evidencias/jflap/fichas/lista2-148.png}{lista2-148 en JFLAP: Multiple Run con las tres cadenas válidas y las
tres inválidas de la tabla anterior.}

\hypertarget{lista-2-ejercicio-149}{%
\subsubsection{Lista 2 · Ejercicio 149}\label{lista-2-ejercicio-149}}

\textbf{Enunciado.} AFD de la tabla (→A, C final). Halle δ(B, 0), δ(C, 1), δ̂(A, 1101), δ̂(A, 01001).

\textbf{Diseño.} Q = \{A, B, C\}, Σ = \{0, 1\}, A0 = A, F = \{C\}.

\textbf{Archivo.} \texttt{automatas/lista2/lista2-149.jff} · Σ = \{0, 1\} · \textbar Q\textbar{} = 3

\begin{longtable}[]{@{}lll@{}}
\caption{Tabla de transiciones (→ inicial, * aceptación, --- sin transición).}\tabularnewline
\toprule\noalign{}
δ & 0 & 1 \\
\midrule\noalign{}
\endfirsthead
\toprule\noalign{}
δ & 0 & 1 \\
\midrule\noalign{}
\endhead
\bottomrule\noalign{}
\endlastfoot
→ A & B & C \\
B & A & C \\
* C & C & A \\
\end{longtable}

\begin{longtable}[]{@{}lclll@{}}
\caption{Cadenas de prueba.}\tabularnewline
\toprule\noalign{}
Cadena & Longitud & Esperado & Obtenido & Tipo de caso \\
\midrule\noalign{}
\endfirsthead
\toprule\noalign{}
Cadena & Longitud & Esperado & Obtenido & Tipo de caso \\
\midrule\noalign{}
\endhead
\bottomrule\noalign{}
\endlastfoot
\texttt{1} & 1 & acepta & acepta & válida \\
\texttt{10} & 2 & acepta & acepta & válida \\
\texttt{1110} & 4 & acepta & acepta & válida \\
λ & 0 & rechaza & rechaza & inválida \\
\texttt{0} & 1 & rechaza & rechaza & inválida \\
\texttt{011} & 3 & rechaza & rechaza & inválida \\
\end{longtable}

\evidenciaaqui[0.95]{../evidencias/jflap/fichas/lista2-149.png}{lista2-149 en JFLAP: Multiple Run con las tres cadenas válidas y las
tres inválidas de la tabla anterior.}

\hypertarget{lista-2-ejercicio-150}{%
\subsubsection{Lista 2 · Ejercicio 150}\label{lista-2-ejercicio-150}}

\textbf{Enunciado.} AFD de la tabla (→A; B, C, D finales). Halle δ̂(B, 10a11), δ̂(B, aa1100), δ̂(A, a01a01), δ̂(C, a11a00).

\textbf{Diseño.} Q = \{A, B, C, D\}, Σ = \{0, 1, a\}, A0 = A, F = \{B, C, D\}.

\textbf{Archivo.} \texttt{automatas/lista2/lista2-150.jff} · Σ = \{0, 1, a\} · \textbar Q\textbar{} = 4

\begin{longtable}[]{@{}llll@{}}
\caption{Tabla de transiciones (→ inicial, * aceptación, --- sin transición).}\tabularnewline
\toprule\noalign{}
δ & 0 & 1 & a \\
\midrule\noalign{}
\endfirsthead
\toprule\noalign{}
δ & 0 & 1 & a \\
\midrule\noalign{}
\endhead
\bottomrule\noalign{}
\endlastfoot
→ A & B & C & D \\
* B & B & C & C \\
* C & A & A & D \\
* D & B & B & B \\
\end{longtable}

\begin{longtable}[]{@{}lclll@{}}
\caption{Cadenas de prueba.}\tabularnewline
\toprule\noalign{}
Cadena & Longitud & Esperado & Obtenido & Tipo de caso \\
\midrule\noalign{}
\endfirsthead
\toprule\noalign{}
Cadena & Longitud & Esperado & Obtenido & Tipo de caso \\
\midrule\noalign{}
\endhead
\bottomrule\noalign{}
\endlastfoot
\texttt{1} & 1 & acepta & acepta & válida \\
\texttt{aa} & 2 & acepta & acepta & válida \\
\texttt{1a01} & 4 & acepta & acepta & válida \\
λ & 0 & rechaza & rechaza & inválida \\
\texttt{01011} & 5 & rechaza & rechaza & inválida \\
\texttt{0aa10a0} & 7 & rechaza & rechaza & inválida \\
\end{longtable}

\evidenciaaqui[0.95]{../evidencias/jflap/fichas/lista2-150.png}{lista2-150 en JFLAP: Multiple Run con las tres cadenas válidas y las
tres inválidas de la tabla anterior.}

\hypertarget{lista-3}{%
\subsection{Lista 3}\label{lista-3}}

\hypertarget{lista-3-ejercicio-151}{%
\subsubsection{Lista 3 · Ejercicio 151}\label{lista-3-ejercicio-151}}

\textbf{Enunciado.} Diagrama y simplificación del AFD de la tabla 151.

\textbf{Diseño.} Primero se eliminan los estados que no se alcanzan desde el inicial; después se refina la partición \{F, Q − F\} hasta que se estabiliza.

\textbf{Archivo.} \texttt{automatas/lista3/lista3-151.jff} · Σ = \{0, 1\} · \textbar Q\textbar{} = 10

\begin{longtable}[]{@{}lll@{}}
\caption{Tabla de transiciones (→ inicial, * aceptación, --- sin transición).}\tabularnewline
\toprule\noalign{}
δ & 0 & 1 \\
\midrule\noalign{}
\endfirsthead
\toprule\noalign{}
δ & 0 & 1 \\
\midrule\noalign{}
\endhead
\bottomrule\noalign{}
\endlastfoot
→ A & C & E \\
B & F & E \\
* C & A & G \\
D & C & E \\
* E & I & B \\
* F & D & H \\
G & F & I \\
H & C & J \\
I & J & E \\
J & I & E \\
\end{longtable}

\begin{longtable}[]{@{}lclll@{}}
\caption{Cadenas de prueba.}\tabularnewline
\toprule\noalign{}
Cadena & Longitud & Esperado & Obtenido & Tipo de caso \\
\midrule\noalign{}
\endfirsthead
\toprule\noalign{}
Cadena & Longitud & Esperado & Obtenido & Tipo de caso \\
\midrule\noalign{}
\endhead
\bottomrule\noalign{}
\endlastfoot
\texttt{1} & 1 & acepta & acepta & válida \\
\texttt{0111} & 4 & acepta & acepta & válida \\
\texttt{1001} & 4 & acepta & acepta & válida \\
λ & 0 & rechaza & rechaza & inválida \\
\texttt{001001000} & 9 & rechaza & rechaza & inválida \\
\texttt{100100000} & 9 & rechaza & rechaza & inválida \\
\end{longtable}

\evidenciaaqui[0.95]{../evidencias/jflap/fichas/lista3-151.png}{lista3-151 en JFLAP: Multiple Run con las tres cadenas válidas y las
tres inválidas de la tabla anterior.}

\hypertarget{lista-3-ejercicio-152}{%
\subsubsection{Lista 3 · Ejercicio 152}\label{lista-3-ejercicio-152}}

\textbf{Enunciado.} Diagrama y simplificación del AFD de la tabla 152.

\textbf{Diseño.} Primero se eliminan los estados que no se alcanzan desde el inicial; después se refina la partición \{F, Q − F\} hasta que se estabiliza.

\textbf{Archivo.} \texttt{automatas/lista3/lista3-152.jff} · Σ = \{0, 1, 2\} · \textbar Q\textbar{} = 6

\begin{longtable}[]{@{}llll@{}}
\caption{Tabla de transiciones (→ inicial, * aceptación, --- sin transición).}\tabularnewline
\toprule\noalign{}
δ & 0 & 1 & 2 \\
\midrule\noalign{}
\endfirsthead
\toprule\noalign{}
δ & 0 & 1 & 2 \\
\midrule\noalign{}
\endhead
\bottomrule\noalign{}
\endlastfoot
→ A & C & E & B \\
B & D & E & B \\
C & A & D & F \\
* D & F & B & C \\
* E & F & B & A \\
* F & D & B & C \\
\end{longtable}

\begin{longtable}[]{@{}lclll@{}}
\caption{Cadenas de prueba.}\tabularnewline
\toprule\noalign{}
Cadena & Longitud & Esperado & Obtenido & Tipo de caso \\
\midrule\noalign{}
\endfirsthead
\toprule\noalign{}
Cadena & Longitud & Esperado & Obtenido & Tipo de caso \\
\midrule\noalign{}
\endhead
\bottomrule\noalign{}
\endlastfoot
\texttt{1} & 1 & acepta & acepta & válida \\
\texttt{02} & 2 & acepta & acepta & válida \\
\texttt{2021} & 4 & acepta & acepta & válida \\
λ & 0 & rechaza & rechaza & inválida \\
\texttt{0} & 1 & rechaza & rechaza & inválida \\
\texttt{211} & 3 & rechaza & rechaza & inválida \\
\end{longtable}

\evidenciaaqui[0.95]{../evidencias/jflap/fichas/lista3-152.png}{lista3-152 en JFLAP: Multiple Run con las tres cadenas válidas y las
tres inválidas de la tabla anterior.}

\hypertarget{lista-3-ejercicio-153}{%
\subsubsection{Lista 3 · Ejercicio 153}\label{lista-3-ejercicio-153}}

\textbf{Enunciado.} Diagrama y simplificación del AFD de la tabla 153.

\textbf{Diseño.} Primero se eliminan los estados que no se alcanzan desde el inicial; después se refina la partición \{F, Q − F\} hasta que se estabiliza.

\textbf{Archivo.} \texttt{automatas/lista3/lista3-153.jff} · Σ = \{0, 1\} · \textbar Q\textbar{} = 7

\begin{longtable}[]{@{}lll@{}}
\caption{Tabla de transiciones (→ inicial, * aceptación, --- sin transición).}\tabularnewline
\toprule\noalign{}
δ & 0 & 1 \\
\midrule\noalign{}
\endfirsthead
\toprule\noalign{}
δ & 0 & 1 \\
\midrule\noalign{}
\endhead
\bottomrule\noalign{}
\endlastfoot
→ A & B & C \\
B & D & E \\
C & F & G \\
* D & D & E \\
E & F & G \\
* F & E & D \\
* G & F & G \\
\end{longtable}

\begin{longtable}[]{@{}
  >{\raggedright\arraybackslash}p{(\columnwidth - 8\tabcolsep) * \real{0.2000}}
  >{\centering\arraybackslash}p{(\columnwidth - 8\tabcolsep) * \real{0.2000}}
  >{\raggedright\arraybackslash}p{(\columnwidth - 8\tabcolsep) * \real{0.2000}}
  >{\raggedright\arraybackslash}p{(\columnwidth - 8\tabcolsep) * \real{0.2000}}
  >{\raggedright\arraybackslash}p{(\columnwidth - 8\tabcolsep) * \real{0.2000}}@{}}
\caption{Cadenas de prueba.}\tabularnewline
\toprule\noalign{}
\begin{minipage}[b]{\linewidth}\raggedright
Cadena
\end{minipage} & \begin{minipage}[b]{\linewidth}\centering
Longitud
\end{minipage} & \begin{minipage}[b]{\linewidth}\raggedright
Esperado
\end{minipage} & \begin{minipage}[b]{\linewidth}\raggedright
Obtenido
\end{minipage} & \begin{minipage}[b]{\linewidth}\raggedright
Tipo de caso
\end{minipage} \\
\midrule\noalign{}
\endfirsthead
\toprule\noalign{}
\begin{minipage}[b]{\linewidth}\raggedright
Cadena
\end{minipage} & \begin{minipage}[b]{\linewidth}\centering
Longitud
\end{minipage} & \begin{minipage}[b]{\linewidth}\raggedright
Esperado
\end{minipage} & \begin{minipage}[b]{\linewidth}\raggedright
Obtenido
\end{minipage} & \begin{minipage}[b]{\linewidth}\raggedright
Tipo de caso
\end{minipage} \\
\midrule\noalign{}
\endhead
\bottomrule\noalign{}
\endlastfoot
\texttt{10} & 2 & acepta & acepta & válida \\
\texttt{101} & 3 & acepta & acepta & válida \\
\texttt{01101} & 5 & acepta & acepta & válida \\
λ & 0 & rechaza & rechaza & inválida \\
\texttt{1} & 1 & rechaza & rechaza & inválida \\
\texttt{100} & 3 & rechaza & rechaza & inválida \\
\texttt{0} & 1 & rechaza & rechaza & frontera: justo por debajo de la longitud mínima \\
\end{longtable}

\evidenciaaqui[0.95]{../evidencias/jflap/fichas/lista3-153.png}{lista3-153 en JFLAP: Multiple Run con las tres cadenas válidas y las
tres inválidas de la tabla anterior.}

\hypertarget{lista-3-ejercicio-154}{%
\subsubsection{Lista 3 · Ejercicio 154}\label{lista-3-ejercicio-154}}

\textbf{Enunciado.} Diagrama y simplificación del AFD de la tabla 154.

\textbf{Diseño.} Primero se eliminan los estados que no se alcanzan desde el inicial; después se refina la partición \{F, Q − F\} hasta que se estabiliza.

\textbf{Archivo.} \texttt{automatas/lista3/lista3-154.jff} · Σ = \{0, 1\} · \textbar Q\textbar{} = 11

\begin{longtable}[]{@{}lll@{}}
\caption{Tabla de transiciones (→ inicial, * aceptación, --- sin transición).}\tabularnewline
\toprule\noalign{}
δ & 0 & 1 \\
\midrule\noalign{}
\endfirsthead
\toprule\noalign{}
δ & 0 & 1 \\
\midrule\noalign{}
\endhead
\bottomrule\noalign{}
\endlastfoot
→ A & D & K \\
B & H & K \\
C & I & E \\
D & C & G \\
* E & K & I \\
F & B & E \\
* G & C & I \\
* H & C & B \\
I & G & K \\
* J & F & B \\
K & C & F \\
\end{longtable}

\begin{longtable}[]{@{}lclll@{}}
\caption{Cadenas de prueba.}\tabularnewline
\toprule\noalign{}
Cadena & Longitud & Esperado & Obtenido & Tipo de caso \\
\midrule\noalign{}
\endfirsthead
\toprule\noalign{}
Cadena & Longitud & Esperado & Obtenido & Tipo de caso \\
\midrule\noalign{}
\endhead
\bottomrule\noalign{}
\endlastfoot
\texttt{01} & 2 & acepta & acepta & válida \\
\texttt{001} & 3 & acepta & acepta & válida \\
\texttt{10110} & 5 & acepta & acepta & válida \\
λ & 0 & rechaza & rechaza & inválida \\
\texttt{1} & 1 & rechaza & rechaza & inválida \\
\texttt{100} & 3 & rechaza & rechaza & inválida \\
\end{longtable}

\evidenciaaqui[0.95]{../evidencias/jflap/fichas/lista3-154.png}{lista3-154 en JFLAP: Multiple Run con las tres cadenas válidas y las
tres inválidas de la tabla anterior.}

\hypertarget{lista-3-ejercicio-155}{%
\subsubsection{Lista 3 · Ejercicio 155}\label{lista-3-ejercicio-155}}

\textbf{Enunciado.} Diagrama y simplificación del AFD de la tabla 155.

\textbf{Diseño.} Primero se eliminan los estados que no se alcanzan desde el inicial; después se refina la partición \{F, Q − F\} hasta que se estabiliza.

\textbf{Archivo.} \texttt{automatas/lista3/lista3-155.jff} · Σ = \{0, 1\} · \textbar Q\textbar{} = 12

\begin{longtable}[]{@{}lll@{}}
\caption{Tabla de transiciones (→ inicial, * aceptación, --- sin transición).}\tabularnewline
\toprule\noalign{}
δ & 0 & 1 \\
\midrule\noalign{}
\endfirsthead
\toprule\noalign{}
δ & 0 & 1 \\
\midrule\noalign{}
\endhead
\bottomrule\noalign{}
\endlastfoot
→ A & B & C \\
B & D & E \\
C & K & G \\
D & H & E \\
E & F & D \\
* F & J & G \\
G & E & G \\
H & B & K \\
I & C & A \\
J & D & G \\
K & F & D \\
L & C & J \\
\end{longtable}

\begin{longtable}[]{@{}
  >{\raggedright\arraybackslash}p{(\columnwidth - 8\tabcolsep) * \real{0.2000}}
  >{\centering\arraybackslash}p{(\columnwidth - 8\tabcolsep) * \real{0.2000}}
  >{\raggedright\arraybackslash}p{(\columnwidth - 8\tabcolsep) * \real{0.2000}}
  >{\raggedright\arraybackslash}p{(\columnwidth - 8\tabcolsep) * \real{0.2000}}
  >{\raggedright\arraybackslash}p{(\columnwidth - 8\tabcolsep) * \real{0.2000}}@{}}
\caption{Cadenas de prueba.}\tabularnewline
\toprule\noalign{}
\begin{minipage}[b]{\linewidth}\raggedright
Cadena
\end{minipage} & \begin{minipage}[b]{\linewidth}\centering
Longitud
\end{minipage} & \begin{minipage}[b]{\linewidth}\raggedright
Esperado
\end{minipage} & \begin{minipage}[b]{\linewidth}\raggedright
Obtenido
\end{minipage} & \begin{minipage}[b]{\linewidth}\raggedright
Tipo de caso
\end{minipage} \\
\midrule\noalign{}
\endfirsthead
\toprule\noalign{}
\begin{minipage}[b]{\linewidth}\raggedright
Cadena
\end{minipage} & \begin{minipage}[b]{\linewidth}\centering
Longitud
\end{minipage} & \begin{minipage}[b]{\linewidth}\raggedright
Esperado
\end{minipage} & \begin{minipage}[b]{\linewidth}\raggedright
Obtenido
\end{minipage} & \begin{minipage}[b]{\linewidth}\raggedright
Tipo de caso
\end{minipage} \\
\midrule\noalign{}
\endhead
\bottomrule\noalign{}
\endlastfoot
\texttt{010} & 3 & acepta & acepta & válida \\
\texttt{1100} & 4 & acepta & acepta & válida \\
\texttt{101010} & 6 & acepta & acepta & válida \\
λ & 0 & rechaza & rechaza & inválida \\
\texttt{0} & 1 & rechaza & rechaza & inválida \\
\texttt{110} & 3 & rechaza & rechaza & inválida \\
\texttt{01} & 2 & rechaza & rechaza & frontera: justo por debajo de la longitud mínima \\
\end{longtable}

\evidenciaaqui[0.95]{../evidencias/jflap/fichas/lista3-155.png}{lista3-155 en JFLAP: Multiple Run con las tres cadenas válidas y las
tres inválidas de la tabla anterior.}

\hypertarget{lista-3-ejercicio-156}{%
\subsubsection{Lista 3 · Ejercicio 156}\label{lista-3-ejercicio-156}}

\textbf{Enunciado.} Diagrama y simplificación del AFD de la tabla 156 (tabla transcrita del diagrama).

\textbf{Diseño.} Primero se eliminan los estados que no se alcanzan desde el inicial; después se refina la partición \{F, Q − F\} hasta que se estabiliza.

\textbf{Archivo.} \texttt{automatas/lista3/lista3-156.jff} · Σ = \{0, 1\} · \textbar Q\textbar{} = 5

\begin{longtable}[]{@{}lll@{}}
\caption{Tabla de transiciones (→ inicial, * aceptación, --- sin transición).}\tabularnewline
\toprule\noalign{}
δ & 0 & 1 \\
\midrule\noalign{}
\endfirsthead
\toprule\noalign{}
δ & 0 & 1 \\
\midrule\noalign{}
\endhead
\bottomrule\noalign{}
\endlastfoot
→ A & B & C \\
B & C & E \\
C & D & C \\
D & C & E \\
* E & B & E \\
\end{longtable}

\begin{longtable}[]{@{}lclll@{}}
\caption{Cadenas de prueba.}\tabularnewline
\toprule\noalign{}
Cadena & Longitud & Esperado & Obtenido & Tipo de caso \\
\midrule\noalign{}
\endfirsthead
\toprule\noalign{}
Cadena & Longitud & Esperado & Obtenido & Tipo de caso \\
\midrule\noalign{}
\endhead
\bottomrule\noalign{}
\endlastfoot
\texttt{01} & 2 & acepta & acepta & válida \\
\texttt{011} & 3 & acepta & acepta & válida \\
\texttt{10001} & 5 & acepta & acepta & válida \\
λ & 0 & rechaza & rechaza & inválida \\
\texttt{1} & 1 & rechaza & rechaza & inválida \\
\texttt{010} & 3 & rechaza & rechaza & inválida \\
\end{longtable}

\evidenciaaqui[0.95]{../evidencias/jflap/fichas/lista3-156.png}{lista3-156 en JFLAP: Multiple Run con las tres cadenas válidas y las
tres inválidas de la tabla anterior.}

\hypertarget{lista-3-ejercicio-157}{%
\subsubsection{Lista 3 · Ejercicio 157}\label{lista-3-ejercicio-157}}

\textbf{Enunciado.} Diagrama y simplificación del AFD de la tabla 157 (tabla transcrita del diagrama).

\textbf{Diseño.} Primero se eliminan los estados que no se alcanzan desde el inicial; después se refina la partición \{F, Q − F\} hasta que se estabiliza.

\textbf{Archivo.} \texttt{automatas/lista3/lista3-157.jff} · Σ = \{0, 1\} · \textbar Q\textbar{} = 6

\begin{longtable}[]{@{}lll@{}}
\caption{Tabla de transiciones (→ inicial, * aceptación, --- sin transición).}\tabularnewline
\toprule\noalign{}
δ & 0 & 1 \\
\midrule\noalign{}
\endfirsthead
\toprule\noalign{}
δ & 0 & 1 \\
\midrule\noalign{}
\endhead
\bottomrule\noalign{}
\endlastfoot
→ A & B & F \\
B & C & E \\
C & D & E \\
D & B & E \\
E & F & D \\
* F & F & B \\
\end{longtable}

\begin{longtable}[]{@{}lclll@{}}
\caption{Cadenas de prueba.}\tabularnewline
\toprule\noalign{}
Cadena & Longitud & Esperado & Obtenido & Tipo de caso \\
\midrule\noalign{}
\endfirsthead
\toprule\noalign{}
Cadena & Longitud & Esperado & Obtenido & Tipo de caso \\
\midrule\noalign{}
\endhead
\bottomrule\noalign{}
\endlastfoot
\texttt{1} & 1 & acepta & acepta & válida \\
\texttt{10} & 2 & acepta & acepta & válida \\
\texttt{0100} & 4 & acepta & acepta & válida \\
λ & 0 & rechaza & rechaza & inválida \\
\texttt{0} & 1 & rechaza & rechaza & inválida \\
\texttt{001} & 3 & rechaza & rechaza & inválida \\
\end{longtable}

\evidenciaaqui[0.95]{../evidencias/jflap/fichas/lista3-157.png}{lista3-157 en JFLAP: Multiple Run con las tres cadenas válidas y las
tres inválidas de la tabla anterior.}

\hypertarget{lista-3-ejercicio-158}{%
\subsubsection{Lista 3 · Ejercicio 158}\label{lista-3-ejercicio-158}}

\textbf{Enunciado.} Describa lo que hace el AFN del diagrama 158.

\textbf{Diseño.} El autómata acepta únicamente la cadena 01: L = \{01\}. Las transiciones que faltan llevan a un rechazo implícito.

\textbf{Archivo.} \texttt{automatas/lista3/lista3-158.jff} · Σ = \{0, 1\} · \textbar Q\textbar{} = 3

\begin{longtable}[]{@{}lll@{}}
\caption{Tabla de transiciones (→ inicial, * aceptación, --- sin transición).}\tabularnewline
\toprule\noalign{}
δ & 0 & 1 \\
\midrule\noalign{}
\endfirsthead
\toprule\noalign{}
δ & 0 & 1 \\
\midrule\noalign{}
\endhead
\bottomrule\noalign{}
\endlastfoot
→ q1 & q2 & --- \\
q2 & --- & q3 \\
* q3 & --- & --- \\
\end{longtable}

\begin{longtable}[]{@{}lclll@{}}
\caption{Cadenas de prueba.}\tabularnewline
\toprule\noalign{}
Cadena & Longitud & Esperado & Obtenido & Tipo de caso \\
\midrule\noalign{}
\endfirsthead
\toprule\noalign{}
Cadena & Longitud & Esperado & Obtenido & Tipo de caso \\
\midrule\noalign{}
\endhead
\bottomrule\noalign{}
\endlastfoot
\texttt{01} & 2 & acepta & acepta & válida \\
λ & 0 & rechaza & rechaza & inválida \\
\texttt{1} & 1 & rechaza & rechaza & inválida \\
\texttt{111} & 3 & rechaza & rechaza & inválida \\
\end{longtable}

\evidenciaaqui[0.95]{../evidencias/jflap/fichas/lista3-158.png}{lista3-158 en JFLAP: Multiple Run con las tres cadenas válidas y las
tres inválidas de la tabla anterior.}

\hypertarget{lista-3-ejercicio-159}{%
\subsubsection{Lista 3 · Ejercicio 159}\label{lista-3-ejercicio-159}}

\textbf{Enunciado.} Describa lo que hace el AFN del diagrama 159.

\textbf{Diseño.} El autómata acepta las cadenas formadas por unos seguidos de 01: L = 1*01. Las transiciones que faltan llevan a un rechazo implícito.

\textbf{Archivo.} \texttt{automatas/lista3/lista3-159.jff} · Σ = \{0, 1\} · \textbar Q\textbar{} = 3

\begin{longtable}[]{@{}lll@{}}
\caption{Tabla de transiciones (→ inicial, * aceptación, --- sin transición).}\tabularnewline
\toprule\noalign{}
δ & 0 & 1 \\
\midrule\noalign{}
\endfirsthead
\toprule\noalign{}
δ & 0 & 1 \\
\midrule\noalign{}
\endhead
\bottomrule\noalign{}
\endlastfoot
→ q1 & q2 & q1 \\
q2 & --- & q3 \\
* q3 & --- & --- \\
\end{longtable}

\begin{longtable}[]{@{}
  >{\raggedright\arraybackslash}p{(\columnwidth - 8\tabcolsep) * \real{0.2000}}
  >{\centering\arraybackslash}p{(\columnwidth - 8\tabcolsep) * \real{0.2000}}
  >{\raggedright\arraybackslash}p{(\columnwidth - 8\tabcolsep) * \real{0.2000}}
  >{\raggedright\arraybackslash}p{(\columnwidth - 8\tabcolsep) * \real{0.2000}}
  >{\raggedright\arraybackslash}p{(\columnwidth - 8\tabcolsep) * \real{0.2000}}@{}}
\caption{Cadenas de prueba.}\tabularnewline
\toprule\noalign{}
\begin{minipage}[b]{\linewidth}\raggedright
Cadena
\end{minipage} & \begin{minipage}[b]{\linewidth}\centering
Longitud
\end{minipage} & \begin{minipage}[b]{\linewidth}\raggedright
Esperado
\end{minipage} & \begin{minipage}[b]{\linewidth}\raggedright
Obtenido
\end{minipage} & \begin{minipage}[b]{\linewidth}\raggedright
Tipo de caso
\end{minipage} \\
\midrule\noalign{}
\endfirsthead
\toprule\noalign{}
\begin{minipage}[b]{\linewidth}\raggedright
Cadena
\end{minipage} & \begin{minipage}[b]{\linewidth}\centering
Longitud
\end{minipage} & \begin{minipage}[b]{\linewidth}\raggedright
Esperado
\end{minipage} & \begin{minipage}[b]{\linewidth}\raggedright
Obtenido
\end{minipage} & \begin{minipage}[b]{\linewidth}\raggedright
Tipo de caso
\end{minipage} \\
\midrule\noalign{}
\endhead
\bottomrule\noalign{}
\endlastfoot
\texttt{01} & 2 & acepta & acepta & válida \\
\texttt{101} & 3 & acepta & acepta & válida \\
\texttt{11101} & 5 & acepta & acepta & válida \\
λ & 0 & rechaza & rechaza & inválida \\
\texttt{0} & 1 & rechaza & rechaza & inválida \\
\texttt{001} & 3 & rechaza & rechaza & inválida \\
\texttt{1} & 1 & rechaza & rechaza & frontera: justo por debajo de la longitud mínima \\
\end{longtable}

\evidenciaaqui[0.95]{../evidencias/jflap/fichas/lista3-159.png}{lista3-159 en JFLAP: Multiple Run con las tres cadenas válidas y las
tres inválidas de la tabla anterior.}

\hypertarget{lista-3-ejercicio-160}{%
\subsubsection{Lista 3 · Ejercicio 160}\label{lista-3-ejercicio-160}}

\textbf{Enunciado.} Describa lo que hace el AFN del diagrama 160.

\textbf{Diseño.} El autómata acepta las cadenas que empiezan con 01: L = 01(0 + 1)*. Las transiciones que faltan llevan a un rechazo implícito.

\textbf{Archivo.} \texttt{automatas/lista3/lista3-160.jff} · Σ = \{0, 1\} · \textbar Q\textbar{} = 3

\begin{longtable}[]{@{}lll@{}}
\caption{Tabla de transiciones (→ inicial, * aceptación, --- sin transición).}\tabularnewline
\toprule\noalign{}
δ & 0 & 1 \\
\midrule\noalign{}
\endfirsthead
\toprule\noalign{}
δ & 0 & 1 \\
\midrule\noalign{}
\endhead
\bottomrule\noalign{}
\endlastfoot
→ q1 & q2 & --- \\
q2 & --- & q3 \\
* q3 & q3 & q3 \\
\end{longtable}

\begin{longtable}[]{@{}lclll@{}}
\caption{Cadenas de prueba.}\tabularnewline
\toprule\noalign{}
Cadena & Longitud & Esperado & Obtenido & Tipo de caso \\
\midrule\noalign{}
\endfirsthead
\toprule\noalign{}
Cadena & Longitud & Esperado & Obtenido & Tipo de caso \\
\midrule\noalign{}
\endhead
\bottomrule\noalign{}
\endlastfoot
\texttt{01} & 2 & acepta & acepta & válida \\
\texttt{011} & 3 & acepta & acepta & válida \\
\texttt{01100} & 5 & acepta & acepta & válida \\
λ & 0 & rechaza & rechaza & inválida \\
\texttt{1} & 1 & rechaza & rechaza & inválida \\
\texttt{101} & 3 & rechaza & rechaza & inválida \\
\end{longtable}

\evidenciaaqui[0.95]{../evidencias/jflap/fichas/lista3-160.png}{lista3-160 en JFLAP: Multiple Run con las tres cadenas válidas y las
tres inválidas de la tabla anterior.}

\hypertarget{lista-3-ejercicio-161}{%
\subsubsection{Lista 3 · Ejercicio 161}\label{lista-3-ejercicio-161}}

\textbf{Enunciado.} Describa lo que hace el AFN del diagrama 161.

\textbf{Diseño.} El autómata acepta las repeticiones del bloque 010, incluida λ: L = (010)*. Las transiciones que faltan llevan a un rechazo implícito.

\textbf{Archivo.} \texttt{automatas/lista3/lista3-161.jff} · Σ = \{0, 1\} · \textbar Q\textbar{} = 3

\begin{longtable}[]{@{}lll@{}}
\caption{Tabla de transiciones (→ inicial, * aceptación, --- sin transición).}\tabularnewline
\toprule\noalign{}
δ & 0 & 1 \\
\midrule\noalign{}
\endfirsthead
\toprule\noalign{}
δ & 0 & 1 \\
\midrule\noalign{}
\endhead
\bottomrule\noalign{}
\endlastfoot
→* q1 & q2 & --- \\
q2 & --- & q3 \\
q3 & q1 & --- \\
\end{longtable}

\begin{longtable}[]{@{}lclll@{}}
\caption{Cadenas de prueba.}\tabularnewline
\toprule\noalign{}
Cadena & Longitud & Esperado & Obtenido & Tipo de caso \\
\midrule\noalign{}
\endfirsthead
\toprule\noalign{}
Cadena & Longitud & Esperado & Obtenido & Tipo de caso \\
\midrule\noalign{}
\endhead
\bottomrule\noalign{}
\endlastfoot
λ & 0 & acepta & acepta & válida \\
\texttt{010} & 3 & acepta & acepta & válida \\
\texttt{010010} & 6 & acepta & acepta & válida \\
\texttt{0} & 1 & rechaza & rechaza & inválida \\
\texttt{01} & 2 & rechaza & rechaza & inválida \\
\texttt{0100} & 4 & rechaza & rechaza & inválida \\
\end{longtable}

\evidenciaaqui[0.95]{../evidencias/jflap/fichas/lista3-161.png}{lista3-161 en JFLAP: Multiple Run con las tres cadenas válidas y las
tres inválidas de la tabla anterior.}

\hypertarget{lista-3-afn-del-problema-097}{%
\subsubsection{Lista 3 · AFN del problema 097}\label{lista-3-afn-del-problema-097}}

\textbf{Enunciado.} AFN para el problema 97: Las palabras de Σ1 cuya longitud es un múltiplo de 3.

\textbf{Diseño.} No hay una versión no determinista más simple: el lenguaje depende de contar (longitud, apariciones o fases). Se presenta el AFD del problema sin el estado muerto; una función de transición parcial es un caso particular de AFN.

\textbf{Archivo.} \texttt{automatas/lista3/lista3-afn-097.jff} · Σ = \{0, 1, 2\} · \textbar Q\textbar{} = 3

\begin{longtable}[]{@{}llll@{}}
\caption{Tabla de transiciones (→ inicial, * aceptación, --- sin transición).}\tabularnewline
\toprule\noalign{}
δ & 0 & 1 & 2 \\
\midrule\noalign{}
\endfirsthead
\toprule\noalign{}
δ & 0 & 1 & 2 \\
\midrule\noalign{}
\endhead
\bottomrule\noalign{}
\endlastfoot
→* r0 & r1 & r1 & r1 \\
r1 & r2 & r2 & r2 \\
r2 & r0 & r0 & r0 \\
\end{longtable}

\begin{longtable}[]{@{}lclll@{}}
\caption{Cadenas de prueba.}\tabularnewline
\toprule\noalign{}
Cadena & Longitud & Esperado & Obtenido & Tipo de caso \\
\midrule\noalign{}
\endfirsthead
\toprule\noalign{}
Cadena & Longitud & Esperado & Obtenido & Tipo de caso \\
\midrule\noalign{}
\endhead
\bottomrule\noalign{}
\endlastfoot
λ & 0 & acepta & acepta & válida \\
\texttt{120} & 3 & acepta & acepta & válida \\
\texttt{200} & 3 & acepta & acepta & válida \\
\texttt{1} & 1 & rechaza & rechaza & inválida \\
\texttt{21} & 2 & rechaza & rechaza & inválida \\
\texttt{2102} & 4 & rechaza & rechaza & inválida \\
\end{longtable}

\evidenciaaqui[0.95]{../evidencias/jflap/fichas/lista3-afn-097.png}{lista3-afn-097 en JFLAP: Multiple Run con las tres cadenas válidas y las
tres inválidas de la tabla anterior.}

\hypertarget{lista-3-afn-del-problema-099}{%
\subsubsection{Lista 3 · AFN del problema 099}\label{lista-3-afn-del-problema-099}}

\textbf{Enunciado.} AFN para el problema 99: Las palabras de Σ1 que terminan con la cadena 001001.

\textbf{Diseño.} Ciclo en c0 y una cadena que adivina dónde empieza el sufijo 001001. El enunciado se resuelve de forma natural con un AFN de 7 estados; el AFD equivalente se obtuvo con la construcción de subconjuntos y la minimización y tiene 7 estados.

\textbf{Archivo.} \texttt{automatas/lista3/lista3-afn-099.jff} · Σ = \{0, 1, 2\} · \textbar Q\textbar{} = 7

\begin{longtable}[]{@{}lllll@{}}
\caption{Tabla de transiciones (→ inicial, * aceptación, --- sin transición).}\tabularnewline
\toprule\noalign{}
δ & 0 & 1 & 2 & Qué recuerda del prefijo leído \\
\midrule\noalign{}
\endfirsthead
\toprule\noalign{}
δ & 0 & 1 & 2 & Qué recuerda del prefijo leído \\
\midrule\noalign{}
\endhead
\bottomrule\noalign{}
\endlastfoot
→ q0 & q1 & q0 & q0 & activos del AFN: \{c0\} \\
q1 & q2 & q0 & q0 & activos del AFN: \{c0,c1\} \\
q2 & q2 & q3 & q0 & activos del AFN: \{c0,c1,c2\} \\
q3 & q4 & q0 & q0 & activos del AFN: \{c0,c3\} \\
q4 & q5 & q0 & q0 & activos del AFN: \{c0,c1,c4\} \\
q5 & q2 & q6 & q0 & activos del AFN: \{c0,c1,c2,c5\} \\
* q6 & q4 & q0 & q0 & activos del AFN: \{c0,c3,c6\} \\
\end{longtable}

\begin{longtable}[]{@{}
  >{\raggedright\arraybackslash}p{(\columnwidth - 8\tabcolsep) * \real{0.2000}}
  >{\centering\arraybackslash}p{(\columnwidth - 8\tabcolsep) * \real{0.2000}}
  >{\raggedright\arraybackslash}p{(\columnwidth - 8\tabcolsep) * \real{0.2000}}
  >{\raggedright\arraybackslash}p{(\columnwidth - 8\tabcolsep) * \real{0.2000}}
  >{\raggedright\arraybackslash}p{(\columnwidth - 8\tabcolsep) * \real{0.2000}}@{}}
\caption{Cadenas de prueba.}\tabularnewline
\toprule\noalign{}
\begin{minipage}[b]{\linewidth}\raggedright
Cadena
\end{minipage} & \begin{minipage}[b]{\linewidth}\centering
Longitud
\end{minipage} & \begin{minipage}[b]{\linewidth}\raggedright
Esperado
\end{minipage} & \begin{minipage}[b]{\linewidth}\raggedright
Obtenido
\end{minipage} & \begin{minipage}[b]{\linewidth}\raggedright
Tipo de caso
\end{minipage} \\
\midrule\noalign{}
\endfirsthead
\toprule\noalign{}
\begin{minipage}[b]{\linewidth}\raggedright
Cadena
\end{minipage} & \begin{minipage}[b]{\linewidth}\centering
Longitud
\end{minipage} & \begin{minipage}[b]{\linewidth}\raggedright
Esperado
\end{minipage} & \begin{minipage}[b]{\linewidth}\raggedright
Obtenido
\end{minipage} & \begin{minipage}[b]{\linewidth}\raggedright
Tipo de caso
\end{minipage} \\
\midrule\noalign{}
\endhead
\bottomrule\noalign{}
\endlastfoot
\texttt{001001} & 6 & acepta & acepta & válida \\
\texttt{1001001} & 7 & acepta & acepta & válida \\
\texttt{121001001} & 9 & acepta & acepta & válida \\
λ & 0 & rechaza & rechaza & inválida \\
\texttt{2} & 1 & rechaza & rechaza & inválida \\
\texttt{120} & 3 & rechaza & rechaza & inválida \\
\texttt{02210} & 5 & rechaza & rechaza & frontera: justo por debajo de la longitud mínima \\
\end{longtable}

\evidenciaaqui[0.95]{../evidencias/jflap/fichas/lista3-afn-099.png}{lista3-afn-099 en JFLAP: Multiple Run con las tres cadenas válidas y las
tres inválidas de la tabla anterior.}

\hypertarget{lista-3-afn-del-problema-102}{%
\subsubsection{Lista 3 · AFN del problema 102}\label{lista-3-afn-del-problema-102}}

\textbf{Enunciado.} AFN para el problema 102: Las palabras de Σ3 que contienen la cadena abbbab.

\textbf{Diseño.} El AFN adivina dónde empieza abbbab: ciclo inicial, el patrón y ciclo final. No necesita los retrocesos del autómata de Knuth-Morris-Pratt. El enunciado se resuelve de forma natural con un AFN de 7 estados; el AFD equivalente se obtuvo con la construcción de subconjuntos y la minimización y tiene 7 estados.

\textbf{Archivo.} \texttt{automatas/lista3/lista3-afn-102.jff} · Σ = \{a, b\} · \textbar Q\textbar{} = 7

\begin{longtable}[]{@{}
  >{\raggedright\arraybackslash}p{(\columnwidth - 6\tabcolsep) * \real{0.1579}}
  >{\raggedright\arraybackslash}p{(\columnwidth - 6\tabcolsep) * \real{0.1579}}
  >{\raggedright\arraybackslash}p{(\columnwidth - 6\tabcolsep) * \real{0.1579}}
  >{\raggedright\arraybackslash}p{(\columnwidth - 6\tabcolsep) * \real{0.5263}}@{}}
\caption{Tabla de transiciones (→ inicial, * aceptación, --- sin transición).}\tabularnewline
\toprule\noalign{}
\begin{minipage}[b]{\linewidth}\raggedright
δ
\end{minipage} & \begin{minipage}[b]{\linewidth}\raggedright
a
\end{minipage} & \begin{minipage}[b]{\linewidth}\raggedright
b
\end{minipage} & \begin{minipage}[b]{\linewidth}\raggedright
Qué recuerda del prefijo leído
\end{minipage} \\
\midrule\noalign{}
\endfirsthead
\toprule\noalign{}
\begin{minipage}[b]{\linewidth}\raggedright
δ
\end{minipage} & \begin{minipage}[b]{\linewidth}\raggedright
a
\end{minipage} & \begin{minipage}[b]{\linewidth}\raggedright
b
\end{minipage} & \begin{minipage}[b]{\linewidth}\raggedright
Qué recuerda del prefijo leído
\end{minipage} \\
\midrule\noalign{}
\endhead
\bottomrule\noalign{}
\endlastfoot
→ q0 & q1 & q0 & activos del AFN: \{c0\} \\
q1 & q1 & q2 & activos del AFN: \{c0,c1\} \\
q2 & q1 & q3 & activos del AFN: \{c0,c2\} \\
q3 & q1 & q4 & activos del AFN: \{c0,c3\} \\
q4 & q5 & q0 & activos del AFN: \{c0,c4\} \\
q5 & q1 & q6 & activos del AFN: \{c0,c1,c5\} \\
* q6 & q6 & q6 & activos del AFN: \{c0,c2,c6\} o \{c0,c1,c6\} o \{c0,c3,c6\} o \{c0,c4,c6\} o \{c0,c1,c5,c6\} o \{c0,c6\} \\
\end{longtable}

\begin{longtable}[]{@{}
  >{\raggedright\arraybackslash}p{(\columnwidth - 8\tabcolsep) * \real{0.2000}}
  >{\centering\arraybackslash}p{(\columnwidth - 8\tabcolsep) * \real{0.2000}}
  >{\raggedright\arraybackslash}p{(\columnwidth - 8\tabcolsep) * \real{0.2000}}
  >{\raggedright\arraybackslash}p{(\columnwidth - 8\tabcolsep) * \real{0.2000}}
  >{\raggedright\arraybackslash}p{(\columnwidth - 8\tabcolsep) * \real{0.2000}}@{}}
\caption{Cadenas de prueba.}\tabularnewline
\toprule\noalign{}
\begin{minipage}[b]{\linewidth}\raggedright
Cadena
\end{minipage} & \begin{minipage}[b]{\linewidth}\centering
Longitud
\end{minipage} & \begin{minipage}[b]{\linewidth}\raggedright
Esperado
\end{minipage} & \begin{minipage}[b]{\linewidth}\raggedright
Obtenido
\end{minipage} & \begin{minipage}[b]{\linewidth}\raggedright
Tipo de caso
\end{minipage} \\
\midrule\noalign{}
\endfirsthead
\toprule\noalign{}
\begin{minipage}[b]{\linewidth}\raggedright
Cadena
\end{minipage} & \begin{minipage}[b]{\linewidth}\centering
Longitud
\end{minipage} & \begin{minipage}[b]{\linewidth}\raggedright
Esperado
\end{minipage} & \begin{minipage}[b]{\linewidth}\raggedright
Obtenido
\end{minipage} & \begin{minipage}[b]{\linewidth}\raggedright
Tipo de caso
\end{minipage} \\
\midrule\noalign{}
\endhead
\bottomrule\noalign{}
\endlastfoot
\texttt{abbbab} & 6 & acepta & acepta & válida \\
\texttt{abbbabb} & 7 & acepta & acepta & válida \\
\texttt{abbbabaaa} & 9 & acepta & acepta & válida \\
λ & 0 & rechaza & rechaza & inválida \\
\texttt{a} & 1 & rechaza & rechaza & inválida \\
\texttt{baa} & 3 & rechaza & rechaza & inválida \\
\texttt{aabab} & 5 & rechaza & rechaza & frontera: justo por debajo de la longitud mínima \\
\end{longtable}

\evidenciaaqui[0.95]{../evidencias/jflap/fichas/lista3-afn-102.png}{lista3-afn-102 en JFLAP: Multiple Run con las tres cadenas válidas y las
tres inválidas de la tabla anterior.}

\hypertarget{lista-3-afn-del-problema-103}{%
\subsubsection{Lista 3 · AFN del problema 103}\label{lista-3-afn-del-problema-103}}

\textbf{Enunciado.} AFN para el problema 103: Las palabras de Σ3 que contienen la cadena aabbba.

\textbf{Diseño.} El AFN adivina dónde empieza aabbba: ciclo inicial, el patrón y ciclo final. No necesita los retrocesos del autómata de Knuth-Morris-Pratt. El enunciado se resuelve de forma natural con un AFN de 7 estados; el AFD equivalente se obtuvo con la construcción de subconjuntos y la minimización y tiene 7 estados.

\textbf{Archivo.} \texttt{automatas/lista3/lista3-afn-103.jff} · Σ = \{a, b\} · \textbar Q\textbar{} = 7

\begin{longtable}[]{@{}
  >{\raggedright\arraybackslash}p{(\columnwidth - 6\tabcolsep) * \real{0.1579}}
  >{\raggedright\arraybackslash}p{(\columnwidth - 6\tabcolsep) * \real{0.1579}}
  >{\raggedright\arraybackslash}p{(\columnwidth - 6\tabcolsep) * \real{0.1579}}
  >{\raggedright\arraybackslash}p{(\columnwidth - 6\tabcolsep) * \real{0.5263}}@{}}
\caption{Tabla de transiciones (→ inicial, * aceptación, --- sin transición).}\tabularnewline
\toprule\noalign{}
\begin{minipage}[b]{\linewidth}\raggedright
δ
\end{minipage} & \begin{minipage}[b]{\linewidth}\raggedright
a
\end{minipage} & \begin{minipage}[b]{\linewidth}\raggedright
b
\end{minipage} & \begin{minipage}[b]{\linewidth}\raggedright
Qué recuerda del prefijo leído
\end{minipage} \\
\midrule\noalign{}
\endfirsthead
\toprule\noalign{}
\begin{minipage}[b]{\linewidth}\raggedright
δ
\end{minipage} & \begin{minipage}[b]{\linewidth}\raggedright
a
\end{minipage} & \begin{minipage}[b]{\linewidth}\raggedright
b
\end{minipage} & \begin{minipage}[b]{\linewidth}\raggedright
Qué recuerda del prefijo leído
\end{minipage} \\
\midrule\noalign{}
\endhead
\bottomrule\noalign{}
\endlastfoot
→ q0 & q1 & q0 & activos del AFN: \{c0\} \\
q1 & q2 & q0 & activos del AFN: \{c0,c1\} \\
q2 & q2 & q3 & activos del AFN: \{c0,c1,c2\} \\
q3 & q1 & q4 & activos del AFN: \{c0,c3\} \\
q4 & q1 & q5 & activos del AFN: \{c0,c4\} \\
q5 & q6 & q0 & activos del AFN: \{c0,c5\} \\
* q6 & q6 & q6 & activos del AFN: \{c0,c1,c6\} o \{c0,c1,c2,c6\} o \{c0,c6\} o \{c0,c3,c6\} o \{c0,c4,c6\} o \{c0,c5,c6\} \\
\end{longtable}

\begin{longtable}[]{@{}
  >{\raggedright\arraybackslash}p{(\columnwidth - 8\tabcolsep) * \real{0.2000}}
  >{\centering\arraybackslash}p{(\columnwidth - 8\tabcolsep) * \real{0.2000}}
  >{\raggedright\arraybackslash}p{(\columnwidth - 8\tabcolsep) * \real{0.2000}}
  >{\raggedright\arraybackslash}p{(\columnwidth - 8\tabcolsep) * \real{0.2000}}
  >{\raggedright\arraybackslash}p{(\columnwidth - 8\tabcolsep) * \real{0.2000}}@{}}
\caption{Cadenas de prueba.}\tabularnewline
\toprule\noalign{}
\begin{minipage}[b]{\linewidth}\raggedright
Cadena
\end{minipage} & \begin{minipage}[b]{\linewidth}\centering
Longitud
\end{minipage} & \begin{minipage}[b]{\linewidth}\raggedright
Esperado
\end{minipage} & \begin{minipage}[b]{\linewidth}\raggedright
Obtenido
\end{minipage} & \begin{minipage}[b]{\linewidth}\raggedright
Tipo de caso
\end{minipage} \\
\midrule\noalign{}
\endfirsthead
\toprule\noalign{}
\begin{minipage}[b]{\linewidth}\raggedright
Cadena
\end{minipage} & \begin{minipage}[b]{\linewidth}\centering
Longitud
\end{minipage} & \begin{minipage}[b]{\linewidth}\raggedright
Esperado
\end{minipage} & \begin{minipage}[b]{\linewidth}\raggedright
Obtenido
\end{minipage} & \begin{minipage}[b]{\linewidth}\raggedright
Tipo de caso
\end{minipage} \\
\midrule\noalign{}
\endhead
\bottomrule\noalign{}
\endlastfoot
\texttt{aabbba} & 6 & acepta & acepta & válida \\
\texttt{baabbba} & 7 & acepta & acepta & válida \\
\texttt{babaabbba} & 9 & acepta & acepta & válida \\
λ & 0 & rechaza & rechaza & inválida \\
\texttt{a} & 1 & rechaza & rechaza & inválida \\
\texttt{aaa} & 3 & rechaza & rechaza & inválida \\
\texttt{ababb} & 5 & rechaza & rechaza & frontera: justo por debajo de la longitud mínima \\
\end{longtable}

\evidenciaaqui[0.95]{../evidencias/jflap/fichas/lista3-afn-103.png}{lista3-afn-103 en JFLAP: Multiple Run con las tres cadenas válidas y las
tres inválidas de la tabla anterior.}

\hypertarget{lista-3-afn-del-problema-106}{%
\subsubsection{Lista 3 · AFN del problema 106}\label{lista-3-afn-del-problema-106}}

\textbf{Enunciado.} AFN para el problema 106: Las palabras de Σ1 que incluyen la subcadena 1101001.

\textbf{Diseño.} El AFN adivina dónde empieza 1101001: ciclo inicial, el patrón y ciclo final. No necesita los retrocesos del autómata de Knuth-Morris-Pratt. El enunciado se resuelve de forma natural con un AFN de 8 estados; el AFD equivalente se obtuvo con la construcción de subconjuntos y la minimización y tiene 8 estados.

\textbf{Archivo.} \texttt{automatas/lista3/lista3-afn-106.jff} · Σ = \{0, 1, 2\} · \textbar Q\textbar{} = 8

\begin{longtable}[]{@{}
  >{\raggedright\arraybackslash}p{(\columnwidth - 8\tabcolsep) * \real{0.1364}}
  >{\raggedright\arraybackslash}p{(\columnwidth - 8\tabcolsep) * \real{0.1364}}
  >{\raggedright\arraybackslash}p{(\columnwidth - 8\tabcolsep) * \real{0.1364}}
  >{\raggedright\arraybackslash}p{(\columnwidth - 8\tabcolsep) * \real{0.1364}}
  >{\raggedright\arraybackslash}p{(\columnwidth - 8\tabcolsep) * \real{0.4545}}@{}}
\caption{Tabla de transiciones (→ inicial, * aceptación, --- sin transición).}\tabularnewline
\toprule\noalign{}
\begin{minipage}[b]{\linewidth}\raggedright
δ
\end{minipage} & \begin{minipage}[b]{\linewidth}\raggedright
0
\end{minipage} & \begin{minipage}[b]{\linewidth}\raggedright
1
\end{minipage} & \begin{minipage}[b]{\linewidth}\raggedright
2
\end{minipage} & \begin{minipage}[b]{\linewidth}\raggedright
Qué recuerda del prefijo leído
\end{minipage} \\
\midrule\noalign{}
\endfirsthead
\toprule\noalign{}
\begin{minipage}[b]{\linewidth}\raggedright
δ
\end{minipage} & \begin{minipage}[b]{\linewidth}\raggedright
0
\end{minipage} & \begin{minipage}[b]{\linewidth}\raggedright
1
\end{minipage} & \begin{minipage}[b]{\linewidth}\raggedright
2
\end{minipage} & \begin{minipage}[b]{\linewidth}\raggedright
Qué recuerda del prefijo leído
\end{minipage} \\
\midrule\noalign{}
\endhead
\bottomrule\noalign{}
\endlastfoot
→ q0 & q0 & q1 & q0 & activos del AFN: \{c0\} \\
q1 & q0 & q2 & q0 & activos del AFN: \{c0,c1\} \\
q2 & q3 & q2 & q0 & activos del AFN: \{c0,c1,c2\} \\
q3 & q0 & q4 & q0 & activos del AFN: \{c0,c3\} \\
q4 & q5 & q2 & q0 & activos del AFN: \{c0,c1,c4\} \\
q5 & q6 & q1 & q0 & activos del AFN: \{c0,c5\} \\
q6 & q0 & q7 & q0 & activos del AFN: \{c0,c6\} \\
* q7 & q7 & q7 & q7 & activos del AFN: \{c0,c1,c7\} o \{c0,c7\} o \{c0,c1,c2,c7\} o \{c0,c3,c7\} o \{c0,c1,c4,c7\} o \{c0,c5,c7\} o \{c0,c6,c7\} \\
\end{longtable}

\begin{longtable}[]{@{}
  >{\raggedright\arraybackslash}p{(\columnwidth - 8\tabcolsep) * \real{0.2000}}
  >{\centering\arraybackslash}p{(\columnwidth - 8\tabcolsep) * \real{0.2000}}
  >{\raggedright\arraybackslash}p{(\columnwidth - 8\tabcolsep) * \real{0.2000}}
  >{\raggedright\arraybackslash}p{(\columnwidth - 8\tabcolsep) * \real{0.2000}}
  >{\raggedright\arraybackslash}p{(\columnwidth - 8\tabcolsep) * \real{0.2000}}@{}}
\caption{Cadenas de prueba.}\tabularnewline
\toprule\noalign{}
\begin{minipage}[b]{\linewidth}\raggedright
Cadena
\end{minipage} & \begin{minipage}[b]{\linewidth}\centering
Longitud
\end{minipage} & \begin{minipage}[b]{\linewidth}\raggedright
Esperado
\end{minipage} & \begin{minipage}[b]{\linewidth}\raggedright
Obtenido
\end{minipage} & \begin{minipage}[b]{\linewidth}\raggedright
Tipo de caso
\end{minipage} \\
\midrule\noalign{}
\endfirsthead
\toprule\noalign{}
\begin{minipage}[b]{\linewidth}\raggedright
Cadena
\end{minipage} & \begin{minipage}[b]{\linewidth}\centering
Longitud
\end{minipage} & \begin{minipage}[b]{\linewidth}\raggedright
Esperado
\end{minipage} & \begin{minipage}[b]{\linewidth}\raggedright
Obtenido
\end{minipage} & \begin{minipage}[b]{\linewidth}\raggedright
Tipo de caso
\end{minipage} \\
\midrule\noalign{}
\endhead
\bottomrule\noalign{}
\endlastfoot
\texttt{1101001} & 7 & acepta & acepta & válida \\
\texttt{11010010} & 8 & acepta & acepta & válida \\
\texttt{1101001110} & 10 & acepta & acepta & válida \\
λ & 0 & rechaza & rechaza & inválida \\
\texttt{2} & 1 & rechaza & rechaza & inválida \\
\texttt{210} & 3 & rechaza & rechaza & inválida \\
\texttt{221200} & 6 & rechaza & rechaza & frontera: justo por debajo de la longitud mínima \\
\end{longtable}

\evidenciaaqui[0.95]{../evidencias/jflap/fichas/lista3-afn-106.png}{lista3-afn-106 en JFLAP: Multiple Run con las tres cadenas válidas y las
tres inválidas de la tabla anterior.}

\hypertarget{lista-3-afn-del-problema-107}{%
\subsubsection{Lista 3 · AFN del problema 107}\label{lista-3-afn-del-problema-107}}

\textbf{Enunciado.} AFN para el problema 107: Las palabras de Σ1 que incluyen al menos una vez 0101 y al menos una vez 021.

\textbf{Diseño.} Intersección de los AFN de «contiene 0101» y «contiene 021» con la construcción de producto; cada estado es un par cᵢcⱼ. Como las dos apariciones pueden traslaparse, no basta con concatenar los patrones. El enunciado se resuelve de forma natural con un AFN de 16 estados; el AFD equivalente se obtuvo con la construcción de subconjuntos y la minimización y tiene 13 estados.

\textbf{Archivo.} \texttt{automatas/lista3/lista3-afn-107.jff} · Σ = \{0, 1, 2\} · \textbar Q\textbar{} = 13

\begin{longtable}[]{@{}
  >{\raggedright\arraybackslash}p{(\columnwidth - 8\tabcolsep) * \real{0.1364}}
  >{\raggedright\arraybackslash}p{(\columnwidth - 8\tabcolsep) * \real{0.1364}}
  >{\raggedright\arraybackslash}p{(\columnwidth - 8\tabcolsep) * \real{0.1364}}
  >{\raggedright\arraybackslash}p{(\columnwidth - 8\tabcolsep) * \real{0.1364}}
  >{\raggedright\arraybackslash}p{(\columnwidth - 8\tabcolsep) * \real{0.4545}}@{}}
\caption{Tabla de transiciones (→ inicial, * aceptación, --- sin transición).}\tabularnewline
\toprule\noalign{}
\begin{minipage}[b]{\linewidth}\raggedright
δ
\end{minipage} & \begin{minipage}[b]{\linewidth}\raggedright
0
\end{minipage} & \begin{minipage}[b]{\linewidth}\raggedright
1
\end{minipage} & \begin{minipage}[b]{\linewidth}\raggedright
2
\end{minipage} & \begin{minipage}[b]{\linewidth}\raggedright
Qué recuerda del prefijo leído
\end{minipage} \\
\midrule\noalign{}
\endfirsthead
\toprule\noalign{}
\begin{minipage}[b]{\linewidth}\raggedright
δ
\end{minipage} & \begin{minipage}[b]{\linewidth}\raggedright
0
\end{minipage} & \begin{minipage}[b]{\linewidth}\raggedright
1
\end{minipage} & \begin{minipage}[b]{\linewidth}\raggedright
2
\end{minipage} & \begin{minipage}[b]{\linewidth}\raggedright
Qué recuerda del prefijo leído
\end{minipage} \\
\midrule\noalign{}
\endhead
\bottomrule\noalign{}
\endlastfoot
→ q0 & q1 & q0 & q0 & activos del AFN: \{c0c0\} \\
q1 & q1 & q2 & q3 & activos del AFN: \{c0c0,c0c1,c1c0,c1c1\} \\
q2 & q4 & q0 & q0 & activos del AFN: \{c0c0,c2c0\} \\
q3 & q1 & q5 & q0 & activos del AFN: \{c0c0,c0c2\} \\
q4 & q1 & q6 & q3 & activos del AFN: \{c0c0,c0c1,c1c0,c1c1,c3c0,c3c1\} \\
q5 & q7 & q5 & q5 & activos del AFN: \{c0c0,c0c3\} o \{c0c0,c0c2,c0c3\} \\
q6 & q8 & q6 & q6 & activos del AFN: \{c0c0,c2c0,c4c0\} o \{c0c0,c4c0\} \\
q7 & q7 & q9 & q5 & activos del AFN: \{c0c0,c0c1,c1c0,c1c1,c0c3,c1c3\} \\
q8 & q8 & q6 & q10 & activos del AFN: \{c0c0,c0c1,c1c0,c1c1,c3c0,c3c1,c4c0,c4c1\} o \{c0c0,c0c1,c1c0,c1c1,c4c0,c4c1\} \\
q9 & q11 & q5 & q5 & activos del AFN: \{c0c0,c2c0,c0c3,c2c3\} \\
q10 & q8 & q12 & q6 & activos del AFN: \{c0c0,c0c2,c4c0,c4c2\} \\
q11 & q7 & q12 & q5 & activos del AFN: \{c0c0,c0c1,c1c0,c1c1,c0c3,c3c0,c3c1,c1c3,c3c3\} \\
* q12 & q12 & q12 & q12 & activos del AFN: \{c0c0,c0c3,c4c0,c4c3\} o \{c0c0,c2c0,c0c3,c4c0,c2c3,c4c3\} o \{c0c0,c0c1,c1c0,c1c1,c0c3,c1c3,c4c0,c4c1,c4c3\} o \{c0c0,c0c1,c1c0,c1c1,c0c3,c3c0,c3c1,c1c3,c4c0,c4c1,c3c3,c4c3\} o \{c0c0,c0c2,c0c3,c4c0,c4c2,c4c3\} \\
\end{longtable}

\begin{longtable}[]{@{}
  >{\raggedright\arraybackslash}p{(\columnwidth - 8\tabcolsep) * \real{0.2000}}
  >{\centering\arraybackslash}p{(\columnwidth - 8\tabcolsep) * \real{0.2000}}
  >{\raggedright\arraybackslash}p{(\columnwidth - 8\tabcolsep) * \real{0.2000}}
  >{\raggedright\arraybackslash}p{(\columnwidth - 8\tabcolsep) * \real{0.2000}}
  >{\raggedright\arraybackslash}p{(\columnwidth - 8\tabcolsep) * \real{0.2000}}@{}}
\caption{Cadenas de prueba.}\tabularnewline
\toprule\noalign{}
\begin{minipage}[b]{\linewidth}\raggedright
Cadena
\end{minipage} & \begin{minipage}[b]{\linewidth}\centering
Longitud
\end{minipage} & \begin{minipage}[b]{\linewidth}\raggedright
Esperado
\end{minipage} & \begin{minipage}[b]{\linewidth}\raggedright
Obtenido
\end{minipage} & \begin{minipage}[b]{\linewidth}\raggedright
Tipo de caso
\end{minipage} \\
\midrule\noalign{}
\endfirsthead
\toprule\noalign{}
\begin{minipage}[b]{\linewidth}\raggedright
Cadena
\end{minipage} & \begin{minipage}[b]{\linewidth}\centering
Longitud
\end{minipage} & \begin{minipage}[b]{\linewidth}\raggedright
Esperado
\end{minipage} & \begin{minipage}[b]{\linewidth}\raggedright
Obtenido
\end{minipage} & \begin{minipage}[b]{\linewidth}\raggedright
Tipo de caso
\end{minipage} \\
\midrule\noalign{}
\endhead
\bottomrule\noalign{}
\endlastfoot
\texttt{0101021} & 7 & acepta & acepta & válida \\
\texttt{20210101} & 8 & acepta & acepta & válida \\
\texttt{0212100101} & 10 & acepta & acepta & válida \\
λ & 0 & rechaza & rechaza & inválida \\
\texttt{0} & 1 & rechaza & rechaza & inválida \\
\texttt{201} & 3 & rechaza & rechaza & inválida \\
\texttt{0210101} & 7 & acepta & acepta & frontera: aceptada más corta \\
\texttt{101002} & 6 & rechaza & rechaza & frontera: justo por debajo de la longitud mínima \\
\end{longtable}

\evidenciaaqui[0.95]{../evidencias/jflap/fichas/lista3-afn-107.png}{lista3-afn-107 en JFLAP: Multiple Run con las tres cadenas válidas y las
tres inválidas de la tabla anterior.}

\hypertarget{lista-3-afn-del-problema-109}{%
\subsubsection{Lista 3 · AFN del problema 109}\label{lista-3-afn-del-problema-109}}

\textbf{Enunciado.} AFN para el problema 109: Las palabras de Σ2 que incluyen abbcb y cbaabb, que sí se pueden traslapar.

\textbf{Diseño.} Intersección de los dos AFN «contiene» por la construcción de producto, que admite traslapes. El enunciado se resuelve de forma natural con un AFN de 26 estados; el AFD equivalente se obtuvo con la construcción de subconjuntos y la minimización y tiene 22 estados.

\textbf{Archivo.} \texttt{automatas/lista3/lista3-afn-109.jff} · Σ = \{0, 1, a, b, c\} · \textbar Q\textbar{} = 22

\begin{longtable}[]{@{}
  >{\raggedright\arraybackslash}p{(\columnwidth - 12\tabcolsep) * \real{0.1071}}
  >{\raggedright\arraybackslash}p{(\columnwidth - 12\tabcolsep) * \real{0.1071}}
  >{\raggedright\arraybackslash}p{(\columnwidth - 12\tabcolsep) * \real{0.1071}}
  >{\raggedright\arraybackslash}p{(\columnwidth - 12\tabcolsep) * \real{0.1071}}
  >{\raggedright\arraybackslash}p{(\columnwidth - 12\tabcolsep) * \real{0.1071}}
  >{\raggedright\arraybackslash}p{(\columnwidth - 12\tabcolsep) * \real{0.1071}}
  >{\raggedright\arraybackslash}p{(\columnwidth - 12\tabcolsep) * \real{0.3571}}@{}}
\caption{Tabla de transiciones (→ inicial, * aceptación, --- sin transición).}\tabularnewline
\toprule\noalign{}
\begin{minipage}[b]{\linewidth}\raggedright
δ
\end{minipage} & \begin{minipage}[b]{\linewidth}\raggedright
0
\end{minipage} & \begin{minipage}[b]{\linewidth}\raggedright
1
\end{minipage} & \begin{minipage}[b]{\linewidth}\raggedright
a
\end{minipage} & \begin{minipage}[b]{\linewidth}\raggedright
b
\end{minipage} & \begin{minipage}[b]{\linewidth}\raggedright
c
\end{minipage} & \begin{minipage}[b]{\linewidth}\raggedright
Qué recuerda del prefijo leído
\end{minipage} \\
\midrule\noalign{}
\endfirsthead
\toprule\noalign{}
\begin{minipage}[b]{\linewidth}\raggedright
δ
\end{minipage} & \begin{minipage}[b]{\linewidth}\raggedright
0
\end{minipage} & \begin{minipage}[b]{\linewidth}\raggedright
1
\end{minipage} & \begin{minipage}[b]{\linewidth}\raggedright
a
\end{minipage} & \begin{minipage}[b]{\linewidth}\raggedright
b
\end{minipage} & \begin{minipage}[b]{\linewidth}\raggedright
c
\end{minipage} & \begin{minipage}[b]{\linewidth}\raggedright
Qué recuerda del prefijo leído
\end{minipage} \\
\midrule\noalign{}
\endhead
\bottomrule\noalign{}
\endlastfoot
→ q0 & q0 & q0 & q1 & q0 & q2 & activos del AFN: \{c0c0\} \\
q1 & q0 & q0 & q1 & q3 & q2 & activos del AFN: \{c0c0,c1c0\} \\
q2 & q0 & q0 & q1 & q4 & q2 & activos del AFN: \{c0c0,c0c1\} \\
q3 & q0 & q0 & q1 & q5 & q2 & activos del AFN: \{c0c0,c2c0\} \\
q4 & q0 & q0 & q6 & q0 & q2 & activos del AFN: \{c0c0,c0c2\} \\
q5 & q0 & q0 & q1 & q0 & q7 & activos del AFN: \{c0c0,c3c0\} \\
q6 & q0 & q0 & q8 & q3 & q2 & activos del AFN: \{c0c0,c1c0,c0c3,c1c3\} \\
q7 & q0 & q0 & q1 & q9 & q2 & activos del AFN: \{c0c0,c0c1,c4c0,c4c1\} \\
q8 & q0 & q0 & q1 & q10 & q2 & activos del AFN: \{c0c0,c1c0,c0c4,c1c4\} \\
q9 & q11 & q11 & q12 & q11 & q13 & activos del AFN: \{c0c0,c0c2,c5c0,c5c2\} \\
q10 & q0 & q0 & q1 & q14 & q2 & activos del AFN: \{c0c0,c2c0,c0c5,c2c5\} \\
q11 & q11 & q11 & q11 & q11 & q13 & activos del AFN: \{c0c0,c5c0\} o \{c0c0,c1c0,c5c0\} o \{c0c0,c2c0,c5c0\} o \{c0c0,c3c0,c5c0\} \\
q12 & q11 & q11 & q15 & q11 & q13 & activos del AFN: \{c0c0,c1c0,c0c3,c1c3,c5c0,c5c3\} \\
q13 & q11 & q11 & q11 & q9 & q13 & activos del AFN: \{c0c0,c0c1,c5c0,c5c1\} o \{c0c0,c0c1,c4c0,c4c1,c5c0,c5c1\} \\
q14 & q16 & q16 & q17 & q16 & q18 & activos del AFN: \{c0c0,c3c0,c0c6,c3c6\} \\
q15 & q11 & q11 & q11 & q19 & q13 & activos del AFN: \{c0c0,c1c0,c0c4,c1c4,c5c0,c5c4\} \\
q16 & q16 & q16 & q17 & q16 & q16 & activos del AFN: \{c0c0,c0c6\} o \{c0c0,c0c1,c0c6\} o \{c0c0,c0c2,c0c6\} \\
q17 & q16 & q16 & q17 & q20 & q16 & activos del AFN: \{c0c0,c1c0,c0c6,c1c6\} o \{c0c0,c1c0,c0c3,c1c3,c0c6,c1c6\} o \{c0c0,c1c0,c0c4,c1c4,c0c6,c1c6\} \\
q18 & q16 & q16 & q17 & q21 & q16 & activos del AFN: \{c0c0,c0c1,c4c0,c4c1,c0c6,c4c6\} \\
q19 & q11 & q11 & q11 & q21 & q13 & activos del AFN: \{c0c0,c2c0,c5c0,c0c5,c2c5,c5c5\} \\
q20 & q16 & q16 & q17 & q14 & q16 & activos del AFN: \{c0c0,c2c0,c0c6,c2c6\} o \{c0c0,c2c0,c0c5,c2c5,c0c6,c2c6\} \\
* q21 & q21 & q21 & q21 & q21 & q21 & activos del AFN: \{c0c0,c0c2,c5c0,c5c2,c0c6,c5c6\} o \{c0c0,c3c0,c5c0,c0c6,c3c6,c5c6\} o \{c0c0,c5c0,c0c6,c5c6\} o \{c0c0,c1c0,c0c3,c1c3,c5c0,c5c3,c0c6,c1c6,c5c6\} o \{c0c0,c0c1,c5c0,c5c1,c0c6,c5c6\} o \{c0c0,c1c0,c5c0,c0c6,c1c6,c5c6\} o \{c0c0,c0c1,c4c0,c4c1,c5c0,c5c1,c0c6,c4c6,c5c6\} o \{c0c0,c1c0,c0c4,c1c4,c5c0,c0c6,c5c4,c1c6,c5c6\} o \{c0c0,c2c0,c5c0,c0c6,c2c6,c5c6\} o \{c0c0,c2c0,c5c0,c0c5,c2c5,c0c6,c5c5,c2c6,c5c6\} \\
\end{longtable}

\begin{longtable}[]{@{}
  >{\raggedright\arraybackslash}p{(\columnwidth - 8\tabcolsep) * \real{0.2000}}
  >{\centering\arraybackslash}p{(\columnwidth - 8\tabcolsep) * \real{0.2000}}
  >{\raggedright\arraybackslash}p{(\columnwidth - 8\tabcolsep) * \real{0.2000}}
  >{\raggedright\arraybackslash}p{(\columnwidth - 8\tabcolsep) * \real{0.2000}}
  >{\raggedright\arraybackslash}p{(\columnwidth - 8\tabcolsep) * \real{0.2000}}@{}}
\caption{Cadenas de prueba.}\tabularnewline
\toprule\noalign{}
\begin{minipage}[b]{\linewidth}\raggedright
Cadena
\end{minipage} & \begin{minipage}[b]{\linewidth}\centering
Longitud
\end{minipage} & \begin{minipage}[b]{\linewidth}\raggedright
Esperado
\end{minipage} & \begin{minipage}[b]{\linewidth}\raggedright
Obtenido
\end{minipage} & \begin{minipage}[b]{\linewidth}\raggedright
Tipo de caso
\end{minipage} \\
\midrule\noalign{}
\endfirsthead
\toprule\noalign{}
\begin{minipage}[b]{\linewidth}\raggedright
Cadena
\end{minipage} & \begin{minipage}[b]{\linewidth}\centering
Longitud
\end{minipage} & \begin{minipage}[b]{\linewidth}\raggedright
Esperado
\end{minipage} & \begin{minipage}[b]{\linewidth}\raggedright
Obtenido
\end{minipage} & \begin{minipage}[b]{\linewidth}\raggedright
Tipo de caso
\end{minipage} \\
\midrule\noalign{}
\endhead
\bottomrule\noalign{}
\endlastfoot
\texttt{cbaabbcb} & 8 & acepta & acepta & válida \\
\texttt{acbaabbcb} & 9 & acepta & acepta & válida \\
\texttt{1acbaabbcba} & 11 & acepta & acepta & válida \\
λ & 0 & rechaza & rechaza & inválida \\
\texttt{a} & 1 & rechaza & rechaza & inválida \\
\texttt{1ac} & 3 & rechaza & rechaza & inválida \\
\texttt{abc0bca} & 7 & rechaza & rechaza & frontera: justo por debajo de la longitud mínima \\
\end{longtable}

\evidenciaaqui[0.95]{../evidencias/jflap/fichas/lista3-afn-109.png}{lista3-afn-109 en JFLAP: Multiple Run con las tres cadenas válidas y las
tres inválidas de la tabla anterior.}

\hypertarget{lista-3-afn-del-problema-111}{%
\subsubsection{Lista 3 · AFN del problema 111}\label{lista-3-afn-del-problema-111}}

\textbf{Enunciado.} AFN para el problema 111: Las palabras de Σ1 que contienen exactamente tres veces la subcadena 01002.

\textbf{Diseño.} No hay una versión no determinista más simple: el lenguaje depende de contar (longitud, apariciones o fases). Se presenta el AFD del problema sin el estado muerto; una función de transición parcial es un caso particular de AFN.

\textbf{Archivo.} \texttt{automatas/lista3/lista3-afn-111.jff} · Σ = \{0, 1, 2\} · \textbar Q\textbar{} = 20

\begin{longtable}[]{@{}llll@{}}
\caption{Tabla de transiciones (→ inicial, * aceptación, --- sin transición).}\tabularnewline
\toprule\noalign{}
δ & 0 & 1 & 2 \\
\midrule\noalign{}
\endfirsthead
\toprule\noalign{}
δ & 0 & 1 & 2 \\
\midrule\noalign{}
\endhead
\bottomrule\noalign{}
\endlastfoot
→ 0·λ & 0·0 & 0·λ & 0·λ \\
0·0 & 0·0 & 0·01 & 0·λ \\
0·01 & 0·010 & 0·λ & 0·λ \\
0·010 & 0·0100 & 0·01 & 0·λ \\
0·0100 & 0·0 & 0·01 & 1·λ \\
1·λ & 1·0 & 1·λ & 1·λ \\
1·0 & 1·0 & 1·01 & 1·λ \\
1·01 & 1·010 & 1·λ & 1·λ \\
1·010 & 1·0100 & 1·01 & 1·λ \\
1·0100 & 1·0 & 1·01 & 2·λ \\
2·λ & 2·0 & 2·λ & 2·λ \\
2·0 & 2·0 & 2·01 & 2·λ \\
2·01 & 2·010 & 2·λ & 2·λ \\
2·010 & 2·0100 & 2·01 & 2·λ \\
2·0100 & 2·0 & 2·01 & 3·λ \\
* 3·λ & 3·0 & 3·λ & 3·λ \\
* 3·0 & 3·0 & 3·01 & 3·λ \\
* 3·01 & 3·010 & 3·λ & 3·λ \\
* 3·010 & 3·0100 & 3·01 & 3·λ \\
* 3·0100 & 3·0 & 3·01 & --- \\
\end{longtable}

\begin{longtable}[]{@{}
  >{\raggedright\arraybackslash}p{(\columnwidth - 8\tabcolsep) * \real{0.2000}}
  >{\centering\arraybackslash}p{(\columnwidth - 8\tabcolsep) * \real{0.2000}}
  >{\raggedright\arraybackslash}p{(\columnwidth - 8\tabcolsep) * \real{0.2000}}
  >{\raggedright\arraybackslash}p{(\columnwidth - 8\tabcolsep) * \real{0.2000}}
  >{\raggedright\arraybackslash}p{(\columnwidth - 8\tabcolsep) * \real{0.2000}}@{}}
\caption{Cadenas de prueba.}\tabularnewline
\toprule\noalign{}
\begin{minipage}[b]{\linewidth}\raggedright
Cadena
\end{minipage} & \begin{minipage}[b]{\linewidth}\centering
Longitud
\end{minipage} & \begin{minipage}[b]{\linewidth}\raggedright
Esperado
\end{minipage} & \begin{minipage}[b]{\linewidth}\raggedright
Obtenido
\end{minipage} & \begin{minipage}[b]{\linewidth}\raggedright
Tipo de caso
\end{minipage} \\
\midrule\noalign{}
\endfirsthead
\toprule\noalign{}
\begin{minipage}[b]{\linewidth}\raggedright
Cadena
\end{minipage} & \begin{minipage}[b]{\linewidth}\centering
Longitud
\end{minipage} & \begin{minipage}[b]{\linewidth}\raggedright
Esperado
\end{minipage} & \begin{minipage}[b]{\linewidth}\raggedright
Obtenido
\end{minipage} & \begin{minipage}[b]{\linewidth}\raggedright
Tipo de caso
\end{minipage} \\
\midrule\noalign{}
\endhead
\bottomrule\noalign{}
\endlastfoot
\texttt{010020100201002} & 15 & acepta & acepta & válida \\
\texttt{0100200100201002} & 16 & acepta & acepta & válida \\
\texttt{010020201002010021} & 18 & acepta & acepta & válida \\
λ & 0 & rechaza & rechaza & inválida \\
\texttt{1} & 1 & rechaza & rechaza & inválida \\
\texttt{022} & 3 & rechaza & rechaza & inválida \\
\texttt{21120022200202} & 14 & rechaza & rechaza & frontera: justo por debajo de la longitud mínima \\
\end{longtable}

\evidenciaaqui[0.95]{../evidencias/jflap/fichas/lista3-afn-111.png}{lista3-afn-111 en JFLAP: Multiple Run con las tres cadenas válidas y las
tres inválidas de la tabla anterior.}

\hypertarget{lista-3-afn-del-problema-112}{%
\subsubsection{Lista 3 · AFN del problema 112}\label{lista-3-afn-del-problema-112}}

\textbf{Enunciado.} AFN para el problema 112: Las palabras de Σ1 que contienen dos veces 010 sin traslaparse (01010 no pertenece, 010010 sí).

\textbf{Diseño.} Σ\emph{010Σ}010Σ*: dos copias de 010 separadas por un ciclo; la segunda empieza después de que termina la primera. El enunciado se resuelve de forma natural con un AFN de 7 estados; el AFD equivalente se obtuvo con la construcción de subconjuntos y la minimización y tiene 7 estados.

\textbf{Archivo.} \texttt{automatas/lista3/lista3-afn-112.jff} · Σ = \{0, 1, 2\} · \textbar Q\textbar{} = 7

\begin{longtable}[]{@{}
  >{\raggedright\arraybackslash}p{(\columnwidth - 8\tabcolsep) * \real{0.1364}}
  >{\raggedright\arraybackslash}p{(\columnwidth - 8\tabcolsep) * \real{0.1364}}
  >{\raggedright\arraybackslash}p{(\columnwidth - 8\tabcolsep) * \real{0.1364}}
  >{\raggedright\arraybackslash}p{(\columnwidth - 8\tabcolsep) * \real{0.1364}}
  >{\raggedright\arraybackslash}p{(\columnwidth - 8\tabcolsep) * \real{0.4545}}@{}}
\caption{Tabla de transiciones (→ inicial, * aceptación, --- sin transición).}\tabularnewline
\toprule\noalign{}
\begin{minipage}[b]{\linewidth}\raggedright
δ
\end{minipage} & \begin{minipage}[b]{\linewidth}\raggedright
0
\end{minipage} & \begin{minipage}[b]{\linewidth}\raggedright
1
\end{minipage} & \begin{minipage}[b]{\linewidth}\raggedright
2
\end{minipage} & \begin{minipage}[b]{\linewidth}\raggedright
Qué recuerda del prefijo leído
\end{minipage} \\
\midrule\noalign{}
\endfirsthead
\toprule\noalign{}
\begin{minipage}[b]{\linewidth}\raggedright
δ
\end{minipage} & \begin{minipage}[b]{\linewidth}\raggedright
0
\end{minipage} & \begin{minipage}[b]{\linewidth}\raggedright
1
\end{minipage} & \begin{minipage}[b]{\linewidth}\raggedright
2
\end{minipage} & \begin{minipage}[b]{\linewidth}\raggedright
Qué recuerda del prefijo leído
\end{minipage} \\
\midrule\noalign{}
\endhead
\bottomrule\noalign{}
\endlastfoot
→ q0 & q1 & q0 & q0 & activos del AFN: \{i\} \\
q1 & q1 & q2 & q0 & activos del AFN: \{i,a1\} \\
q2 & q3 & q0 & q0 & activos del AFN: \{i,a2\} \\
q3 & q4 & q3 & q3 & activos del AFN: \{i,a1,a3\} o \{i,a2,a3\} o \{i,a3\} \\
q4 & q4 & q5 & q3 & activos del AFN: \{i,a1,a3,b1\} \\
q5 & q6 & q3 & q3 & activos del AFN: \{i,a2,a3,b2\} \\
* q6 & q6 & q6 & q6 & activos del AFN: \{i,a1,a3,b1,b3\} o \{i,a2,a3,b2,b3\} o \{i,a3,b3\} \\
\end{longtable}

\begin{longtable}[]{@{}
  >{\raggedright\arraybackslash}p{(\columnwidth - 8\tabcolsep) * \real{0.2000}}
  >{\centering\arraybackslash}p{(\columnwidth - 8\tabcolsep) * \real{0.2000}}
  >{\raggedright\arraybackslash}p{(\columnwidth - 8\tabcolsep) * \real{0.2000}}
  >{\raggedright\arraybackslash}p{(\columnwidth - 8\tabcolsep) * \real{0.2000}}
  >{\raggedright\arraybackslash}p{(\columnwidth - 8\tabcolsep) * \real{0.2000}}@{}}
\caption{Cadenas de prueba.}\tabularnewline
\toprule\noalign{}
\begin{minipage}[b]{\linewidth}\raggedright
Cadena
\end{minipage} & \begin{minipage}[b]{\linewidth}\centering
Longitud
\end{minipage} & \begin{minipage}[b]{\linewidth}\raggedright
Esperado
\end{minipage} & \begin{minipage}[b]{\linewidth}\raggedright
Obtenido
\end{minipage} & \begin{minipage}[b]{\linewidth}\raggedright
Tipo de caso
\end{minipage} \\
\midrule\noalign{}
\endfirsthead
\toprule\noalign{}
\begin{minipage}[b]{\linewidth}\raggedright
Cadena
\end{minipage} & \begin{minipage}[b]{\linewidth}\centering
Longitud
\end{minipage} & \begin{minipage}[b]{\linewidth}\raggedright
Esperado
\end{minipage} & \begin{minipage}[b]{\linewidth}\raggedright
Obtenido
\end{minipage} & \begin{minipage}[b]{\linewidth}\raggedright
Tipo de caso
\end{minipage} \\
\midrule\noalign{}
\endhead
\bottomrule\noalign{}
\endlastfoot
\texttt{010010} & 6 & acepta & acepta & válida \\
\texttt{0100010} & 7 & acepta & acepta & válida \\
\texttt{100100100} & 9 & acepta & acepta & válida \\
λ & 0 & rechaza & rechaza & inválida \\
\texttt{1} & 1 & rechaza & rechaza & inválida \\
\texttt{122} & 3 & rechaza & rechaza & inválida \\
\texttt{10002} & 5 & rechaza & rechaza & frontera: justo por debajo de la longitud mínima \\
\end{longtable}

\evidenciaaqui[0.95]{../evidencias/jflap/fichas/lista3-afn-112.png}{lista3-afn-112 en JFLAP: Multiple Run con las tres cadenas válidas y las
tres inválidas de la tabla anterior.}

\hypertarget{lista-3-afn-del-problema-113}{%
\subsubsection{Lista 3 · AFN del problema 113}\label{lista-3-afn-del-problema-113}}

\textbf{Enunciado.} AFN para el problema 113: Las palabras de Σ1 que contienen dos veces 010, que sí se pueden traslapar.

\textbf{Diseño.} Σ\emph{(010Σ}010 + 01010)Σ*: dos apariciones disjuntas o una sola cadena 01010 en la que comparten el 0 central. El enunciado se resuelve de forma natural con un AFN de 11 estados; el AFD equivalente se obtuvo con la construcción de subconjuntos y la minimización y tiene 7 estados.

\textbf{Archivo.} \texttt{automatas/lista3/lista3-afn-113.jff} · Σ = \{0, 1, 2\} · \textbar Q\textbar{} = 7

\begin{longtable}[]{@{}
  >{\raggedright\arraybackslash}p{(\columnwidth - 8\tabcolsep) * \real{0.1364}}
  >{\raggedright\arraybackslash}p{(\columnwidth - 8\tabcolsep) * \real{0.1364}}
  >{\raggedright\arraybackslash}p{(\columnwidth - 8\tabcolsep) * \real{0.1364}}
  >{\raggedright\arraybackslash}p{(\columnwidth - 8\tabcolsep) * \real{0.1364}}
  >{\raggedright\arraybackslash}p{(\columnwidth - 8\tabcolsep) * \real{0.4545}}@{}}
\caption{Tabla de transiciones (→ inicial, * aceptación, --- sin transición).}\tabularnewline
\toprule\noalign{}
\begin{minipage}[b]{\linewidth}\raggedright
δ
\end{minipage} & \begin{minipage}[b]{\linewidth}\raggedright
0
\end{minipage} & \begin{minipage}[b]{\linewidth}\raggedright
1
\end{minipage} & \begin{minipage}[b]{\linewidth}\raggedright
2
\end{minipage} & \begin{minipage}[b]{\linewidth}\raggedright
Qué recuerda del prefijo leído
\end{minipage} \\
\midrule\noalign{}
\endfirsthead
\toprule\noalign{}
\begin{minipage}[b]{\linewidth}\raggedright
δ
\end{minipage} & \begin{minipage}[b]{\linewidth}\raggedright
0
\end{minipage} & \begin{minipage}[b]{\linewidth}\raggedright
1
\end{minipage} & \begin{minipage}[b]{\linewidth}\raggedright
2
\end{minipage} & \begin{minipage}[b]{\linewidth}\raggedright
Qué recuerda del prefijo leído
\end{minipage} \\
\midrule\noalign{}
\endhead
\bottomrule\noalign{}
\endlastfoot
→ q0 & q1 & q0 & q0 & activos del AFN: \{i\} \\
q1 & q1 & q2 & q0 & activos del AFN: \{i,a1,o1\} \\
q2 & q3 & q0 & q0 & activos del AFN: \{i,a2,o2\} \\
q3 & q3 & q4 & q5 & activos del AFN: \{i,a1,a3,o1,o3\} o \{i,a1,a3,b1,o1\} \\
q4 & q6 & q5 & q5 & activos del AFN: \{i,a2,a3,o2,o4\} o \{i,a2,a3,b2,o2\} \\
q5 & q3 & q5 & q5 & activos del AFN: \{i,a3\} \\
* q6 & q6 & q6 & q6 & activos del AFN: \{i,a1,a3,b1,b3,o1,o3\} o \{i,a1,a3,b1,b3,o1\} o \{i,a2,a3,b2,b3,o2,o4\} o \{i,a3,b3\} o \{i,a2,a3,b2,b3,o2\} \\
\end{longtable}

\begin{longtable}[]{@{}
  >{\raggedright\arraybackslash}p{(\columnwidth - 8\tabcolsep) * \real{0.2000}}
  >{\centering\arraybackslash}p{(\columnwidth - 8\tabcolsep) * \real{0.2000}}
  >{\raggedright\arraybackslash}p{(\columnwidth - 8\tabcolsep) * \real{0.2000}}
  >{\raggedright\arraybackslash}p{(\columnwidth - 8\tabcolsep) * \real{0.2000}}
  >{\raggedright\arraybackslash}p{(\columnwidth - 8\tabcolsep) * \real{0.2000}}@{}}
\caption{Cadenas de prueba.}\tabularnewline
\toprule\noalign{}
\begin{minipage}[b]{\linewidth}\raggedright
Cadena
\end{minipage} & \begin{minipage}[b]{\linewidth}\centering
Longitud
\end{minipage} & \begin{minipage}[b]{\linewidth}\raggedright
Esperado
\end{minipage} & \begin{minipage}[b]{\linewidth}\raggedright
Obtenido
\end{minipage} & \begin{minipage}[b]{\linewidth}\raggedright
Tipo de caso
\end{minipage} \\
\midrule\noalign{}
\endfirsthead
\toprule\noalign{}
\begin{minipage}[b]{\linewidth}\raggedright
Cadena
\end{minipage} & \begin{minipage}[b]{\linewidth}\centering
Longitud
\end{minipage} & \begin{minipage}[b]{\linewidth}\raggedright
Esperado
\end{minipage} & \begin{minipage}[b]{\linewidth}\raggedright
Obtenido
\end{minipage} & \begin{minipage}[b]{\linewidth}\raggedright
Tipo de caso
\end{minipage} \\
\midrule\noalign{}
\endhead
\bottomrule\noalign{}
\endlastfoot
\texttt{01010} & 5 & acepta & acepta & válida \\
\texttt{010101} & 6 & acepta & acepta & válida \\
\texttt{20101001} & 8 & acepta & acepta & válida \\
λ & 0 & rechaza & rechaza & inválida \\
\texttt{0} & 1 & rechaza & rechaza & inválida \\
\texttt{222} & 3 & rechaza & rechaza & inválida \\
\texttt{0111} & 4 & rechaza & rechaza & frontera: justo por debajo de la longitud mínima \\
\end{longtable}

\evidenciaaqui[0.95]{../evidencias/jflap/fichas/lista3-afn-113.png}{lista3-afn-113 en JFLAP: Multiple Run con las tres cadenas válidas y las
tres inválidas de la tabla anterior.}

\hypertarget{lista-3-afn-del-problema-116}{%
\subsubsection{Lista 3 · AFN del problema 116}\label{lista-3-afn-del-problema-116}}

\textbf{Enunciado.} AFN para el problema 116: L = \{w ∈ Σ2* : w = aⁿbᵐcⁿ⁺¹, m, n ∈ N\}.

El lenguaje no es regular; tampoco existe AFN.

\hypertarget{lista-3-afn-del-problema-118}{%
\subsubsection{Lista 3 · AFN del problema 118}\label{lista-3-afn-del-problema-118}}

\textbf{Enunciado.} AFN para el problema 118: L = \{w ∈ Σ3* : w = axabya, x, y ∈ Σ3*\}.

\textbf{Diseño.} Sigue literalmente axabya: a, ciclo x, ab, ciclo y, a final. El enunciado se resuelve de forma natural con un AFN de 5 estados; el AFD equivalente se obtuvo con la construcción de subconjuntos y la minimización y tiene 5 estados.

\textbf{Archivo.} \texttt{automatas/lista3/lista3-afn-118.jff} · Σ = \{a, b\} · \textbar Q\textbar{} = 5

\begin{longtable}[]{@{}llll@{}}
\caption{Tabla de transiciones (→ inicial, * aceptación, --- sin transición).}\tabularnewline
\toprule\noalign{}
δ & a & b & Qué recuerda del prefijo leído \\
\midrule\noalign{}
\endfirsthead
\toprule\noalign{}
δ & a & b & Qué recuerda del prefijo leído \\
\midrule\noalign{}
\endhead
\bottomrule\noalign{}
\endlastfoot
→ q0 & q1 & --- & activos del AFN: \{q0\} \\
q1 & q2 & q1 & activos del AFN: \{q1\} \\
q2 & q2 & q3 & activos del AFN: \{q1,q2\} \\
q3 & q4 & q3 & activos del AFN: \{q1,q3\} \\
* q4 & q4 & q3 & activos del AFN: \{q1,q2,q3,q4\} \\
\end{longtable}

\begin{longtable}[]{@{}lclll@{}}
\caption{Cadenas de prueba.}\tabularnewline
\toprule\noalign{}
Cadena & Longitud & Esperado & Obtenido & Tipo de caso \\
\midrule\noalign{}
\endfirsthead
\toprule\noalign{}
Cadena & Longitud & Esperado & Obtenido & Tipo de caso \\
\midrule\noalign{}
\endhead
\bottomrule\noalign{}
\endlastfoot
\texttt{aaba} & 4 & acepta & acepta & válida \\
\texttt{aabba} & 5 & acepta & acepta & válida \\
\texttt{ababaaa} & 7 & acepta & acepta & válida \\
λ & 0 & rechaza & rechaza & inválida \\
\texttt{a} & 1 & rechaza & rechaza & inválida \\
\texttt{baa} & 3 & rechaza & rechaza & inválida \\
\end{longtable}

\evidenciaaqui[0.95]{../evidencias/jflap/fichas/lista3-afn-118.png}{lista3-afn-118 en JFLAP: Multiple Run con las tres cadenas válidas y las
tres inválidas de la tabla anterior.}

\hypertarget{lista-3-afn-del-problema-120}{%
\subsubsection{Lista 3 · AFN del problema 120}\label{lista-3-afn-del-problema-120}}

\textbf{Enunciado.} AFN para el problema 120: Las palabras de Σ1 tales que el primer carácter y el último son distintos.

\textbf{Diseño.} El estado sₓ recuerda el primer símbolo x; el AFN adivina el último símbolo, que debe ser distinto de x. El enunciado se resuelve de forma natural con un AFN de 5 estados; el AFD equivalente se obtuvo con la construcción de subconjuntos y la minimización y tiene 7 estados.

\textbf{Archivo.} \texttt{automatas/lista3/lista3-afn-120.jff} · Σ = \{0, 1, 2\} · \textbar Q\textbar{} = 7

\begin{longtable}[]{@{}lllll@{}}
\caption{Tabla de transiciones (→ inicial, * aceptación, --- sin transición).}\tabularnewline
\toprule\noalign{}
δ & 0 & 1 & 2 & Qué recuerda del prefijo leído \\
\midrule\noalign{}
\endfirsthead
\toprule\noalign{}
δ & 0 & 1 & 2 & Qué recuerda del prefijo leído \\
\midrule\noalign{}
\endhead
\bottomrule\noalign{}
\endlastfoot
→ q0 & q1 & q2 & q3 & activos del AFN: \{q0\} \\
q1 & q1 & q4 & q4 & activos del AFN: \{s0\} \\
q2 & q5 & q2 & q5 & activos del AFN: \{s1\} \\
q3 & q6 & q6 & q3 & activos del AFN: \{s2\} \\
* q4 & q1 & q4 & q4 & activos del AFN: \{s0,f\} \\
* q5 & q5 & q2 & q5 & activos del AFN: \{s1,f\} \\
* q6 & q6 & q6 & q3 & activos del AFN: \{s2,f\} \\
\end{longtable}

\begin{longtable}[]{@{}
  >{\raggedright\arraybackslash}p{(\columnwidth - 8\tabcolsep) * \real{0.2000}}
  >{\centering\arraybackslash}p{(\columnwidth - 8\tabcolsep) * \real{0.2000}}
  >{\raggedright\arraybackslash}p{(\columnwidth - 8\tabcolsep) * \real{0.2000}}
  >{\raggedright\arraybackslash}p{(\columnwidth - 8\tabcolsep) * \real{0.2000}}
  >{\raggedright\arraybackslash}p{(\columnwidth - 8\tabcolsep) * \real{0.2000}}@{}}
\caption{Cadenas de prueba.}\tabularnewline
\toprule\noalign{}
\begin{minipage}[b]{\linewidth}\raggedright
Cadena
\end{minipage} & \begin{minipage}[b]{\linewidth}\centering
Longitud
\end{minipage} & \begin{minipage}[b]{\linewidth}\raggedright
Esperado
\end{minipage} & \begin{minipage}[b]{\linewidth}\raggedright
Obtenido
\end{minipage} & \begin{minipage}[b]{\linewidth}\raggedright
Tipo de caso
\end{minipage} \\
\midrule\noalign{}
\endfirsthead
\toprule\noalign{}
\begin{minipage}[b]{\linewidth}\raggedright
Cadena
\end{minipage} & \begin{minipage}[b]{\linewidth}\centering
Longitud
\end{minipage} & \begin{minipage}[b]{\linewidth}\raggedright
Esperado
\end{minipage} & \begin{minipage}[b]{\linewidth}\raggedright
Obtenido
\end{minipage} & \begin{minipage}[b]{\linewidth}\raggedright
Tipo de caso
\end{minipage} \\
\midrule\noalign{}
\endhead
\bottomrule\noalign{}
\endlastfoot
\texttt{12} & 2 & acepta & acepta & válida \\
\texttt{211} & 3 & acepta & acepta & válida \\
\texttt{12002} & 5 & acepta & acepta & válida \\
λ & 0 & rechaza & rechaza & inválida \\
\texttt{1} & 1 & rechaza & rechaza & inválida \\
\texttt{222} & 3 & rechaza & rechaza & inválida \\
\texttt{02} & 2 & acepta & acepta & frontera: aceptada más corta \\
\texttt{2} & 1 & rechaza & rechaza & frontera: justo por debajo de la longitud mínima \\
\end{longtable}

\evidenciaaqui[0.95]{../evidencias/jflap/fichas/lista3-afn-120.png}{lista3-afn-120 en JFLAP: Multiple Run con las tres cadenas válidas y las
tres inválidas de la tabla anterior.}

\hypertarget{lista-3-afn-del-problema-121}{%
\subsubsection{Lista 3 · AFN del problema 121}\label{lista-3-afn-del-problema-121}}

\textbf{Enunciado.} AFN para el problema 121: Las palabras de Σ1 de longitud 3n tales que cada bloque de longitud 3 tiene al menos un 0.

\textbf{Diseño.} No hay una versión no determinista más simple: el lenguaje depende de contar (longitud, apariciones o fases). Se presenta el AFD del problema sin el estado muerto; una función de transición parcial es un caso particular de AFN.

\textbf{Archivo.} \texttt{automatas/lista3/lista3-afn-121.jff} · Σ = \{0, 1, 2\} · \textbar Q\textbar{} = 5

\begin{longtable}[]{@{}llll@{}}
\caption{Tabla de transiciones (→ inicial, * aceptación, --- sin transición).}\tabularnewline
\toprule\noalign{}
δ & 0 & 1 & 2 \\
\midrule\noalign{}
\endfirsthead
\toprule\noalign{}
δ & 0 & 1 & 2 \\
\midrule\noalign{}
\endhead
\bottomrule\noalign{}
\endlastfoot
→* b0 & b1+0 & b1 & b1 \\
b1+0 & b2+0 & b2+0 & b2+0 \\
b1 & b2+0 & b2 & b2 \\
b2+0 & b0 & b0 & b0 \\
b2 & b0 & --- & --- \\
\end{longtable}

\begin{longtable}[]{@{}lclll@{}}
\caption{Cadenas de prueba.}\tabularnewline
\toprule\noalign{}
Cadena & Longitud & Esperado & Obtenido & Tipo de caso \\
\midrule\noalign{}
\endfirsthead
\toprule\noalign{}
Cadena & Longitud & Esperado & Obtenido & Tipo de caso \\
\midrule\noalign{}
\endhead
\bottomrule\noalign{}
\endlastfoot
λ & 0 & acepta & acepta & válida \\
\texttt{021} & 3 & acepta & acepta & válida \\
\texttt{020001} & 6 & acepta & acepta & válida \\
\texttt{2} & 1 & rechaza & rechaza & inválida \\
\texttt{01} & 2 & rechaza & rechaza & inválida \\
\texttt{0022} & 4 & rechaza & rechaza & inválida \\
\end{longtable}

\evidenciaaqui[0.95]{../evidencias/jflap/fichas/lista3-afn-121.png}{lista3-afn-121 en JFLAP: Multiple Run con las tres cadenas válidas y las
tres inválidas de la tabla anterior.}

\hypertarget{lista-3-afn-del-problema-122}{%
\subsubsection{Lista 3 · AFN del problema 122}\label{lista-3-afn-del-problema-122}}

\textbf{Enunciado.} AFN para el problema 122: Las palabras de Σ1 tales que cualquier cadena de tres símbolos consecutivos tiene un 0.

\textbf{Diseño.} No hay una versión no determinista más simple: el lenguaje depende de contar (longitud, apariciones o fases). Se presenta el AFD del problema sin el estado muerto; una función de transición parcial es un caso particular de AFN.

\textbf{Archivo.} \texttt{automatas/lista3/lista3-afn-122.jff} · Σ = \{0, 1, 2\} · \textbar Q\textbar{} = 3

\begin{longtable}[]{@{}llll@{}}
\caption{Tabla de transiciones (→ inicial, * aceptación, --- sin transición).}\tabularnewline
\toprule\noalign{}
δ & 0 & 1 & 2 \\
\midrule\noalign{}
\endfirsthead
\toprule\noalign{}
δ & 0 & 1 & 2 \\
\midrule\noalign{}
\endhead
\bottomrule\noalign{}
\endlastfoot
→* 0 sin 0 & 0 sin 0 & 1 sin 0 & 1 sin 0 \\
* 1 sin 0 & 0 sin 0 & 2 sin 0 & 2 sin 0 \\
* 2 sin 0 & 0 sin 0 & --- & --- \\
\end{longtable}

\begin{longtable}[]{@{}lclll@{}}
\caption{Cadenas de prueba.}\tabularnewline
\toprule\noalign{}
Cadena & Longitud & Esperado & Obtenido & Tipo de caso \\
\midrule\noalign{}
\endfirsthead
\toprule\noalign{}
Cadena & Longitud & Esperado & Obtenido & Tipo de caso \\
\midrule\noalign{}
\endhead
\bottomrule\noalign{}
\endlastfoot
λ & 0 & acepta & acepta & válida \\
\texttt{1} & 1 & acepta & acepta & válida \\
\texttt{012} & 3 & acepta & acepta & válida \\
\texttt{112} & 3 & rechaza & rechaza & inválida \\
\texttt{1120} & 4 & rechaza & rechaza & inválida \\
\texttt{202211} & 6 & rechaza & rechaza & inválida \\
\end{longtable}

\evidenciaaqui[0.95]{../evidencias/jflap/fichas/lista3-afn-122.png}{lista3-afn-122 en JFLAP: Multiple Run con las tres cadenas válidas y las
tres inválidas de la tabla anterior.}

\hypertarget{lista-3-afn-del-problema-123}{%
\subsubsection{Lista 3 · AFN del problema 123}\label{lista-3-afn-del-problema-123}}

\textbf{Enunciado.} AFN para el problema 123: Las palabras de Σ3 de la forma wxw donde w ∈ Σ3² y x ∈ Σ3*.

\textbf{Diseño.} Una rama por cada w ∈ Σ3²: se lee w, un ciclo para x y otra vez w. El AFN adivina cuándo empieza la copia final. El enunciado se resuelve de forma natural con un AFN de 14 estados; el AFD equivalente se obtuvo con la construcción de subconjuntos y la minimización y tiene 15 estados.

\textbf{Archivo.} \texttt{automatas/lista3/lista3-afn-123.jff} · Σ = \{a, b\} · \textbar Q\textbar{} = 15

\begin{longtable}[]{@{}llll@{}}
\caption{Tabla de transiciones (→ inicial, * aceptación, --- sin transición).}\tabularnewline
\toprule\noalign{}
δ & a & b & Qué recuerda del prefijo leído \\
\midrule\noalign{}
\endfirsthead
\toprule\noalign{}
δ & a & b & Qué recuerda del prefijo leído \\
\midrule\noalign{}
\endhead
\bottomrule\noalign{}
\endlastfoot
→ q0 & q1 & q2 & activos del AFN: \{q0\} \\
q1 & q3 & q4 & activos del AFN: \{aa1,ab1\} \\
q2 & q5 & q6 & activos del AFN: \{ba1,bb1\} \\
q3 & q7 & q3 & activos del AFN: \{aa2\} \\
q4 & q8 & q4 & activos del AFN: \{ab2\} \\
q5 & q5 & q9 & activos del AFN: \{ba2\} \\
q6 & q6 & q10 & activos del AFN: \{bb2\} \\
q7 & q11 & q3 & activos del AFN: \{aa2,aa3\} \\
q8 & q8 & q12 & activos del AFN: \{ab2,ab3\} \\
q9 & q13 & q9 & activos del AFN: \{ba2,ba3\} \\
q10 & q6 & q14 & activos del AFN: \{bb2,bb3\} \\
* q11 & q11 & q3 & activos del AFN: \{aa2,aa3,F\} \\
* q12 & q8 & q4 & activos del AFN: \{ab2,F\} \\
* q13 & q5 & q9 & activos del AFN: \{ba2,F\} \\
* q14 & q6 & q14 & activos del AFN: \{bb2,bb3,F\} \\
\end{longtable}

\begin{longtable}[]{@{}
  >{\raggedright\arraybackslash}p{(\columnwidth - 8\tabcolsep) * \real{0.2000}}
  >{\centering\arraybackslash}p{(\columnwidth - 8\tabcolsep) * \real{0.2000}}
  >{\raggedright\arraybackslash}p{(\columnwidth - 8\tabcolsep) * \real{0.2000}}
  >{\raggedright\arraybackslash}p{(\columnwidth - 8\tabcolsep) * \real{0.2000}}
  >{\raggedright\arraybackslash}p{(\columnwidth - 8\tabcolsep) * \real{0.2000}}@{}}
\caption{Cadenas de prueba.}\tabularnewline
\toprule\noalign{}
\begin{minipage}[b]{\linewidth}\raggedright
Cadena
\end{minipage} & \begin{minipage}[b]{\linewidth}\centering
Longitud
\end{minipage} & \begin{minipage}[b]{\linewidth}\raggedright
Esperado
\end{minipage} & \begin{minipage}[b]{\linewidth}\raggedright
Obtenido
\end{minipage} & \begin{minipage}[b]{\linewidth}\raggedright
Tipo de caso
\end{minipage} \\
\midrule\noalign{}
\endfirsthead
\toprule\noalign{}
\begin{minipage}[b]{\linewidth}\raggedright
Cadena
\end{minipage} & \begin{minipage}[b]{\linewidth}\centering
Longitud
\end{minipage} & \begin{minipage}[b]{\linewidth}\raggedright
Esperado
\end{minipage} & \begin{minipage}[b]{\linewidth}\raggedright
Obtenido
\end{minipage} & \begin{minipage}[b]{\linewidth}\raggedright
Tipo de caso
\end{minipage} \\
\midrule\noalign{}
\endhead
\bottomrule\noalign{}
\endlastfoot
\texttt{abab} & 4 & acepta & acepta & válida \\
\texttt{aaaaa} & 5 & acepta & acepta & válida \\
\texttt{aababaa} & 7 & acepta & acepta & válida \\
λ & 0 & rechaza & rechaza & inválida \\
\texttt{a} & 1 & rechaza & rechaza & inválida \\
\texttt{bba} & 3 & rechaza & rechaza & inválida \\
\texttt{aaaa} & 4 & acepta & acepta & frontera: aceptada más corta \\
\texttt{aab} & 3 & rechaza & rechaza & frontera: justo por debajo de la longitud mínima \\
\end{longtable}

\evidenciaaqui[0.95]{../evidencias/jflap/fichas/lista3-afn-123.png}{lista3-afn-123 en JFLAP: Multiple Run con las tres cadenas válidas y las
tres inválidas de la tabla anterior.}

\hypertarget{lista-3-ejercicio-189}{%
\subsubsection{Lista 3 · Ejercicio 189}\label{lista-3-ejercicio-189}}

\textbf{Enunciado.} AFN sobre \{0, 1\} para las palabras que contienen 00101 o 00110; recorrido de 1001011.

\textbf{Diseño.} Los dos patrones comparten el prefijo 001, así que se escribe una sola vez y se bifurca en 01 y 10. El enunciado se resuelve de forma natural con un AFN de 7 estados; el AFD equivalente se obtuvo con la construcción de subconjuntos y la minimización y tiene 7 estados.

\textbf{Archivo.} \texttt{automatas/lista3/lista3-189.jff} · Σ = \{0, 1\} · \textbar Q\textbar{} = 7

\begin{longtable}[]{@{}
  >{\raggedright\arraybackslash}p{(\columnwidth - 6\tabcolsep) * \real{0.1579}}
  >{\raggedright\arraybackslash}p{(\columnwidth - 6\tabcolsep) * \real{0.1579}}
  >{\raggedright\arraybackslash}p{(\columnwidth - 6\tabcolsep) * \real{0.1579}}
  >{\raggedright\arraybackslash}p{(\columnwidth - 6\tabcolsep) * \real{0.5263}}@{}}
\caption{Tabla de transiciones (→ inicial, * aceptación, --- sin transición).}\tabularnewline
\toprule\noalign{}
\begin{minipage}[b]{\linewidth}\raggedright
δ
\end{minipage} & \begin{minipage}[b]{\linewidth}\raggedright
0
\end{minipage} & \begin{minipage}[b]{\linewidth}\raggedright
1
\end{minipage} & \begin{minipage}[b]{\linewidth}\raggedright
Qué recuerda del prefijo leído
\end{minipage} \\
\midrule\noalign{}
\endfirsthead
\toprule\noalign{}
\begin{minipage}[b]{\linewidth}\raggedright
δ
\end{minipage} & \begin{minipage}[b]{\linewidth}\raggedright
0
\end{minipage} & \begin{minipage}[b]{\linewidth}\raggedright
1
\end{minipage} & \begin{minipage}[b]{\linewidth}\raggedright
Qué recuerda del prefijo leído
\end{minipage} \\
\midrule\noalign{}
\endhead
\bottomrule\noalign{}
\endlastfoot
→ q0 & q1 & q0 & activos del AFN: \{i\} \\
q1 & q2 & q0 & activos del AFN: \{i,c1\} \\
q2 & q2 & q3 & activos del AFN: \{i,c1,c2\} \\
q3 & q4 & q5 & activos del AFN: \{i,c3\} \\
q4 & q2 & q6 & activos del AFN: \{i,c1,u4\} \\
q5 & q6 & q0 & activos del AFN: \{i,d4\} \\
* q6 & q6 & q6 & activos del AFN: \{i,f\} o \{i,c1,f\} o \{i,c1,c2,f\} o \{i,c3,f\} o \{i,c1,u4,f\} o \{i,f,d4\} \\
\end{longtable}

\begin{longtable}[]{@{}
  >{\raggedright\arraybackslash}p{(\columnwidth - 8\tabcolsep) * \real{0.2000}}
  >{\centering\arraybackslash}p{(\columnwidth - 8\tabcolsep) * \real{0.2000}}
  >{\raggedright\arraybackslash}p{(\columnwidth - 8\tabcolsep) * \real{0.2000}}
  >{\raggedright\arraybackslash}p{(\columnwidth - 8\tabcolsep) * \real{0.2000}}
  >{\raggedright\arraybackslash}p{(\columnwidth - 8\tabcolsep) * \real{0.2000}}@{}}
\caption{Cadenas de prueba.}\tabularnewline
\toprule\noalign{}
\begin{minipage}[b]{\linewidth}\raggedright
Cadena
\end{minipage} & \begin{minipage}[b]{\linewidth}\centering
Longitud
\end{minipage} & \begin{minipage}[b]{\linewidth}\raggedright
Esperado
\end{minipage} & \begin{minipage}[b]{\linewidth}\raggedright
Obtenido
\end{minipage} & \begin{minipage}[b]{\linewidth}\raggedright
Tipo de caso
\end{minipage} \\
\midrule\noalign{}
\endfirsthead
\toprule\noalign{}
\begin{minipage}[b]{\linewidth}\raggedright
Cadena
\end{minipage} & \begin{minipage}[b]{\linewidth}\centering
Longitud
\end{minipage} & \begin{minipage}[b]{\linewidth}\raggedright
Esperado
\end{minipage} & \begin{minipage}[b]{\linewidth}\raggedright
Obtenido
\end{minipage} & \begin{minipage}[b]{\linewidth}\raggedright
Tipo de caso
\end{minipage} \\
\midrule\noalign{}
\endhead
\bottomrule\noalign{}
\endlastfoot
\texttt{00110} & 5 & acepta & acepta & válida \\
\texttt{100110} & 6 & acepta & acepta & válida \\
\texttt{00101100} & 8 & acepta & acepta & válida \\
λ & 0 & rechaza & rechaza & inválida \\
\texttt{0} & 1 & rechaza & rechaza & inválida \\
\texttt{101} & 3 & rechaza & rechaza & inválida \\
\texttt{00101} & 5 & acepta & acepta & frontera: aceptada más corta \\
\texttt{0110} & 4 & rechaza & rechaza & frontera: justo por debajo de la longitud mínima \\
\end{longtable}

\evidenciaaqui[0.95]{../evidencias/jflap/fichas/lista3-189.png}{lista3-189 en JFLAP: Multiple Run con las tres cadenas válidas y las
tres inválidas de la tabla anterior.}

\hypertarget{lista-3-ejercicio-190}{%
\subsubsection{Lista 3 · Ejercicio 190}\label{lista-3-ejercicio-190}}

\textbf{Enunciado.} AFN sobre \{0, 1\} para las palabras que contienen 11010 y 10100; recorrido de 01101010100.

\textbf{Diseño.} Producto de los AFN «contiene 11010» y «contiene 10100»; los patrones pueden traslaparse (110100 contiene ambos), por eso no basta concatenarlos. El enunciado se resuelve de forma natural con un AFN de 26 estados; el AFD equivalente se obtuvo con la construcción de subconjuntos y la minimización y tiene 19 estados.

\textbf{Archivo.} \texttt{automatas/lista3/lista3-190.jff} · Σ = \{0, 1\} · \textbar Q\textbar{} = 19

\begin{longtable}[]{@{}
  >{\raggedright\arraybackslash}p{(\columnwidth - 6\tabcolsep) * \real{0.1579}}
  >{\raggedright\arraybackslash}p{(\columnwidth - 6\tabcolsep) * \real{0.1579}}
  >{\raggedright\arraybackslash}p{(\columnwidth - 6\tabcolsep) * \real{0.1579}}
  >{\raggedright\arraybackslash}p{(\columnwidth - 6\tabcolsep) * \real{0.5263}}@{}}
\caption{Tabla de transiciones (→ inicial, * aceptación, --- sin transición).}\tabularnewline
\toprule\noalign{}
\begin{minipage}[b]{\linewidth}\raggedright
δ
\end{minipage} & \begin{minipage}[b]{\linewidth}\raggedright
0
\end{minipage} & \begin{minipage}[b]{\linewidth}\raggedright
1
\end{minipage} & \begin{minipage}[b]{\linewidth}\raggedright
Qué recuerda del prefijo leído
\end{minipage} \\
\midrule\noalign{}
\endfirsthead
\toprule\noalign{}
\begin{minipage}[b]{\linewidth}\raggedright
δ
\end{minipage} & \begin{minipage}[b]{\linewidth}\raggedright
0
\end{minipage} & \begin{minipage}[b]{\linewidth}\raggedright
1
\end{minipage} & \begin{minipage}[b]{\linewidth}\raggedright
Qué recuerda del prefijo leído
\end{minipage} \\
\midrule\noalign{}
\endhead
\bottomrule\noalign{}
\endlastfoot
→ q0 & q0 & q1 & activos del AFN: \{c0c0\} \\
q1 & q2 & q3 & activos del AFN: \{c0c0,c0c1,c1c0,c1c1\} \\
q2 & q0 & q4 & activos del AFN: \{c0c0,c0c2\} \\
q3 & q5 & q3 & activos del AFN: \{c0c0,c0c1,c1c0,c1c1,c2c0,c2c1\} \\
q4 & q6 & q3 & activos del AFN: \{c0c0,c0c1,c1c0,c1c1,c0c3,c1c3\} \\
q5 & q0 & q7 & activos del AFN: \{c0c0,c0c2,c3c0,c3c2\} \\
q6 & q8 & q4 & activos del AFN: \{c0c0,c0c2,c0c4\} \\
q7 & q9 & q3 & activos del AFN: \{c0c0,c0c1,c1c0,c1c1,c0c3,c1c3,c4c0,c4c1,c4c3\} \\
q8 & q8 & q10 & activos del AFN: \{c0c0,c0c5\} o \{c0c0,c0c2,c0c5\} o \{c0c0,c0c2,c0c4,c0c5\} \\
q9 & q11 & q12 & activos del AFN: \{c0c0,c0c2,c0c4,c5c0,c5c2,c5c4\} \\
q10 & q8 & q13 & activos del AFN: \{c0c0,c0c1,c1c0,c1c1,c0c5,c1c5\} o \{c0c0,c0c1,c1c0,c1c1,c0c3,c1c3,c0c5,c1c5\} \\
* q11 & q11 & q11 & activos del AFN: \{c0c0,c0c5,c5c0,c5c5\} o \{c0c0,c0c1,c1c0,c1c1,c0c5,c5c0,c1c5,c5c1,c5c5\} o \{c0c0,c0c2,c0c5,c5c0,c5c2,c5c5\} o \{c0c0,c0c1,c1c0,c1c1,c2c0,c2c1,c0c5,c5c0,c1c5,c5c1,c5c5,c2c5\} o \{c0c0,c0c1,c1c0,c1c1,c0c3,c1c3,c0c5,c5c0,c1c5,c5c1,c5c3,c5c5\} o \{c0c0,c0c2,c3c0,c3c2,c0c5,c5c0,c5c2,c5c5,c3c5\} o \{c0c0,c0c2,c0c4,c0c5,c5c0,c5c2,c5c4,c5c5\} o \{c0c0,c0c1,c1c0,c1c1,c0c3,c1c3,c4c0,c4c1,c4c3,c0c5,c5c0,c1c5,c5c1,c5c3,c5c5,c4c5\} \\
q12 & q9 & q14 & activos del AFN: \{c0c0,c0c1,c1c0,c1c1,c0c3,c1c3,c5c0,c5c1,c5c3\} o \{c0c0,c0c1,c1c0,c1c1,c0c3,c1c3,c4c0,c4c1,c4c3,c5c0,c5c1,c5c3\} \\
q13 & q15 & q13 & activos del AFN: \{c0c0,c0c1,c1c0,c1c1,c2c0,c2c1,c0c5,c1c5,c2c5\} \\
q14 & q16 & q14 & activos del AFN: \{c0c0,c0c1,c1c0,c1c1,c2c0,c2c1,c5c0,c5c1\} o \{c0c0,c0c1,c1c0,c1c1,c5c0,c5c1\} \\
q15 & q8 & q17 & activos del AFN: \{c0c0,c0c2,c3c0,c3c2,c0c5,c3c5\} \\
q16 & q18 & q12 & activos del AFN: \{c0c0,c0c2,c3c0,c3c2,c5c0,c5c2\} o \{c0c0,c0c2,c5c0,c5c2\} \\
q17 & q11 & q13 & activos del AFN: \{c0c0,c0c1,c1c0,c1c1,c0c3,c1c3,c4c0,c4c1,c4c3,c0c5,c1c5,c4c5\} \\
q18 & q18 & q14 & activos del AFN: \{c0c0,c5c0\} \\
\end{longtable}

\begin{longtable}[]{@{}
  >{\raggedright\arraybackslash}p{(\columnwidth - 8\tabcolsep) * \real{0.2000}}
  >{\centering\arraybackslash}p{(\columnwidth - 8\tabcolsep) * \real{0.2000}}
  >{\raggedright\arraybackslash}p{(\columnwidth - 8\tabcolsep) * \real{0.2000}}
  >{\raggedright\arraybackslash}p{(\columnwidth - 8\tabcolsep) * \real{0.2000}}
  >{\raggedright\arraybackslash}p{(\columnwidth - 8\tabcolsep) * \real{0.2000}}@{}}
\caption{Cadenas de prueba.}\tabularnewline
\toprule\noalign{}
\begin{minipage}[b]{\linewidth}\raggedright
Cadena
\end{minipage} & \begin{minipage}[b]{\linewidth}\centering
Longitud
\end{minipage} & \begin{minipage}[b]{\linewidth}\raggedright
Esperado
\end{minipage} & \begin{minipage}[b]{\linewidth}\raggedright
Obtenido
\end{minipage} & \begin{minipage}[b]{\linewidth}\raggedright
Tipo de caso
\end{minipage} \\
\midrule\noalign{}
\endfirsthead
\toprule\noalign{}
\begin{minipage}[b]{\linewidth}\raggedright
Cadena
\end{minipage} & \begin{minipage}[b]{\linewidth}\centering
Longitud
\end{minipage} & \begin{minipage}[b]{\linewidth}\raggedright
Esperado
\end{minipage} & \begin{minipage}[b]{\linewidth}\raggedright
Obtenido
\end{minipage} & \begin{minipage}[b]{\linewidth}\raggedright
Tipo de caso
\end{minipage} \\
\midrule\noalign{}
\endhead
\bottomrule\noalign{}
\endlastfoot
\texttt{110100} & 6 & acepta & acepta & válida \\
\texttt{1101001} & 7 & acepta & acepta & válida \\
\texttt{110100010} & 9 & acepta & acepta & válida \\
λ & 0 & rechaza & rechaza & inválida \\
\texttt{0} & 1 & rechaza & rechaza & inválida \\
\texttt{001} & 3 & rechaza & rechaza & inválida \\
\texttt{00010} & 5 & rechaza & rechaza & frontera: justo por debajo de la longitud mínima \\
\end{longtable}

\evidenciaaqui[0.95]{../evidencias/jflap/fichas/lista3-190.png}{lista3-190 en JFLAP: Multiple Run con las tres cadenas válidas y las
tres inválidas de la tabla anterior.}

\hypertarget{lista-3-ejercicio-191}{%
\subsubsection{Lista 3 · Ejercicio 191}\label{lista-3-ejercicio-191}}

\textbf{Enunciado.} Instrucciones de 192-195: AFD equivalente, δ̂(Z, ω) y δ̂\_D(Z, ω).

Ejercicio de cálculo sobre un AFN dado; se resuelve en la Tarea 3.

\hypertarget{lista-3-ejercicio-192}{%
\subsubsection{Lista 3 · Ejercicio 192}\label{lista-3-ejercicio-192}}

\textbf{Enunciado.} AFN de la tabla 192: AFD equivalente y δ̂(Z, ω).

\textbf{Diseño.} Construcción de subconjuntos (sólo alcanzables). El AFD asociado tiene un estado por subconjunto, así que δ̂\_D(Z, ω) es el subconjunto δ̂(Z, ω). El enunciado se resuelve de forma natural con un AFN de 3 estados; el AFD equivalente se obtuvo con la construcción de subconjuntos y la minimización y tiene 5 estados.

\textbf{Archivo.} \texttt{automatas/lista3/lista3-192.jff} · Σ = \{0, 1\} · \textbar Q\textbar{} = 5

\begin{longtable}[]{@{}llll@{}}
\caption{Tabla de transiciones (→ inicial, * aceptación, --- sin transición).}\tabularnewline
\toprule\noalign{}
δ & 0 & 1 & Qué recuerda del prefijo leído \\
\midrule\noalign{}
\endfirsthead
\toprule\noalign{}
δ & 0 & 1 & Qué recuerda del prefijo leído \\
\midrule\noalign{}
\endhead
\bottomrule\noalign{}
\endlastfoot
→ q0 & q1 & q2 & activos del AFN: \{A\} \\
q1 & q3 & --- & activos del AFN: \{B\} \\
* q2 & q4 & --- & activos del AFN: \{B,C\} \\
q3 & q3 & q2 & activos del AFN: \{A,B\} \\
* q4 & q4 & q2 & activos del AFN: \{A,B,C\} \\
\end{longtable}

\begin{longtable}[]{@{}lclll@{}}
\caption{Cadenas de prueba.}\tabularnewline
\toprule\noalign{}
Cadena & Longitud & Esperado & Obtenido & Tipo de caso \\
\midrule\noalign{}
\endfirsthead
\toprule\noalign{}
Cadena & Longitud & Esperado & Obtenido & Tipo de caso \\
\midrule\noalign{}
\endhead
\bottomrule\noalign{}
\endlastfoot
\texttt{1} & 1 & acepta & acepta & válida \\
\texttt{10} & 2 & acepta & acepta & válida \\
\texttt{1010} & 4 & acepta & acepta & válida \\
λ & 0 & rechaza & rechaza & inválida \\
\texttt{0} & 1 & rechaza & rechaza & inválida \\
\texttt{111} & 3 & rechaza & rechaza & inválida \\
\end{longtable}

\evidenciaaqui[0.95]{../evidencias/jflap/fichas/lista3-192.png}{lista3-192 en JFLAP: Multiple Run con las tres cadenas válidas y las
tres inválidas de la tabla anterior.}

\hypertarget{lista-3-ejercicio-193}{%
\subsubsection{Lista 3 · Ejercicio 193}\label{lista-3-ejercicio-193}}

\textbf{Enunciado.} AFN de la tabla 193: AFD equivalente y δ̂(Z, ω).

\textbf{Diseño.} Construcción de subconjuntos (sólo alcanzables). El AFD asociado tiene un estado por subconjunto, así que δ̂\_D(Z, ω) es el subconjunto δ̂(Z, ω). El enunciado se resuelve de forma natural con un AFN de 4 estados; el AFD equivalente se obtuvo con la construcción de subconjuntos y la minimización y tiene 6 estados.

\textbf{Archivo.} \texttt{automatas/lista3/lista3-193.jff} · Σ = \{0, 1\} · \textbar Q\textbar{} = 6

\begin{longtable}[]{@{}llll@{}}
\caption{Tabla de transiciones (→ inicial, * aceptación, --- sin transición).}\tabularnewline
\toprule\noalign{}
δ & 0 & 1 & Qué recuerda del prefijo leído \\
\midrule\noalign{}
\endfirsthead
\toprule\noalign{}
δ & 0 & 1 & Qué recuerda del prefijo leído \\
\midrule\noalign{}
\endhead
\bottomrule\noalign{}
\endlastfoot
→ q0 & --- & q1 & activos del AFN: \{A\} \\
q1 & q2 & q3 & activos del AFN: \{B,C\} o \{A,B,C\} \\
* q2 & q4 & q3 & activos del AFN: \{A,C,D\} \\
* q3 & q2 & q3 & activos del AFN: \{A,B,C,D\} \\
q4 & q4 & q5 & activos del AFN: \{C\} \\
q5 & q4 & q1 & activos del AFN: \{A,C\} \\
\end{longtable}

\begin{longtable}[]{@{}
  >{\raggedright\arraybackslash}p{(\columnwidth - 8\tabcolsep) * \real{0.2000}}
  >{\centering\arraybackslash}p{(\columnwidth - 8\tabcolsep) * \real{0.2000}}
  >{\raggedright\arraybackslash}p{(\columnwidth - 8\tabcolsep) * \real{0.2000}}
  >{\raggedright\arraybackslash}p{(\columnwidth - 8\tabcolsep) * \real{0.2000}}
  >{\raggedright\arraybackslash}p{(\columnwidth - 8\tabcolsep) * \real{0.2000}}@{}}
\caption{Cadenas de prueba.}\tabularnewline
\toprule\noalign{}
\begin{minipage}[b]{\linewidth}\raggedright
Cadena
\end{minipage} & \begin{minipage}[b]{\linewidth}\centering
Longitud
\end{minipage} & \begin{minipage}[b]{\linewidth}\raggedright
Esperado
\end{minipage} & \begin{minipage}[b]{\linewidth}\raggedright
Obtenido
\end{minipage} & \begin{minipage}[b]{\linewidth}\raggedright
Tipo de caso
\end{minipage} \\
\midrule\noalign{}
\endfirsthead
\toprule\noalign{}
\begin{minipage}[b]{\linewidth}\raggedright
Cadena
\end{minipage} & \begin{minipage}[b]{\linewidth}\centering
Longitud
\end{minipage} & \begin{minipage}[b]{\linewidth}\raggedright
Esperado
\end{minipage} & \begin{minipage}[b]{\linewidth}\raggedright
Obtenido
\end{minipage} & \begin{minipage}[b]{\linewidth}\raggedright
Tipo de caso
\end{minipage} \\
\midrule\noalign{}
\endhead
\bottomrule\noalign{}
\endlastfoot
\texttt{11} & 2 & acepta & acepta & válida \\
\texttt{110} & 3 & acepta & acepta & válida \\
\texttt{10101} & 5 & acepta & acepta & válida \\
λ & 0 & rechaza & rechaza & inválida \\
\texttt{1} & 1 & rechaza & rechaza & inválida \\
\texttt{000} & 3 & rechaza & rechaza & inválida \\
\texttt{10} & 2 & acepta & acepta & frontera: aceptada más corta \\
\texttt{0} & 1 & rechaza & rechaza & frontera: justo por debajo de la longitud mínima \\
\end{longtable}

\evidenciaaqui[0.95]{../evidencias/jflap/fichas/lista3-193.png}{lista3-193 en JFLAP: Multiple Run con las tres cadenas válidas y las
tres inválidas de la tabla anterior.}

\hypertarget{lista-3-ejercicio-194}{%
\subsubsection{Lista 3 · Ejercicio 194}\label{lista-3-ejercicio-194}}

\textbf{Enunciado.} AFN de la tabla 194: AFD equivalente y δ̂(Z, ω).

\textbf{Diseño.} Construcción de subconjuntos (sólo alcanzables). El AFD asociado tiene un estado por subconjunto, así que δ̂\_D(Z, ω) es el subconjunto δ̂(Z, ω). El enunciado se resuelve de forma natural con un AFN de 4 estados; el AFD equivalente se obtuvo con la construcción de subconjuntos y la minimización y tiene 2 estados.

\textbf{Archivo.} \texttt{automatas/lista3/lista3-194.jff} · Σ = \{0, 1, 2\} · \textbar Q\textbar{} = 2

\begin{longtable}[]{@{}lllll@{}}
\caption{Tabla de transiciones (→ inicial, * aceptación, --- sin transición).}\tabularnewline
\toprule\noalign{}
δ & 0 & 1 & 2 & Qué recuerda del prefijo leído \\
\midrule\noalign{}
\endfirsthead
\toprule\noalign{}
δ & 0 & 1 & 2 & Qué recuerda del prefijo leído \\
\midrule\noalign{}
\endhead
\bottomrule\noalign{}
\endlastfoot
→ q0 & q0 & q1 & q0 & activos del AFN: \{A\} o \{A,B\} \\
* q1 & q0 & q1 & q0 & activos del AFN: \{A,C\} o \{A,B,C,D\} \\
\end{longtable}

\begin{longtable}[]{@{}lclll@{}}
\caption{Cadenas de prueba.}\tabularnewline
\toprule\noalign{}
Cadena & Longitud & Esperado & Obtenido & Tipo de caso \\
\midrule\noalign{}
\endfirsthead
\toprule\noalign{}
Cadena & Longitud & Esperado & Obtenido & Tipo de caso \\
\midrule\noalign{}
\endhead
\bottomrule\noalign{}
\endlastfoot
\texttt{1} & 1 & acepta & acepta & válida \\
\texttt{11} & 2 & acepta & acepta & válida \\
\texttt{0011} & 4 & acepta & acepta & válida \\
λ & 0 & rechaza & rechaza & inválida \\
\texttt{2} & 1 & rechaza & rechaza & inválida \\
\texttt{220} & 3 & rechaza & rechaza & inválida \\
\end{longtable}

\evidenciaaqui[0.95]{../evidencias/jflap/fichas/lista3-194.png}{lista3-194 en JFLAP: Multiple Run con las tres cadenas válidas y las
tres inválidas de la tabla anterior.}

\hypertarget{lista-3-ejercicio-195}{%
\subsubsection{Lista 3 · Ejercicio 195}\label{lista-3-ejercicio-195}}

\textbf{Enunciado.} AFN de la tabla 195: AFD equivalente y δ̂(Z, ω).

\textbf{Diseño.} Construcción de subconjuntos (sólo alcanzables). El AFD asociado tiene un estado por subconjunto, así que δ̂\_D(Z, ω) es el subconjunto δ̂(Z, ω). El enunciado se resuelve de forma natural con un AFN de 5 estados; el AFD equivalente se obtuvo con la construcción de subconjuntos y la minimización y tiene 7 estados.

\textbf{Archivo.} \texttt{automatas/lista3/lista3-195.jff} · Σ = \{0, 1\} · \textbar Q\textbar{} = 7

\begin{longtable}[]{@{}llll@{}}
\caption{Tabla de transiciones (→ inicial, * aceptación, --- sin transición).}\tabularnewline
\toprule\noalign{}
δ & 0 & 1 & Qué recuerda del prefijo leído \\
\midrule\noalign{}
\endfirsthead
\toprule\noalign{}
δ & 0 & 1 & Qué recuerda del prefijo leído \\
\midrule\noalign{}
\endhead
\bottomrule\noalign{}
\endlastfoot
→ q0 & q1 & q0 & activos del AFN: \{A\} \\
q1 & q1 & q2 & activos del AFN: \{A,B\} \\
* q2 & q1 & q3 & activos del AFN: \{A,C\} \\
q3 & q3 & q4 & activos del AFN: \{A,D\} o \{A,B,D\} \\
* q4 & q4 & q5 & activos del AFN: \{A,E\} o \{A,C,E\} o \{A,B,E\} \\
* q5 & q3 & q6 & activos del AFN: \{A,C,D\} \\
* q6 & q6 & q6 & activos del AFN: \{A,D,E\} o \{A,B,D,E\} o \{A,C,D,E\} \\
\end{longtable}

\begin{longtable}[]{@{}lclll@{}}
\caption{Cadenas de prueba.}\tabularnewline
\toprule\noalign{}
Cadena & Longitud & Esperado & Obtenido & Tipo de caso \\
\midrule\noalign{}
\endfirsthead
\toprule\noalign{}
Cadena & Longitud & Esperado & Obtenido & Tipo de caso \\
\midrule\noalign{}
\endhead
\bottomrule\noalign{}
\endlastfoot
\texttt{01} & 2 & acepta & acepta & válida \\
\texttt{101} & 3 & acepta & acepta & válida \\
\texttt{01110} & 5 & acepta & acepta & válida \\
λ & 0 & rechaza & rechaza & inválida \\
\texttt{1} & 1 & rechaza & rechaza & inválida \\
\texttt{011} & 3 & rechaza & rechaza & inválida \\
\end{longtable}

\evidenciaaqui[0.95]{../evidencias/jflap/fichas/lista3-195.png}{lista3-195 en JFLAP: Multiple Run con las tres cadenas válidas y las
tres inválidas de la tabla anterior.}

\hypertarget{lista-3-ejercicio-201-235}{%
\subsubsection{Lista 3 · Ejercicio 201-235}\label{lista-3-ejercicio-201-235}}

\textbf{Enunciado.} Indique si las afirmaciones 201-235 son verdaderas (V) o falsas (F).

Afirmaciones verdadero/falso sobre expresiones regulares; no piden autómata.

\hypertarget{lista-3-ejercicio-236}{%
\subsubsection{Lista 3 · Ejercicio 236}\label{lista-3-ejercicio-236}}

\textbf{Enunciado.} Expresión regular sobre \{0, 1\} para: L = \{x001y : x, y ∈ Σ*\}.

\textbf{Diseño.} La expresión se convierte en AFN-λ con la construcción de Thompson y se verifica contra la definición del lenguaje. El enunciado se resuelve de forma natural con un AFN-λ de 22 estados; el AFD equivalente se obtuvo con la construcción de subconjuntos y la minimización y tiene 4 estados.

\textbf{Archivo.} \texttt{automatas/lista3/lista3-236.jff} · Σ = \{0, 1\} · \textbar Q\textbar{} = 4

\begin{longtable}[]{@{}
  >{\raggedright\arraybackslash}p{(\columnwidth - 6\tabcolsep) * \real{0.1579}}
  >{\raggedright\arraybackslash}p{(\columnwidth - 6\tabcolsep) * \real{0.1579}}
  >{\raggedright\arraybackslash}p{(\columnwidth - 6\tabcolsep) * \real{0.1579}}
  >{\raggedright\arraybackslash}p{(\columnwidth - 6\tabcolsep) * \real{0.5263}}@{}}
\caption{Tabla de transiciones (→ inicial, * aceptación, --- sin transición).}\tabularnewline
\toprule\noalign{}
\begin{minipage}[b]{\linewidth}\raggedright
δ
\end{minipage} & \begin{minipage}[b]{\linewidth}\raggedright
0
\end{minipage} & \begin{minipage}[b]{\linewidth}\raggedright
1
\end{minipage} & \begin{minipage}[b]{\linewidth}\raggedright
Qué recuerda del prefijo leído
\end{minipage} \\
\midrule\noalign{}
\endfirsthead
\toprule\noalign{}
\begin{minipage}[b]{\linewidth}\raggedright
δ
\end{minipage} & \begin{minipage}[b]{\linewidth}\raggedright
0
\end{minipage} & \begin{minipage}[b]{\linewidth}\raggedright
1
\end{minipage} & \begin{minipage}[b]{\linewidth}\raggedright
Qué recuerda del prefijo leído
\end{minipage} \\
\midrule\noalign{}
\endhead
\bottomrule\noalign{}
\endlastfoot
→ q0 & q1 & q0 & activos del AFN: \{q6\} o \{q4,q0,q2,q3,q5,q7,q8\} \\
q1 & q2 & q0 & activos del AFN: \{q4,q0,q1,q2,q5,q7,q8,q9,q10\} \\
q2 & q2 & q3 & activos del AFN: \{q4,q0,q1,q2,q5,q7,q8,q9,q10,q11,q12\} \\
* q3 & q3 & q3 & activos del AFN: \{q4,q0,q2,q3,q5,q7,q8,q13,q20,q18,q14,q16,q21\} o \{q4,q0,q1,q2,q5,q7,q8,q9,q10,q18,q14,q15,q16,q19,q21\} o \{q4,q0,q2,q3,q5,q7,q8,q18,q14,q16,q17,q19,q21\} o \{q4,q0,q1,q2,q5,q7,q8,q9,q10,q11,q12,q18,q14,q15,q16,q19,q21\} o \{q4,q0,q2,q3,q5,q7,q8,q13,q20,q18,q14,q16,q17,q19,q21\} \\
\end{longtable}

\begin{longtable}[]{@{}
  >{\raggedright\arraybackslash}p{(\columnwidth - 8\tabcolsep) * \real{0.2000}}
  >{\centering\arraybackslash}p{(\columnwidth - 8\tabcolsep) * \real{0.2000}}
  >{\raggedright\arraybackslash}p{(\columnwidth - 8\tabcolsep) * \real{0.2000}}
  >{\raggedright\arraybackslash}p{(\columnwidth - 8\tabcolsep) * \real{0.2000}}
  >{\raggedright\arraybackslash}p{(\columnwidth - 8\tabcolsep) * \real{0.2000}}@{}}
\caption{Cadenas de prueba.}\tabularnewline
\toprule\noalign{}
\begin{minipage}[b]{\linewidth}\raggedright
Cadena
\end{minipage} & \begin{minipage}[b]{\linewidth}\centering
Longitud
\end{minipage} & \begin{minipage}[b]{\linewidth}\raggedright
Esperado
\end{minipage} & \begin{minipage}[b]{\linewidth}\raggedright
Obtenido
\end{minipage} & \begin{minipage}[b]{\linewidth}\raggedright
Tipo de caso
\end{minipage} \\
\midrule\noalign{}
\endfirsthead
\toprule\noalign{}
\begin{minipage}[b]{\linewidth}\raggedright
Cadena
\end{minipage} & \begin{minipage}[b]{\linewidth}\centering
Longitud
\end{minipage} & \begin{minipage}[b]{\linewidth}\raggedright
Esperado
\end{minipage} & \begin{minipage}[b]{\linewidth}\raggedright
Obtenido
\end{minipage} & \begin{minipage}[b]{\linewidth}\raggedright
Tipo de caso
\end{minipage} \\
\midrule\noalign{}
\endhead
\bottomrule\noalign{}
\endlastfoot
\texttt{001} & 3 & acepta & acepta & válida \\
\texttt{0001} & 4 & acepta & acepta & válida \\
\texttt{001010} & 6 & acepta & acepta & válida \\
λ & 0 & rechaza & rechaza & inválida \\
\texttt{0} & 1 & rechaza & rechaza & inválida \\
\texttt{100} & 3 & rechaza & rechaza & inválida \\
\texttt{00} & 2 & rechaza & rechaza & frontera: justo por debajo de la longitud mínima \\
\end{longtable}

\evidenciaaqui[0.95]{../evidencias/jflap/fichas/lista3-236.png}{lista3-236 en JFLAP: Multiple Run con las tres cadenas válidas y las
tres inválidas de la tabla anterior.}

\hypertarget{lista-3-ejercicio-237}{%
\subsubsection{Lista 3 · Ejercicio 237}\label{lista-3-ejercicio-237}}

\textbf{Enunciado.} Expresión regular sobre \{0, 1\} para: L = \{x010 : x ∈ Σ*\}.

\textbf{Diseño.} La expresión se convierte en AFN-λ con la construcción de Thompson y se verifica contra la definición del lenguaje. El enunciado se resuelve de forma natural con un AFN-λ de 14 estados; el AFD equivalente se obtuvo con la construcción de subconjuntos y la minimización y tiene 4 estados.

\textbf{Archivo.} \texttt{automatas/lista3/lista3-237.jff} · Σ = \{0, 1\} · \textbar Q\textbar{} = 4

\begin{longtable}[]{@{}llll@{}}
\caption{Tabla de transiciones (→ inicial, * aceptación, --- sin transición).}\tabularnewline
\toprule\noalign{}
δ & 0 & 1 & Qué recuerda del prefijo leído \\
\midrule\noalign{}
\endfirsthead
\toprule\noalign{}
δ & 0 & 1 & Qué recuerda del prefijo leído \\
\midrule\noalign{}
\endhead
\bottomrule\noalign{}
\endlastfoot
→ q0 & q1 & q0 & activos del AFN: \{q6\} o \{q4,q0,q2,q3,q5,q7,q8\} \\
q1 & q1 & q2 & activos del AFN: \{q4,q0,q1,q2,q5,q7,q8,q9,q10\} \\
q2 & q3 & q0 & activos del AFN: \{q4,q0,q2,q3,q5,q7,q8,q11,q12\} \\
* q3 & q1 & q2 & activos del AFN: \{q4,q0,q1,q2,q5,q7,q8,q9,q10,q13\} \\
\end{longtable}

\begin{longtable}[]{@{}
  >{\raggedright\arraybackslash}p{(\columnwidth - 8\tabcolsep) * \real{0.2000}}
  >{\centering\arraybackslash}p{(\columnwidth - 8\tabcolsep) * \real{0.2000}}
  >{\raggedright\arraybackslash}p{(\columnwidth - 8\tabcolsep) * \real{0.2000}}
  >{\raggedright\arraybackslash}p{(\columnwidth - 8\tabcolsep) * \real{0.2000}}
  >{\raggedright\arraybackslash}p{(\columnwidth - 8\tabcolsep) * \real{0.2000}}@{}}
\caption{Cadenas de prueba.}\tabularnewline
\toprule\noalign{}
\begin{minipage}[b]{\linewidth}\raggedright
Cadena
\end{minipage} & \begin{minipage}[b]{\linewidth}\centering
Longitud
\end{minipage} & \begin{minipage}[b]{\linewidth}\raggedright
Esperado
\end{minipage} & \begin{minipage}[b]{\linewidth}\raggedright
Obtenido
\end{minipage} & \begin{minipage}[b]{\linewidth}\raggedright
Tipo de caso
\end{minipage} \\
\midrule\noalign{}
\endfirsthead
\toprule\noalign{}
\begin{minipage}[b]{\linewidth}\raggedright
Cadena
\end{minipage} & \begin{minipage}[b]{\linewidth}\centering
Longitud
\end{minipage} & \begin{minipage}[b]{\linewidth}\raggedright
Esperado
\end{minipage} & \begin{minipage}[b]{\linewidth}\raggedright
Obtenido
\end{minipage} & \begin{minipage}[b]{\linewidth}\raggedright
Tipo de caso
\end{minipage} \\
\midrule\noalign{}
\endhead
\bottomrule\noalign{}
\endlastfoot
\texttt{010} & 3 & acepta & acepta & válida \\
\texttt{0010} & 4 & acepta & acepta & válida \\
\texttt{111010} & 6 & acepta & acepta & válida \\
λ & 0 & rechaza & rechaza & inválida \\
\texttt{0} & 1 & rechaza & rechaza & inválida \\
\texttt{110} & 3 & rechaza & rechaza & inválida \\
\texttt{11} & 2 & rechaza & rechaza & frontera: justo por debajo de la longitud mínima \\
\end{longtable}

\evidenciaaqui[0.95]{../evidencias/jflap/fichas/lista3-237.png}{lista3-237 en JFLAP: Multiple Run con las tres cadenas válidas y las
tres inválidas de la tabla anterior.}

\hypertarget{lista-3-ejercicio-238}{%
\subsubsection{Lista 3 · Ejercicio 238}\label{lista-3-ejercicio-238}}

\textbf{Enunciado.} Expresión regular sobre \{0, 1\} para: L = Σ³ ∪ Σ⁵.

\textbf{Diseño.} La expresión se convierte en AFN-λ con la construcción de Thompson y se verifica contra la definición del lenguaje. El enunciado se resuelve de forma natural con un AFN-λ de 34 estados; el AFD equivalente se obtuvo con la construcción de subconjuntos y la minimización y tiene 6 estados.

\textbf{Archivo.} \texttt{automatas/lista3/lista3-238.jff} · Σ = \{0, 1\} · \textbar Q\textbar{} = 6

\begin{longtable}[]{@{}
  >{\raggedright\arraybackslash}p{(\columnwidth - 6\tabcolsep) * \real{0.1579}}
  >{\raggedright\arraybackslash}p{(\columnwidth - 6\tabcolsep) * \real{0.1579}}
  >{\raggedright\arraybackslash}p{(\columnwidth - 6\tabcolsep) * \real{0.1579}}
  >{\raggedright\arraybackslash}p{(\columnwidth - 6\tabcolsep) * \real{0.5263}}@{}}
\caption{Tabla de transiciones (→ inicial, * aceptación, --- sin transición).}\tabularnewline
\toprule\noalign{}
\begin{minipage}[b]{\linewidth}\raggedright
δ
\end{minipage} & \begin{minipage}[b]{\linewidth}\raggedright
0
\end{minipage} & \begin{minipage}[b]{\linewidth}\raggedright
1
\end{minipage} & \begin{minipage}[b]{\linewidth}\raggedright
Qué recuerda del prefijo leído
\end{minipage} \\
\midrule\noalign{}
\endfirsthead
\toprule\noalign{}
\begin{minipage}[b]{\linewidth}\raggedright
δ
\end{minipage} & \begin{minipage}[b]{\linewidth}\raggedright
0
\end{minipage} & \begin{minipage}[b]{\linewidth}\raggedright
1
\end{minipage} & \begin{minipage}[b]{\linewidth}\raggedright
Qué recuerda del prefijo leído
\end{minipage} \\
\midrule\noalign{}
\endhead
\bottomrule\noalign{}
\endlastfoot
→ q0 & q1 & q1 & activos del AFN: \{q4\} \\
q1 & q2 & q2 & activos del AFN: \{q1,q5,q10,q6,q8\} o \{q3,q5,q10,q6,q8\} \\
q2 & q3 & q3 & activos del AFN: \{q7,q11,q16,q12,q14\} o \{q9,q11,q16,q12,q14\} \\
* q3 & q4 & q4 & activos del AFN: \{q13,q17,q32,q18,q19,q24,q20,q22,q33\} o \{q15,q17,q32,q18,q19,q24,q20,q22,q33\} \\
q4 & q5 & q5 & activos del AFN: \{q21,q25,q30,q26,q28\} o \{q23,q25,q30,q26,q28\} \\
* q5 & --- & --- & activos del AFN: \{q27,q31,q33\} o \{q29,q31,q33\} \\
\end{longtable}

\begin{longtable}[]{@{}
  >{\raggedright\arraybackslash}p{(\columnwidth - 8\tabcolsep) * \real{0.2000}}
  >{\centering\arraybackslash}p{(\columnwidth - 8\tabcolsep) * \real{0.2000}}
  >{\raggedright\arraybackslash}p{(\columnwidth - 8\tabcolsep) * \real{0.2000}}
  >{\raggedright\arraybackslash}p{(\columnwidth - 8\tabcolsep) * \real{0.2000}}
  >{\raggedright\arraybackslash}p{(\columnwidth - 8\tabcolsep) * \real{0.2000}}@{}}
\caption{Cadenas de prueba.}\tabularnewline
\toprule\noalign{}
\begin{minipage}[b]{\linewidth}\raggedright
Cadena
\end{minipage} & \begin{minipage}[b]{\linewidth}\centering
Longitud
\end{minipage} & \begin{minipage}[b]{\linewidth}\raggedright
Esperado
\end{minipage} & \begin{minipage}[b]{\linewidth}\raggedright
Obtenido
\end{minipage} & \begin{minipage}[b]{\linewidth}\raggedright
Tipo de caso
\end{minipage} \\
\midrule\noalign{}
\endfirsthead
\toprule\noalign{}
\begin{minipage}[b]{\linewidth}\raggedright
Cadena
\end{minipage} & \begin{minipage}[b]{\linewidth}\centering
Longitud
\end{minipage} & \begin{minipage}[b]{\linewidth}\raggedright
Esperado
\end{minipage} & \begin{minipage}[b]{\linewidth}\raggedright
Obtenido
\end{minipage} & \begin{minipage}[b]{\linewidth}\raggedright
Tipo de caso
\end{minipage} \\
\midrule\noalign{}
\endhead
\bottomrule\noalign{}
\endlastfoot
\texttt{000} & 3 & acepta & acepta & válida \\
\texttt{001} & 3 & acepta & acepta & válida \\
\texttt{101} & 3 & acepta & acepta & válida \\
λ & 0 & rechaza & rechaza & inválida \\
\texttt{0} & 1 & rechaza & rechaza & inválida \\
\texttt{0100} & 4 & rechaza & rechaza & inválida \\
\texttt{00} & 2 & rechaza & rechaza & frontera: justo por debajo de la longitud mínima \\
\end{longtable}

\evidenciaaqui[0.95]{../evidencias/jflap/fichas/lista3-238.png}{lista3-238 en JFLAP: Multiple Run con las tres cadenas válidas y las
tres inválidas de la tabla anterior.}

\hypertarget{lista-3-ejercicio-239}{%
\subsubsection{Lista 3 · Ejercicio 239}\label{lista-3-ejercicio-239}}

\textbf{Enunciado.} Expresión regular sobre \{0, 1\} para: Cadenas que comienzan con 01 o terminan con 10.

\textbf{Diseño.} La expresión se convierte en AFN-λ con la construcción de Thompson y se verifica contra la definición del lenguaje. El enunciado se resuelve de forma natural con un AFN-λ de 26 estados; el AFD equivalente se obtuvo con la construcción de subconjuntos y la minimización y tiene 6 estados.

\textbf{Archivo.} \texttt{automatas/lista3/lista3-239.jff} · Σ = \{0, 1\} · \textbar Q\textbar{} = 6

\begin{longtable}[]{@{}
  >{\raggedright\arraybackslash}p{(\columnwidth - 6\tabcolsep) * \real{0.1579}}
  >{\raggedright\arraybackslash}p{(\columnwidth - 6\tabcolsep) * \real{0.1579}}
  >{\raggedright\arraybackslash}p{(\columnwidth - 6\tabcolsep) * \real{0.1579}}
  >{\raggedright\arraybackslash}p{(\columnwidth - 6\tabcolsep) * \real{0.5263}}@{}}
\caption{Tabla de transiciones (→ inicial, * aceptación, --- sin transición).}\tabularnewline
\toprule\noalign{}
\begin{minipage}[b]{\linewidth}\raggedright
δ
\end{minipage} & \begin{minipage}[b]{\linewidth}\raggedright
0
\end{minipage} & \begin{minipage}[b]{\linewidth}\raggedright
1
\end{minipage} & \begin{minipage}[b]{\linewidth}\raggedright
Qué recuerda del prefijo leído
\end{minipage} \\
\midrule\noalign{}
\endfirsthead
\toprule\noalign{}
\begin{minipage}[b]{\linewidth}\raggedright
δ
\end{minipage} & \begin{minipage}[b]{\linewidth}\raggedright
0
\end{minipage} & \begin{minipage}[b]{\linewidth}\raggedright
1
\end{minipage} & \begin{minipage}[b]{\linewidth}\raggedright
Qué recuerda del prefijo leído
\end{minipage} \\
\midrule\noalign{}
\endhead
\bottomrule\noalign{}
\endlastfoot
→ q0 & q1 & q2 & activos del AFN: \{q24\} \\
q1 & q3 & q4 & activos del AFN: \{q1,q2,q16,q12,q13,q14,q17,q19,q20\} \\
q2 & q5 & q2 & activos del AFN: \{q16,q12,q14,q15,q17,q19,q20,q21,q22\} \\
q3 & q3 & q2 & activos del AFN: \{q16,q12,q13,q14,q17,q19,q20\} \\
* q4 & q4 & q4 & activos del AFN: \{q3,q10,q8,q4,q6,q11,q16,q12,q14,q15,q17,q19,q20,q21,q22,q25\} o \{q8,q4,q5,q6,q9,q11,q16,q12,q13,q14,q17,q19,q20,q23,q25\} o \{q8,q4,q6,q7,q9,q11,q16,q12,q14,q15,q17,q19,q20,q21,q22,q25\} o \{q8,q4,q5,q6,q9,q11,q16,q12,q13,q14,q17,q19,q20,q25\} \\
* q5 & q3 & q2 & activos del AFN: \{q16,q12,q13,q14,q17,q19,q20,q23,q25\} \\
\end{longtable}

\begin{longtable}[]{@{}lclll@{}}
\caption{Cadenas de prueba.}\tabularnewline
\toprule\noalign{}
Cadena & Longitud & Esperado & Obtenido & Tipo de caso \\
\midrule\noalign{}
\endfirsthead
\toprule\noalign{}
Cadena & Longitud & Esperado & Obtenido & Tipo de caso \\
\midrule\noalign{}
\endhead
\bottomrule\noalign{}
\endlastfoot
\texttt{01} & 2 & acepta & acepta & válida \\
\texttt{010} & 3 & acepta & acepta & válida \\
\texttt{10010} & 5 & acepta & acepta & válida \\
λ & 0 & rechaza & rechaza & inválida \\
\texttt{0} & 1 & rechaza & rechaza & inválida \\
\texttt{001} & 3 & rechaza & rechaza & inválida \\
\end{longtable}

\evidenciaaqui[0.95]{../evidencias/jflap/fichas/lista3-239.png}{lista3-239 en JFLAP: Multiple Run con las tres cadenas válidas y las
tres inválidas de la tabla anterior.}

\hypertarget{lista-3-ejercicio-240}{%
\subsubsection{Lista 3 · Ejercicio 240}\label{lista-3-ejercicio-240}}

\textbf{Enunciado.} Expresión regular sobre \{0, 1\} para: Cadenas de longitud par que contienen 01.

\textbf{Diseño.} La expresión se convierte en AFN-λ con la construcción de Thompson y se verifica contra la definición del lenguaje. El enunciado se resuelve de forma natural con un AFN-λ de 78 estados; el AFD equivalente se obtuvo con la construcción de subconjuntos y la minimización y tiene 6 estados.

\textbf{Archivo.} \texttt{automatas/lista3/lista3-240.jff} · Σ = \{0, 1\} · \textbar Q\textbar{} = 6

\begin{longtable}[]{@{}
  >{\raggedright\arraybackslash}p{(\columnwidth - 6\tabcolsep) * \real{0.1579}}
  >{\raggedright\arraybackslash}p{(\columnwidth - 6\tabcolsep) * \real{0.1579}}
  >{\raggedright\arraybackslash}p{(\columnwidth - 6\tabcolsep) * \real{0.1579}}
  >{\raggedright\arraybackslash}p{(\columnwidth - 6\tabcolsep) * \real{0.5263}}@{}}
\caption{Tabla de transiciones (→ inicial, * aceptación, --- sin transición).}\tabularnewline
\toprule\noalign{}
\begin{minipage}[b]{\linewidth}\raggedright
δ
\end{minipage} & \begin{minipage}[b]{\linewidth}\raggedright
0
\end{minipage} & \begin{minipage}[b]{\linewidth}\raggedright
1
\end{minipage} & \begin{minipage}[b]{\linewidth}\raggedright
Qué recuerda del prefijo leído
\end{minipage} \\
\midrule\noalign{}
\endfirsthead
\toprule\noalign{}
\begin{minipage}[b]{\linewidth}\raggedright
δ
\end{minipage} & \begin{minipage}[b]{\linewidth}\raggedright
0
\end{minipage} & \begin{minipage}[b]{\linewidth}\raggedright
1
\end{minipage} & \begin{minipage}[b]{\linewidth}\raggedright
Qué recuerda del prefijo leído
\end{minipage} \\
\midrule\noalign{}
\endhead
\bottomrule\noalign{}
\endlastfoot
→ q0 & q1 & q2 & activos del AFN: \{q76\} o \{q4,q0,q2,q9,q11,q13,q14,q41,q43,q48,q44,q46\} \\
q1 & q3 & q4 & activos del AFN: \{q1,q5,q10,q6,q8,q15,q16,q33,q37,q50,q42,q38,q40,q51,q52\} o \{q1,q5,q10,q6,q8,q15,q16,q42,q38,q40,q45,q49,q51,q52\} \\
q2 & q3 & q0 & activos del AFN: \{q3,q5,q10,q6,q8,q35,q37,q50,q42,q38,q40,q51,q52\} o \{q3,q5,q10,q6,q8,q42,q38,q40,q47,q49,q51,q52\} \\
q3 & q1 & q5 & activos del AFN: \{q4,q0,q2,q7,q11,q13,q14,q39,q43,q48,q44,q46,q53,q54\} \\
* q4 & q5 & q5 & activos del AFN: \{q4,q0,q2,q9,q11,q13,q14,q17,q30,q22,q18,q20,q31,q41,q43,q48,q44,q46,q77\} o \{q4,q0,q2,q7,q11,q13,q14,q39,q43,q48,q44,q46,q53,q54,q57,q61,q66,q62,q64,q71,q75,q77\} o \{q4,q0,q2,q9,q11,q13,q14,q41,q43,q48,q44,q46,q59,q61,q66,q62,q64,q73,q75,q77\} o \{q4,q0,q2,q7,q11,q13,q14,q22,q18,q20,q25,q29,q31,q39,q43,q48,q44,q46,q53,q54,q77\} o \{q4,q0,q2,q9,q11,q13,q14,q17,q30,q22,q18,q20,q27,q29,q31,q41,q43,q48,q44,q46,q77\} o \{q4,q0,q2,q9,q11,q13,q14,q22,q18,q20,q27,q29,q31,q41,q43,q48,q44,q46,q77\} o \{q4,q0,q2,q9,q11,q13,q14,q17,q30,q22,q18,q20,q31,q41,q43,q48,q44,q46,q59,q61,q66,q62,q64,q73,q75,q77\} o \{q4,q0,q2,q7,q11,q13,q14,q22,q18,q20,q25,q29,q31,q39,q43,q48,q44,q46,q53,q54,q57,q61,q66,q62,q64,q71,q75,q77\} o \{q4,q0,q2,q9,q11,q13,q14,q22,q18,q20,q27,q29,q31,q41,q43,q48,q44,q46,q59,q61,q66,q62,q64,q73,q75,q77\} o \{q4,q0,q2,q9,q11,q13,q14,q17,q30,q22,q18,q20,q27,q29,q31,q41,q43,q48,q44,q46,q59,q61,q66,q62,q64,q73,q75,q77\} \\
q5 & q4 & q4 & activos del AFN: \{q3,q5,q10,q6,q8,q42,q38,q40,q47,q49,q51,q52,q55,q68,q60,q56,q58,q69,q74,q70,q72\} o \{q1,q5,q10,q6,q8,q15,q16,q19,q23,q28,q24,q26,q42,q38,q40,q45,q49,q51,q52\} o \{q3,q5,q10,q6,q8,q21,q23,q28,q24,q26,q42,q38,q40,q47,q49,q51,q52\} o \{q1,q5,q10,q6,q8,q15,q16,q42,q38,q40,q45,q49,q51,q52,q60,q56,q58,q63,q67,q69,q74,q70,q72\} o \{q3,q5,q10,q6,q8,q42,q38,q40,q47,q49,q51,q52,q55,q68,q60,q56,q58,q65,q67,q69,q74,q70,q72\} o \{q3,q5,q10,q6,q8,q42,q38,q40,q47,q49,q51,q52,q60,q56,q58,q65,q67,q69,q74,q70,q72\} o \{q3,q5,q10,q6,q8,q21,q23,q28,q24,q26,q42,q38,q40,q47,q49,q51,q52,q55,q68,q60,q56,q58,q69,q74,q70,q72\} o \{q1,q5,q10,q6,q8,q15,q16,q19,q23,q28,q24,q26,q42,q38,q40,q45,q49,q51,q52,q60,q56,q58,q63,q67,q69,q74,q70,q72\} o \{q3,q5,q10,q6,q8,q21,q23,q28,q24,q26,q42,q38,q40,q47,q49,q51,q52,q60,q56,q58,q65,q67,q69,q74,q70,q72\} o \{q3,q5,q10,q6,q8,q21,q23,q28,q24,q26,q42,q38,q40,q47,q49,q51,q52,q55,q68,q60,q56,q58,q65,q67,q69,q74,q70,q72\} \\
\end{longtable}

\begin{longtable}[]{@{}lclll@{}}
\caption{Cadenas de prueba.}\tabularnewline
\toprule\noalign{}
Cadena & Longitud & Esperado & Obtenido & Tipo de caso \\
\midrule\noalign{}
\endfirsthead
\toprule\noalign{}
Cadena & Longitud & Esperado & Obtenido & Tipo de caso \\
\midrule\noalign{}
\endhead
\bottomrule\noalign{}
\endlastfoot
\texttt{01} & 2 & acepta & acepta & válida \\
\texttt{100010} & 6 & acepta & acepta & válida \\
\texttt{00001001} & 8 & acepta & acepta & válida \\
λ & 0 & rechaza & rechaza & inválida \\
\texttt{1} & 1 & rechaza & rechaza & inválida \\
\texttt{001} & 3 & rechaza & rechaza & inválida \\
\end{longtable}

\evidenciaaqui[0.95]{../evidencias/jflap/fichas/lista3-240.png}{lista3-240 en JFLAP: Multiple Run con las tres cadenas válidas y las
tres inválidas de la tabla anterior.}

\hypertarget{lista-3-ejercicio-254}{%
\subsubsection{Lista 3 · Ejercicio 254}\label{lista-3-ejercicio-254}}

\textbf{Enunciado.} Expresión regular para el autómata de la tabla 254.

\textbf{Diseño.} Eliminación de estados sobre el autómata generalizado (inicial S y final F nuevos), en el orden de Q.

\textbf{Archivo.} \texttt{automatas/lista3/lista3-254.jff} · Σ = \{0, 1\} · \textbar Q\textbar{} = 2

\begin{longtable}[]{@{}lll@{}}
\caption{Tabla de transiciones (→ inicial, * aceptación, --- sin transición).}\tabularnewline
\toprule\noalign{}
δ & 0 & 1 \\
\midrule\noalign{}
\endfirsthead
\toprule\noalign{}
δ & 0 & 1 \\
\midrule\noalign{}
\endhead
\bottomrule\noalign{}
\endlastfoot
→ A & A & B \\
* B & A & B \\
\end{longtable}

\begin{longtable}[]{@{}lclll@{}}
\caption{Cadenas de prueba.}\tabularnewline
\toprule\noalign{}
Cadena & Longitud & Esperado & Obtenido & Tipo de caso \\
\midrule\noalign{}
\endfirsthead
\toprule\noalign{}
Cadena & Longitud & Esperado & Obtenido & Tipo de caso \\
\midrule\noalign{}
\endhead
\bottomrule\noalign{}
\endlastfoot
\texttt{1} & 1 & acepta & acepta & válida \\
\texttt{01} & 2 & acepta & acepta & válida \\
\texttt{0001} & 4 & acepta & acepta & válida \\
λ & 0 & rechaza & rechaza & inválida \\
\texttt{0} & 1 & rechaza & rechaza & inválida \\
\texttt{010} & 3 & rechaza & rechaza & inválida \\
\end{longtable}

\evidenciaaqui[0.95]{../evidencias/jflap/fichas/lista3-254.png}{lista3-254 en JFLAP: Multiple Run con las tres cadenas válidas y las
tres inválidas de la tabla anterior.}

\hypertarget{lista-3-ejercicio-255}{%
\subsubsection{Lista 3 · Ejercicio 255}\label{lista-3-ejercicio-255}}

\textbf{Enunciado.} Expresión regular para el autómata de la tabla 255.

\textbf{Diseño.} Eliminación de estados sobre el autómata generalizado (inicial S y final F nuevos), en el orden de Q.

\textbf{Archivo.} \texttt{automatas/lista3/lista3-255.jff} · Σ = \{0, 1\} · \textbar Q\textbar{} = 2

\begin{longtable}[]{@{}lll@{}}
\caption{Tabla de transiciones (→ inicial, * aceptación, --- sin transición).}\tabularnewline
\toprule\noalign{}
δ & 0 & 1 \\
\midrule\noalign{}
\endfirsthead
\toprule\noalign{}
δ & 0 & 1 \\
\midrule\noalign{}
\endhead
\bottomrule\noalign{}
\endlastfoot
→* A & B & B \\
B & A & B \\
\end{longtable}

\begin{longtable}[]{@{}lclll@{}}
\caption{Cadenas de prueba.}\tabularnewline
\toprule\noalign{}
Cadena & Longitud & Esperado & Obtenido & Tipo de caso \\
\midrule\noalign{}
\endfirsthead
\toprule\noalign{}
Cadena & Longitud & Esperado & Obtenido & Tipo de caso \\
\midrule\noalign{}
\endhead
\bottomrule\noalign{}
\endlastfoot
λ & 0 & acepta & acepta & válida \\
\texttt{00} & 2 & acepta & acepta & válida \\
\texttt{1111010} & 7 & acepta & acepta & válida \\
\texttt{1} & 1 & rechaza & rechaza & inválida \\
\texttt{01} & 2 & rechaza & rechaza & inválida \\
\texttt{1101} & 4 & rechaza & rechaza & inválida \\
\end{longtable}

\evidenciaaqui[0.95]{../evidencias/jflap/fichas/lista3-255.png}{lista3-255 en JFLAP: Multiple Run con las tres cadenas válidas y las
tres inválidas de la tabla anterior.}

\hypertarget{lista-3-ejercicio-256}{%
\subsubsection{Lista 3 · Ejercicio 256}\label{lista-3-ejercicio-256}}

\textbf{Enunciado.} Expresión regular para el autómata de la tabla 256.

\textbf{Diseño.} Eliminación de estados sobre el autómata generalizado (inicial S y final F nuevos), en el orden de Q.

\textbf{Archivo.} \texttt{automatas/lista3/lista3-256.jff} · Σ = \{0, 1, 2\} · \textbar Q\textbar{} = 2

\begin{longtable}[]{@{}llll@{}}
\caption{Tabla de transiciones (→ inicial, * aceptación, --- sin transición).}\tabularnewline
\toprule\noalign{}
δ & 0 & 1 & 2 \\
\midrule\noalign{}
\endfirsthead
\toprule\noalign{}
δ & 0 & 1 & 2 \\
\midrule\noalign{}
\endhead
\bottomrule\noalign{}
\endlastfoot
→ A & B & A & A \\
* B & B & B & A \\
\end{longtable}

\begin{longtable}[]{@{}lclll@{}}
\caption{Cadenas de prueba.}\tabularnewline
\toprule\noalign{}
Cadena & Longitud & Esperado & Obtenido & Tipo de caso \\
\midrule\noalign{}
\endfirsthead
\toprule\noalign{}
Cadena & Longitud & Esperado & Obtenido & Tipo de caso \\
\midrule\noalign{}
\endhead
\bottomrule\noalign{}
\endlastfoot
\texttt{0} & 1 & acepta & acepta & válida \\
\texttt{01} & 2 & acepta & acepta & válida \\
\texttt{0000} & 4 & acepta & acepta & válida \\
λ & 0 & rechaza & rechaza & inválida \\
\texttt{2} & 1 & rechaza & rechaza & inválida \\
\texttt{202} & 3 & rechaza & rechaza & inválida \\
\end{longtable}

\evidenciaaqui[0.95]{../evidencias/jflap/fichas/lista3-256.png}{lista3-256 en JFLAP: Multiple Run con las tres cadenas válidas y las
tres inválidas de la tabla anterior.}

\hypertarget{lista-3-ejercicio-257}{%
\subsubsection{Lista 3 · Ejercicio 257}\label{lista-3-ejercicio-257}}

\textbf{Enunciado.} Expresión regular para el autómata de la tabla 257.

\textbf{Diseño.} Eliminación de estados sobre el autómata generalizado (inicial S y final F nuevos), en el orden de Q.

\textbf{Archivo.} \texttt{automatas/lista3/lista3-257.jff} · Σ = \{0, 1\} · \textbar Q\textbar{} = 3

\begin{longtable}[]{@{}lll@{}}
\caption{Tabla de transiciones (→ inicial, * aceptación, --- sin transición).}\tabularnewline
\toprule\noalign{}
δ & 0 & 1 \\
\midrule\noalign{}
\endfirsthead
\toprule\noalign{}
δ & 0 & 1 \\
\midrule\noalign{}
\endhead
\bottomrule\noalign{}
\endlastfoot
→ A & B & C \\
* B & C & B \\
* C & A & A \\
\end{longtable}

\begin{longtable}[]{@{}lclll@{}}
\caption{Cadenas de prueba.}\tabularnewline
\toprule\noalign{}
Cadena & Longitud & Esperado & Obtenido & Tipo de caso \\
\midrule\noalign{}
\endfirsthead
\toprule\noalign{}
Cadena & Longitud & Esperado & Obtenido & Tipo de caso \\
\midrule\noalign{}
\endhead
\bottomrule\noalign{}
\endlastfoot
\texttt{0} & 1 & acepta & acepta & válida \\
\texttt{00} & 2 & acepta & acepta & válida \\
\texttt{1101} & 4 & acepta & acepta & válida \\
λ & 0 & rechaza & rechaza & inválida \\
\texttt{001} & 3 & rechaza & rechaza & inválida \\
\texttt{00010001} & 8 & rechaza & rechaza & inválida \\
\end{longtable}

\evidenciaaqui[0.95]{../evidencias/jflap/fichas/lista3-257.png}{lista3-257 en JFLAP: Multiple Run con las tres cadenas válidas y las
tres inválidas de la tabla anterior.}

\hypertarget{lista-3-ejercicio-258}{%
\subsubsection{Lista 3 · Ejercicio 258}\label{lista-3-ejercicio-258}}

\textbf{Enunciado.} Expresión regular para el autómata de la tabla 258.

\textbf{Diseño.} Eliminación de estados sobre el autómata generalizado (inicial S y final F nuevos), en el orden de Q.

\textbf{Archivo.} \texttt{automatas/lista3/lista3-258.jff} · Σ = \{0, 1\} · \textbar Q\textbar{} = 3

\begin{longtable}[]{@{}lll@{}}
\caption{Tabla de transiciones (→ inicial, * aceptación, --- sin transición).}\tabularnewline
\toprule\noalign{}
δ & 0 & 1 \\
\midrule\noalign{}
\endfirsthead
\toprule\noalign{}
δ & 0 & 1 \\
\midrule\noalign{}
\endhead
\bottomrule\noalign{}
\endlastfoot
→ A & C & C \\
B & A & C \\
* C & B & A \\
\end{longtable}

\begin{longtable}[]{@{}lclll@{}}
\caption{Cadenas de prueba.}\tabularnewline
\toprule\noalign{}
Cadena & Longitud & Esperado & Obtenido & Tipo de caso \\
\midrule\noalign{}
\endfirsthead
\toprule\noalign{}
Cadena & Longitud & Esperado & Obtenido & Tipo de caso \\
\midrule\noalign{}
\endhead
\bottomrule\noalign{}
\endlastfoot
\texttt{1} & 1 & acepta & acepta & válida \\
\texttt{01100011} & 8 & acepta & acepta & válida \\
\texttt{1110100010} & 10 & acepta & acepta & válida \\
λ & 0 & rechaza & rechaza & inválida \\
\texttt{0000101} & 7 & rechaza & rechaza & inválida \\
\texttt{00000100} & 8 & rechaza & rechaza & inválida \\
\end{longtable}

\evidenciaaqui[0.95]{../evidencias/jflap/fichas/lista3-258.png}{lista3-258 en JFLAP: Multiple Run con las tres cadenas válidas y las
tres inválidas de la tabla anterior.}

\hypertarget{lista-3-ejercicio-260}{%
\subsubsection{Lista 3 · Ejercicio 260}\label{lista-3-ejercicio-260}}

\textbf{Enunciado.} AFD equivalente al AFN de la tabla 260; describa cómo se obtuvo.

\textbf{Diseño.} Construcción de subconjuntos: se parte de \{q0\} y, para cada subconjunto nuevo y cada símbolo, se reúnen los destinos; sólo se crean subconjuntos alcanzables. El enunciado se resuelve de forma natural con un AFN de 5 estados; el AFD equivalente se obtuvo con la construcción de subconjuntos y la minimización y tiene 6 estados.

\textbf{Archivo.} \texttt{automatas/lista3/lista3-260.jff} · Σ = \{0, 1\} · \textbar Q\textbar{} = 6

\begin{longtable}[]{@{}llll@{}}
\caption{Tabla de transiciones (→ inicial, * aceptación, --- sin transición).}\tabularnewline
\toprule\noalign{}
δ & 0 & 1 & Qué recuerda del prefijo leído \\
\midrule\noalign{}
\endfirsthead
\toprule\noalign{}
δ & 0 & 1 & Qué recuerda del prefijo leído \\
\midrule\noalign{}
\endhead
\bottomrule\noalign{}
\endlastfoot
→ q0 & --- & q1 & activos del AFN: \{A\} \\
q1 & q2 & q3 & activos del AFN: \{B,C\} \\
* q2 & q2 & q2 & activos del AFN: \{A,B,C,D\} o \{A,B,C,D,E\} \\
* q3 & q4 & q2 & activos del AFN: \{A,C,D,E\} \\
q4 & q4 & q5 & activos del AFN: \{A,C,E\} \\
q5 & q2 & q2 & activos del AFN: \{A,B,C,E\} \\
\end{longtable}

\begin{longtable}[]{@{}lclll@{}}
\caption{Cadenas de prueba.}\tabularnewline
\toprule\noalign{}
Cadena & Longitud & Esperado & Obtenido & Tipo de caso \\
\midrule\noalign{}
\endfirsthead
\toprule\noalign{}
Cadena & Longitud & Esperado & Obtenido & Tipo de caso \\
\midrule\noalign{}
\endhead
\bottomrule\noalign{}
\endlastfoot
\texttt{10} & 2 & acepta & acepta & válida \\
\texttt{101} & 3 & acepta & acepta & válida \\
\texttt{10011} & 5 & acepta & acepta & válida \\
λ & 0 & rechaza & rechaza & inválida \\
\texttt{1} & 1 & rechaza & rechaza & inválida \\
\texttt{011} & 3 & rechaza & rechaza & inválida \\
\end{longtable}

\evidenciaaqui[0.95]{../evidencias/jflap/fichas/lista3-260.png}{lista3-260 en JFLAP: Multiple Run con las tres cadenas válidas y las
tres inválidas de la tabla anterior.}

\hypertarget{lista-3-ejercicio-261}{%
\subsubsection{Lista 3 · Ejercicio 261}\label{lista-3-ejercicio-261}}

\textbf{Enunciado.} AFD equivalente al AFN de la figura 261 (transcrito).

\textbf{Diseño.} Construcción de subconjuntos: se parte de \{q0\} y, para cada subconjunto nuevo y cada símbolo, se reúnen los destinos; sólo se crean subconjuntos alcanzables. El enunciado se resuelve de forma natural con un AFN de 5 estados; el AFD equivalente se obtuvo con la construcción de subconjuntos y la minimización y tiene 3 estados.

\textbf{Archivo.} \texttt{automatas/lista3/lista3-261.jff} · Σ = \{0, 1\} · \textbar Q\textbar{} = 3

\begin{longtable}[]{@{}llll@{}}
\caption{Tabla de transiciones (→ inicial, * aceptación, --- sin transición).}\tabularnewline
\toprule\noalign{}
δ & 0 & 1 & Qué recuerda del prefijo leído \\
\midrule\noalign{}
\endfirsthead
\toprule\noalign{}
δ & 0 & 1 & Qué recuerda del prefijo leído \\
\midrule\noalign{}
\endhead
\bottomrule\noalign{}
\endlastfoot
→ q0 & q1 & q0 & activos del AFN: \{q0\} o \{q0,q3\} o \{q0,q3,q4\} \\
q1 & q2 & q0 & activos del AFN: \{q0,q1\} \\
* q2 & q2 & q0 & activos del AFN: \{q0,q1,q2\} \\
\end{longtable}

\begin{longtable}[]{@{}lclll@{}}
\caption{Cadenas de prueba.}\tabularnewline
\toprule\noalign{}
Cadena & Longitud & Esperado & Obtenido & Tipo de caso \\
\midrule\noalign{}
\endfirsthead
\toprule\noalign{}
Cadena & Longitud & Esperado & Obtenido & Tipo de caso \\
\midrule\noalign{}
\endhead
\bottomrule\noalign{}
\endlastfoot
\texttt{00} & 2 & acepta & acepta & válida \\
\texttt{000} & 3 & acepta & acepta & válida \\
\texttt{01000} & 5 & acepta & acepta & válida \\
λ & 0 & rechaza & rechaza & inválida \\
\texttt{0} & 1 & rechaza & rechaza & inválida \\
\texttt{010} & 3 & rechaza & rechaza & inválida \\
\end{longtable}

\evidenciaaqui[0.95]{../evidencias/jflap/fichas/lista3-261.png}{lista3-261 en JFLAP: Multiple Run con las tres cadenas válidas y las
tres inválidas de la tabla anterior.}

\hypertarget{lista-3-ejercicio-262}{%
\subsubsection{Lista 3 · Ejercicio 262}\label{lista-3-ejercicio-262}}

\textbf{Enunciado.} Diseñe un AFN que acepte \{ab, ba\} y obtenga un AFD.

\textbf{Diseño.} Construcción de subconjuntos: se parte de \{q0\} y, para cada subconjunto nuevo y cada símbolo, se reúnen los destinos; sólo se crean subconjuntos alcanzables. El AFN tiene dos ramas desde q0, una por palabra.

\textbf{Archivo.} \texttt{automatas/lista3/lista3-262.jff} · Σ = \{a, b\} · \textbar Q\textbar{} = 5

\begin{longtable}[]{@{}lll@{}}
\caption{Tabla de transiciones (→ inicial, * aceptación, --- sin transición).}\tabularnewline
\toprule\noalign{}
δ & a & b \\
\midrule\noalign{}
\endfirsthead
\toprule\noalign{}
δ & a & b \\
\midrule\noalign{}
\endhead
\bottomrule\noalign{}
\endlastfoot
→ q0 & q1 & q3 \\
q1 & --- & q2 \\
* q2 & --- & --- \\
q3 & q4 & --- \\
* q4 & --- & --- \\
\end{longtable}

\begin{longtable}[]{@{}
  >{\raggedright\arraybackslash}p{(\columnwidth - 8\tabcolsep) * \real{0.2000}}
  >{\centering\arraybackslash}p{(\columnwidth - 8\tabcolsep) * \real{0.2000}}
  >{\raggedright\arraybackslash}p{(\columnwidth - 8\tabcolsep) * \real{0.2000}}
  >{\raggedright\arraybackslash}p{(\columnwidth - 8\tabcolsep) * \real{0.2000}}
  >{\raggedright\arraybackslash}p{(\columnwidth - 8\tabcolsep) * \real{0.2000}}@{}}
\caption{Cadenas de prueba.}\tabularnewline
\toprule\noalign{}
\begin{minipage}[b]{\linewidth}\raggedright
Cadena
\end{minipage} & \begin{minipage}[b]{\linewidth}\centering
Longitud
\end{minipage} & \begin{minipage}[b]{\linewidth}\raggedright
Esperado
\end{minipage} & \begin{minipage}[b]{\linewidth}\raggedright
Obtenido
\end{minipage} & \begin{minipage}[b]{\linewidth}\raggedright
Tipo de caso
\end{minipage} \\
\midrule\noalign{}
\endfirsthead
\toprule\noalign{}
\begin{minipage}[b]{\linewidth}\raggedright
Cadena
\end{minipage} & \begin{minipage}[b]{\linewidth}\centering
Longitud
\end{minipage} & \begin{minipage}[b]{\linewidth}\raggedright
Esperado
\end{minipage} & \begin{minipage}[b]{\linewidth}\raggedright
Obtenido
\end{minipage} & \begin{minipage}[b]{\linewidth}\raggedright
Tipo de caso
\end{minipage} \\
\midrule\noalign{}
\endhead
\bottomrule\noalign{}
\endlastfoot
\texttt{ab} & 2 & acepta & acepta & válida \\
\texttt{ba} & 2 & acepta & acepta & válida \\
λ & 0 & rechaza & rechaza & inválida \\
\texttt{a} & 1 & rechaza & rechaza & inválida \\
\texttt{bba} & 3 & rechaza & rechaza & inválida \\
\texttt{b} & 1 & rechaza & rechaza & frontera: justo por debajo de la longitud mínima \\
\end{longtable}

\evidenciaaqui[0.95]{../evidencias/jflap/fichas/lista3-262.png}{lista3-262 en JFLAP: Multiple Run con las tres cadenas válidas y las
tres inválidas de la tabla anterior.}

\hypertarget{lista-3-ejercicio-263}{%
\subsubsection{Lista 3 · Ejercicio 263}\label{lista-3-ejercicio-263}}

\textbf{Enunciado.} AFD equivalente al AFN de la figura 263 (transcrito).

\textbf{Diseño.} Construcción de subconjuntos: se parte de \{q0\} y, para cada subconjunto nuevo y cada símbolo, se reúnen los destinos; sólo se crean subconjuntos alcanzables. El enunciado se resuelve de forma natural con un AFN de 3 estados; el AFD equivalente se obtuvo con la construcción de subconjuntos y la minimización y tiene 3 estados.

\textbf{Archivo.} \texttt{automatas/lista3/lista3-263.jff} · Σ = \{0, 1\} · \textbar Q\textbar{} = 3

\begin{longtable}[]{@{}llll@{}}
\caption{Tabla de transiciones (→ inicial, * aceptación, --- sin transición).}\tabularnewline
\toprule\noalign{}
δ & 0 & 1 & Qué recuerda del prefijo leído \\
\midrule\noalign{}
\endfirsthead
\toprule\noalign{}
δ & 0 & 1 & Qué recuerda del prefijo leído \\
\midrule\noalign{}
\endhead
\bottomrule\noalign{}
\endlastfoot
→ q0 & q1 & q0 & activos del AFN: \{q0\} \\
q1 & q1 & q2 & activos del AFN: \{q0,q1\} \\
* q2 & q1 & q0 & activos del AFN: \{q0,q2\} \\
\end{longtable}

\begin{longtable}[]{@{}lclll@{}}
\caption{Cadenas de prueba.}\tabularnewline
\toprule\noalign{}
Cadena & Longitud & Esperado & Obtenido & Tipo de caso \\
\midrule\noalign{}
\endfirsthead
\toprule\noalign{}
Cadena & Longitud & Esperado & Obtenido & Tipo de caso \\
\midrule\noalign{}
\endhead
\bottomrule\noalign{}
\endlastfoot
\texttt{01} & 2 & acepta & acepta & válida \\
\texttt{001} & 3 & acepta & acepta & válida \\
\texttt{00001} & 5 & acepta & acepta & válida \\
λ & 0 & rechaza & rechaza & inválida \\
\texttt{0} & 1 & rechaza & rechaza & inválida \\
\texttt{100} & 3 & rechaza & rechaza & inválida \\
\end{longtable}

\evidenciaaqui[0.95]{../evidencias/jflap/fichas/lista3-263.png}{lista3-263 en JFLAP: Multiple Run con las tres cadenas válidas y las
tres inválidas de la tabla anterior.}

\endgroup

\end{document}
